% Produire des écrits utilitaires : le CV — Français Tle Tech — Cours signé OnBuch+
% Programme officiel MINESEC. Compiler avec : tectonic N23-produire-ecrits-utilitaires-cv.tex
\newcommand{\DOCMATIERE}{Français}
\newcommand{\DOCNIVEAU}{Tle Tech}
\newcommand{\DOCMODULE}{Français – classe de Terminale technique}
\newcommand{\DOCLECON}{N23}
\newcommand{\DOCTITRE}{Produire des écrits utilitaires : le CV}
% OnBuch+ — préambule « cours signés » ENRICHI (compilé avec tectonic / XeLaTeX)
% Système visuel « Pop » d'OnBuch+ : fond crème (jamais blanc), bordures encre
% épaisses, ombres « dures » décalées, titres Archivo Black en capitales.
% Version MATHÉMATIQUES : graphes (pgfplots/tikz), boîtes pédagogiques
% (définition, propriété, méthode, attention, le savais-tu, exercice résolu,
% exercice, TP, bilan), page de garde et signature OnBuch+ ancrée en bas.
%
% Macros attendues dans main.tex AVANT \input{preamble} :
%   \DOCMATIERE  \DOCNIVEAU  \DOCMODULE  \DOCLECON (numéro)  \DOCTITRE
\documentclass[11pt]{article}

\usepackage{fontspec}
\usepackage[french]{babel}
\usepackage[a4paper,margin=2cm,top=3.4cm,bottom=2.8cm,headheight=1.4cm,footskip=1.3cm]{geometry}
\usepackage{amsmath,amssymb,amsfonts}
\usepackage{mathtools} % \xmapsto, \coloneqq... souvent utilisés par les modèles
\usepackage{cancel} % simplifications barrées (bilans de forces)
% \abs{z} (module, valeur absolue) et \norm{u} (norme), utilisés par les modèles
\providecommand{\abs}{}\let\abs\relax\DeclarePairedDelimiter{\abs}{\lvert}{\rvert}
\providecommand{\norm}{}\let\norm\relax\DeclarePairedDelimiter{\norm}{\lVert}{\rVert}
\usepackage[table]{xcolor} % [table] : fournit aussi \rowcolors (alternance de lignes)
\usepackage{pagecolor}
\usepackage{tikz}
\usetikzlibrary{arrows.meta,positioning,calc,shapes.geometric,shapes.misc,shapes.arrows,decorations.pathmorphing,decorations.pathreplacing,decorations.markings,patterns,shadows,mindmap,trees,fit,backgrounds,matrix,intersections,angles,quotes}
\usepackage{pgfplots}
\usepackage[version=4]{mhchem} % équations nucléaires \ce{^{235}_{92}U}
\usepackage[european,siunitx]{circuitikz} % schémas électriques
\pgfplotsset{compat=1.18}
\usepgfplotslibrary{fillbetween,groupplots} % aires entre courbes (calcul intégral)
\usepackage{enumitem}
\usepackage{fancyhdr}
% [most] charge toutes les bibliothèques tcolorbox (listings, minted, raster,
% documentation...) et gonfle inutilement la mémoire TeX sur un document long
% (~300 boîtes) : on ne charge que ce qui est réellement utilisé.
\usepackage[skins,breakable]{tcolorbox}
\usepackage{graphicx}
\usepackage{colortbl}
\usepackage{booktabs}
\usepackage{tabularx}
\usepackage{multirow}
\usepackage{diagbox} % case d'en-tête diagonale des tableaux à double entrée
% Colonnes extensibles centrée / gauche / droite (C, L, R), souvent utilisées par les modèles
\newcolumntype{C}{>{\centering\arraybackslash}X}
\newcolumntype{L}{>{\raggedright\arraybackslash}X}
\newcolumntype{R}{>{\raggedleft\arraybackslash}X}
\usepackage{array}
\usepackage{multicol}
\usepackage{siunitx}
% parse-numbers=false : accepte aussi \SI{2,5 \times 10^{-3}}{...}
\sisetup{locale=FR,output-decimal-marker={,},group-minimum-digits=4,parse-numbers=false}
\DeclareSIUnit\ohm{\ensuremath{\symup{\Omega}}} % U+2126 absent de la police
\DeclareSIUnit\division{div} % réglages d'oscilloscope (ms/div, V/div)
\DeclareSIUnit\tr{tr} % tours (vitesse de rotation, ex. \SI{1500}{\tr\per\minute})
\DeclareSIUnit\molar{M} % concentration molaire (ex. \SI{3}{\molar}, \SI{1}{\micro\molar})
\DeclareSIUnit\an{an} % année (le modèle confond parfois avec \year, un compteur TeX) — ex. \SI{5}{\kilo\meter\per\an}
\DeclareSIUnit\FCFA{FCFA} % franc CFA (ex. \SI{500}{\FCFA\per\kilo\gram})
\DeclareSIUnit\degreeBrix{\textdegree Brix} % teneur en sucre d'un sirop/jus (réfractométrie)
\DeclareSIUnit\month{mois} % le modèle utilise parfois \month, un compteur TeX, comme unité de durée
\DeclareSIUnit\microliter{\micro\liter} % le modèle écrit parfois \microliter en un seul token
\DeclareSIUnit\million{millions} % dénombrement (ex. \SI{15}{\million\per\mL} spermatozoïdes)
\usepackage{qrcode}
% \degree : commande fréquemment utilisée par les modèles pour un angle ou une
% température en dehors de \SI{}{} (non fournie par les paquets chargés).
\providecommand{\degree}{^{\circ}}
% \qed (fin de démonstration), utilisé par les modèles sans amsthm
\providecommand{\qed}{\hfill$\square$}
% \barre : double barre des valeurs interdites dans un tableau de variations (tkz-tab)
\providecommand{\barre}{\ensuremath{\|}}
% environnement proof (amsthm non chargé) utilisé par les modèles
\ifdefined\proof\else\newenvironment{proof}[1][Démonstration]{\par\noindent\textit{#1.}\ }{\qed\par}\fi
\usepackage{lastpage}
\usepackage{hyperref}
\hypersetup{hidelinks}

% ============================================================
% Palette « Pop » OnBuch+ — mêmes hex que la classe P (plus_ui.dart)
% ============================================================
\definecolor{poporange}{HTML}{F59321}
\definecolor{popdark}{HTML}{E07A0C}
\definecolor{popcream}{HTML}{FFF6E8}
\definecolor{popink}{HTML}{1C1714}
\definecolor{poppurple}{HTML}{7A5AE0}
\definecolor{popgreen}{HTML}{2E9E6A}
\definecolor{poppink}{HTML}{E0457B}
\definecolor{popblue}{HTML}{2F7DE1}
\definecolor{popgold}{HTML}{C98A12}
\definecolor{popmuted}{HTML}{8A7F76}
\definecolor{popline}{HTML}{EADFD2}
\definecolor{popgray}{HTML}{8A7F76}
\colorlet{popgrey}{popgray}
% Alias tolérants (le modèle nomme parfois les couleurs différemment)
\colorlet{popwhite}{white}
\colorlet{popblack}{popink}
\colorlet{popred}{poppink}
\colorlet{popyellow}{popgold}
% Teintes claires (fonds de tableaux/encadrés) — le modèle nomme parfois
% les couleurs avec un suffixe L (ex. popblueL) sans les avoir définies.
\colorlet{poporangeL}{poporange!20}
\colorlet{popdarkL}{popdark!20}
\colorlet{poppurpleL}{poppurple!20}
\colorlet{popgreenL}{popgreen!20}
\colorlet{poppinkL}{poppink!20}
\colorlet{popblueL}{popblue!20}
\colorlet{popgoldL}{popgold!20}
\colorlet{popredL}{popred!20}
\colorlet{popgrayL}{popgray!20}
% Teintes claires pour fonds de tableaux / zones
\colorlet{popblueL}{popblue!10!white}
\colorlet{poporangeL}{poporange!14!white}
\colorlet{popgreenL}{popgreen!12!white}
\colorlet{poppurpleL}{poppurple!12!white}
\colorlet{poppinkL}{poppink!12!white}
\colorlet{popgoldL}{popgold!14!white}

\pagecolor{popcream}

% ============================================================
% Typographie
% ============================================================
\defaultfontfeatures{Ligatures=TeX}
\setmainfont{PlusJakartaSans-Medium.ttf}[Path=../../../fonts/,BoldFont=PlusJakartaSans-Bold.ttf]
% Maths en sans-serif (Fira Math) avec chiffres et lettres droites de Plus
% Jakarta, pour que les formules se fondent dans le texte.
\usepackage[math-style=ISO,bold-style=ISO]{unicode-math}
\setmathfont{FiraMath-Regular.otf}[Path=../../../fonts/]
\setmathfont{PlusJakartaSans-Medium.ttf}[Path=../../../fonts/,range={up/num,up/{Latin,latin},\mathup}]
\usepackage{icomma} % virgule décimale sans espace en mode math
% Caractères mathématiques Unicode absents des polices de texte (Plus Jakarta,
% Archivo Black) : rendus par leur équivalent mathématique, en texte comme en
% formule — sinon ils s'affichent en carré vide (ex. « Arithmétique dans ℤ »).
\usepackage{newunicodechar}
\newunicodechar{ℝ}{\ensuremath{\mathbb{R}}}
\newunicodechar{ℤ}{\ensuremath{\mathbb{Z}}}
\newunicodechar{ℂ}{\ensuremath{\mathbb{C}}}
\newunicodechar{ℕ}{\ensuremath{\mathbb{N}}}
\newunicodechar{ℚ}{\ensuremath{\mathbb{Q}}}
\newunicodechar{𝔻}{\ensuremath{\mathbb{D}}}
\newunicodechar{α}{\ensuremath{\alpha}}
\newunicodechar{β}{\ensuremath{\beta}}
\newunicodechar{γ}{\ensuremath{\gamma}}
\newunicodechar{δ}{\ensuremath{\delta}}
\newunicodechar{ε}{\ensuremath{\varepsilon}}
\newunicodechar{θ}{\ensuremath{\theta}}
\newunicodechar{λ}{\ensuremath{\lambda}}
\newunicodechar{μ}{\ensuremath{\mu}}
\newunicodechar{σ}{\ensuremath{\sigma}}
\newunicodechar{τ}{\ensuremath{\tau}}
\newunicodechar{φ}{\ensuremath{\varphi}}
\newunicodechar{ψ}{\ensuremath{\psi}}
\newunicodechar{ω}{\ensuremath{\omega}}
\newunicodechar{Σ}{\ensuremath{\Sigma}}
% majuscules produites par \MakeUppercase dans les titres (α-aminés -> Α)
\newunicodechar{Α}{\ensuremath{\alpha}}
\newunicodechar{Β}{\ensuremath{\beta}}
\newunicodechar{Γ}{\ensuremath{\Gamma}}
\newunicodechar{Δ}{\ensuremath{\Delta}}
\newunicodechar{Ε}{\ensuremath{\varepsilon}}
\newunicodechar{Θ}{\ensuremath{\Theta}}
\newunicodechar{Λ}{\ensuremath{\Lambda}}
\newunicodechar{Μ}{\ensuremath{\mu}}
\newunicodechar{Τ}{\ensuremath{\tau}}
\newunicodechar{Φ}{\ensuremath{\Phi}}
\newunicodechar{Ψ}{\ensuremath{\Psi}}
\newunicodechar{Ω}{\ensuremath{\Omega}}
\newunicodechar{↦}{\ensuremath{\mapsto}}
\newunicodechar{∈}{\ensuremath{\in}}
\newunicodechar{∉}{\ensuremath{\notin}}
\newunicodechar{∧}{\ensuremath{\wedge}}
\newunicodechar{⇔}{\ensuremath{\Leftrightarrow}}
\newunicodechar{⟺}{\ensuremath{\Longleftrightarrow}}
\newunicodechar{⇒}{\ensuremath{\Rightarrow}}
\newunicodechar{⟹}{\ensuremath{\Longrightarrow}}
\newunicodechar{∘}{\ensuremath{\circ}}
\newunicodechar{∪}{\ensuremath{\cup}}
\newunicodechar{∩}{\ensuremath{\cap}}
\newunicodechar{≡}{\ensuremath{\equiv}}
\newunicodechar{⊂}{\ensuremath{\subset}}
\newunicodechar{⊥}{\ensuremath{\perp}}
\newunicodechar{∃}{\ensuremath{\exists}}
\newunicodechar{∀}{\ensuremath{\forall}}
\newunicodechar{∛}{\ensuremath{\sqrt[3]{\,}}}
\newunicodechar{∜}{\ensuremath{\sqrt[4]{\,}}}
\newunicodechar{⃗}{\ensuremath{{}^{\to}}}
\newunicodechar{ⁿ}{\ifmmode^{n}\else\textsuperscript{\ensuremath{n}}\fi}
\newunicodechar{ˣ}{\ifmmode^{x}\else\textsuperscript{\ensuremath{x}}\fi}
\newunicodechar{ᵇ}{\ifmmode^{b}\else\textsuperscript{\ensuremath{b}}\fi}
\newunicodechar{ʳ}{\ifmmode^{r}\else\textsuperscript{\ensuremath{r}}\fi}
\newunicodechar{ᵘ}{\ifmmode^{u}\else\textsuperscript{\ensuremath{u}}\fi}
\newunicodechar{ʸ}{\ifmmode^{y}\else\textsuperscript{\ensuremath{y}}\fi}
\newunicodechar{ᵃ}{\ifmmode^{a}\else\textsuperscript{\ensuremath{a}}\fi}
\newunicodechar{ᵉ}{\ifmmode^{e}\else\textsuperscript{\ensuremath{e}}\fi}
\newunicodechar{ᵏ}{\ifmmode^{k}\else\textsuperscript{\ensuremath{k}}\fi}
\newunicodechar{ᵖ}{\ifmmode^{p}\else\textsuperscript{\ensuremath{p}}\fi}
\newunicodechar{ᴺ}{\ifmmode^{N}\else\textsuperscript{\ensuremath{N}}\fi}
\newunicodechar{ᶜ}{\ifmmode^{c}\else\textsuperscript{\ensuremath{c}}\fi}
\newunicodechar{ʰ}{\ifmmode^{h}\else\textsuperscript{\ensuremath{h}}\fi}
\newunicodechar{ᵠ}{\ifmmode^{q}\else\textsuperscript{\ensuremath{q}}\fi}
\newunicodechar{ₙ}{\ifmmode_{n}\else\textsubscript{\ensuremath{n}}\fi}
\newunicodechar{ₐ}{\ifmmode_{a}\else\textsubscript{\ensuremath{a}}\fi}
\newunicodechar{ᵢ}{\ifmmode_{i}\else\textsubscript{\ensuremath{i}}\fi}
\newunicodechar{ᵣ}{\ifmmode_{r}\else\textsubscript{\ensuremath{r}}\fi}
\newunicodechar{ₖ}{\ifmmode_{k}\else\textsubscript{\ensuremath{k}}\fi}
\newunicodechar{ᵦ}{\ifmmode_{\beta}\else\textsubscript{\ensuremath{\beta}}\fi}
\newunicodechar{ₓ}{\ifmmode_{x}\else\textsubscript{\ensuremath{x}}\fi}
\newfontfamily\popdisplay{ArchivoBlack-Regular.ttf}[Path=../../../fonts/]
\newfontfamily\poplogo{SpaceGrotesk-Bold.ttf}[Path=../../../fonts/]
\newfontfamily\popsemi{PlusJakartaSans-SemiBold.ttf}[Path=../../../fonts/]
\newfontfamily\popxbold{PlusJakartaSans-ExtraBold.ttf}[Path=../../../fonts/]
\newfontfamily\popxboldls{PlusJakartaSans-ExtraBold.ttf}[Path=../../../fonts/,LetterSpace=3.0]

\setlist[enumerate]{leftmargin=1.7em,itemsep=5pt,topsep=4pt}
\setlist[itemize]{leftmargin=1.7em,itemsep=4pt,topsep=4pt,label={\color{poporange}\textbullet}}
\setlength{\parindent}{0pt}
\setlength{\parskip}{5pt}
\renewcommand{\arraystretch}{1.25}

% pgfplots : style Pop par défaut
\pgfplotsset{
  every axis/.append style={axis line style={popink,line width=1pt},tick style={popink},
    grid style={popline,line width=0.5pt},label style={font=\small},tick label style={font=\footnotesize}},
  popcurve/.style={poporange,line width=1.8pt,no markers,smooth},
}
% Même style disponible dans un \draw TikZ simple (hors pgfplots)
\tikzset{popcurve/.style={poporange,line width=1.8pt}}
\tikzset{popgrid/.style={popline,very thin}}
\colorlet{popcurve}{poporange} % le nom du style est parfois utilisé comme couleur

% ============================================================
% Pill / squiggle
% ============================================================
\newcommand{\poppill}[3]{%
  \tikz[baseline=-0.55ex]\node[draw=popink,line width=1.2pt,rounded corners=2.6pt,%
    fill=#2,inner xsep=6.5pt,inner ysep=3pt]{%
    \popxboldls\fontsize{7}{7}\selectfont\color{#3}\MakeUppercase{#1}};%
}
\newcommand{\popsquiggle}[1]{%
  \tikz[baseline=-0.4ex]{
    \draw[#1,line width=2.6pt,line cap=round]
      plot [smooth] coordinates {(0,0.09) (0.34,0.01) (0.68,0.21) (1.02,0.01) (1.36,0.10)};
  }%
}
% Mot-clé surligné (marqueur orange sous le texte)
\newcommand{\cle}[1]{{\popxbold\color{popdark}#1}}

% ============================================================
% En-tête / pied
% ============================================================
\pagestyle{fancy}
\fancyhf{}
\renewcommand{\headrulewidth}{0pt}
\renewcommand{\footrulewidth}{0pt}
\lhead{\raisebox{-0.15ex}{\poplogo\fontsize{13}{13}\selectfont\color{popink}OnBuch\color{poporange}+}}
\rhead{\poppill{\DOCMATIERE\ \textbullet\ \DOCNIVEAU\ \textbullet\ Leçon \DOCLECON}{white}{popink}}
\lfoot{\popxbold\fontsize{7.6}{7.6}\selectfont\color{popmuted}\MakeUppercase{Cours signé OnBuch+}\ \textbullet\ {\popsemi\DOCTITRE}}
\rfoot{\tikz[baseline=-0.6ex]\node[rounded corners=8.5pt,fill=popink,minimum height=17pt,minimum width=17pt,inner xsep=5pt,inner sep=0]{\popxbold\fontsize{8}{8}\selectfont\color{popcream}\thepage\,/\,\pageref*{LastPage}};}

% ============================================================
% Titres
% ============================================================
\newcommand{\chaptitre}[2]{%
  \begin{tcolorbox}[enhanced,colback=poporange,colframe=popink,boxrule=2.6pt,arc=11pt,
      drop shadow=popink,left=15pt,right=15pt,top=12pt,bottom=11pt,before skip=3pt,after skip=18pt]
    \poppill{#2}{popink}{popcream}\par\vspace{9pt}
    {\raggedright\hyphenpenalty=10000\exhyphenpenalty=10000\popdisplay\fontsize{22}{24}\selectfont\color{popcream}\MakeUppercase{#1}\par}
  \end{tcolorbox}
}
% Section numérotée : pastille orange avec le numéro + titre Archivo
\newcounter{popsec}
\newcommand{\coursec}[1]{%
  \stepcounter{popsec}\setcounter{popsub}{0}%
  \par\vspace{16pt}\noindent
  \begin{minipage}{\linewidth}
  \tikz[baseline=-0.7ex]\node[draw=popink,line width=1.6pt,fill=poporange,rounded corners=4pt,minimum size=22pt,inner sep=0]{\popdisplay\fontsize{12}{12}\selectfont\color{popcream}\thepopsec};%
  \hspace{8pt}{\popdisplay\fontsize{15}{17}\selectfont\color{popink}\MakeUppercase{#1}}\par
  \vspace{4pt}\noindent{\color{poporange}\rule{36pt}{3.6pt}}
  \end{minipage}\par\nopagebreak\vspace{8pt}
}
\newcounter{popsub}
\newcommand{\courssub}[1]{%
  \stepcounter{popsub}%
  \par\vspace{9pt}\noindent{\popxbold\fontsize{11.5}{13}\selectfont\color{popink}{\color{poporange}\thepopsec.\thepopsub}\hspace{6pt}#1}\par\nopagebreak\vspace{4pt}
}

% ============================================================
% Boîtes pédagogiques — même recette : fond blanc, bordure épaisse colorée,
% ombre dure encre, badge plein. Titre optionnel [#1].
% ============================================================
\tcbset{popbase/.style={breakable,enhanced,colback=white,boxrule=2.3pt,arc=9pt,
  shadow={3pt}{-3pt}{0pt}{popink},left=14pt,right=14pt,top=11pt,bottom=11pt,before skip=11pt,after skip=15pt,
  fontupper=\popsemi\fontsize{10.6}{14.5}\selectfont\color{popink},
  fontlower=\popsemi\fontsize{10.6}{14.5}\selectfont\color{popink}}}
% Coche dessinée (le glyphe ✓ n'existe pas dans les polices du document)
\newcommand{\popcheck}[1]{\tikz[baseline=-0.5ex]\draw[#1,line width=1.6pt,line cap=round,line join=round](-0.13,0.0)--(-0.04,-0.09)--(0.13,0.1);}
\providecommand{\coche}{\popcheck{popgreen}}
\providecommand{\resultat}[1]{\par\smallskip\noindent\fcolorbox{popgreen}{popgreenL}{\parbox{\dimexpr\linewidth-2\fboxsep-2\fboxrule\relax}{\textbf{Résultat.} #1}}\par\smallskip}
\newcommand{\popboxhead}[3]{\poppill{#1}{#2}{white}\ifx\relax#3\relax\else\hspace{6pt}{\popxbold\fontsize{10.6}{12}\selectfont\color{popink}#3}\fi\par\vspace{7pt}}

\newtcolorbox{aretenir}[1][]{popbase,colframe=popgreen,before upper={\popboxhead{À retenir}{popgreen}{#1}}}
\newtcolorbox{exemplebox}[1][]{popbase,colframe=poporange,before upper={\popboxhead{Exemple}{poporange}{#1}}}
\newtcolorbox{definition}[1][]{popbase,colframe=popblue,before upper={\popboxhead{Définition}{popblue}{#1}}}
\newtcolorbox{propriete}[1][]{popbase,colframe=poppurple,before upper={\popboxhead{Propriété}{poppurple}{#1}}}
\newtcolorbox{methode}[1][]{popbase,colframe=popgold,before upper={\popboxhead{Méthode}{popgold}{#1}}}
\newtcolorbox{attention}[1][]{popbase,colframe=poppink,before upper={\popboxhead{Attention}{poppink}{#1}}}
\newtcolorbox{experience}[1][]{popbase,colframe=popblue,colback=popblueL,before upper={\popboxhead{Expérience}{popblue}{#1}}}
% « Le savais-tu ? » : carte violette pleine (contexte, culture, Cameroun)
\newtcolorbox{savaistu}[1][]{popbase,colback=poppurple,colframe=popink,
  fontupper=\popsemi\fontsize{10.4}{14.2}\selectfont\color{white},
  before upper={\poppill{Le savais-tu\,?}{white}{poppurple}\ifx\relax#1\relax\else\hspace{6pt}{\popxbold\color{white}#1}\fi\par\vspace{7pt}}}
% Exercice résolu : énoncé en haut, \tcblower puis la solution sur fond crème
\newcounter{popexo}\newcounter{popexr}
\newtcolorbox{exoresolu}[1][]{popbase,colframe=poporange,colbacklower=popcream,
  segmentation style={popink,line width=1.2pt,dashed},
  before upper={\stepcounter{popexr}\popboxhead{Exercice résolu \thepopexr}{poporange}{#1}},
  before lower={\poppill{Solution}{popink}{popcream}\par\vspace{6pt}}}
% Exercice (non résolu en place) — la correction est dans la partie Corrigés
\newtcolorbox{exercice}[1][]{popbase,colframe=popink,
  before upper={\stepcounter{popexo}\popboxhead{Exercice \thepopexo}{popink}{#1}}}
% Corrigé
\newtcolorbox{corrige}[1][]{popbase,colframe=popgreen,colback=popgreenL,
  before upper={\popboxhead{Corrigé}{popgreen}{#1}}}
% Objectifs / prérequis / bilan
\newtcolorbox{objectifs}{popbase,colframe=popink,colback=poporangeL,
  before upper={\popboxhead{Objectifs de la leçon}{popink}{}}}
\newtcolorbox{prerequis}{popbase,colframe=popink,colback=popblueL,
  before upper={\popboxhead{Prérequis}{popblue}{}}}
\newtcolorbox{bilan}{popbase,colframe=popink,colback=white,boxrule=2.8pt,
  before upper={\popboxhead{Fiche bilan}{poporange}{L'essentiel à connaître pour l'examen}}}
% Figure encadrée avec légende
\newtcolorbox{popfigure}[1][]{enhanced,breakable=false,colback=white,colframe=popline,boxrule=1.4pt,arc=8pt,
  left=10pt,right=10pt,top=10pt,bottom=8pt,before skip=10pt,after skip=12pt,halign=center,
  lower separated=false,halign lower=center,
  fontlower=\popsemi\fontsize{9}{11.5}\selectfont\color{popmuted},
  #1}
\newcommand{\legende}[1]{\par\vspace{6pt}{\popsemi\fontsize{9}{11.5}\selectfont\color{popmuted}#1}}

% ============================================================
% Signature OnBuch+
% ============================================================
\newcommand{\poplogoinline}[1]{{\poplogo\fontsize{#1}{#1}\selectfont\color{popink}OnBuch\color{poporange}+}}
\newcommand{\signatureonbuch}{%
  \begin{tcolorbox}[enhanced,colback=popink,colframe=popink,boxrule=2.2pt,arc=10pt,
      drop shadow=poporange,left=14pt,right=14pt,top=10pt,bottom=10pt,before skip=0pt,after skip=0pt]
    \tikz[baseline=-0.7ex]\node[circle,fill=popgreen,draw=popcream,line width=1.3pt,minimum size=22pt,inner sep=0]{\popcheck{white}};%
    \hspace{9pt}\begin{minipage}[c]{0.62\linewidth}
      {\popxbold\fontsize{12}{13}\selectfont\color{popcream}Signé\ \textbullet\ {\poplogo OnBuch\color{poporange}+}}\par\vspace{2pt}
      {\raggedright\popsemi\fontsize{8}{10.5}\selectfont\color{popline}\MakeUppercase{Équipe pédagogique OnBuch+}\\\MakeUppercase{Conforme au programme officiel MINESEC}\par}
    \end{minipage}\hfill
    {\poplogo\fontsize{22}{22}\selectfont\color{popcream}OnBuch\color{poporange}+}
  \end{tcolorbox}%
}

% ============================================================
% Page de garde
% #1 titre · #2 matière · #3 niveau · #4 module · #5 n° leçon · #6 durée ·
% #7 description / objectifs (texte ou itemize)
% ============================================================
\newcommand{\pagegarde}[7]{%
  \thispagestyle{empty}
  \newgeometry{a4paper,margin=1.7cm,top=1.6cm,bottom=1.4cm}%
  \begin{tikzpicture}[remember picture,overlay]
    \fill[poporange!18] ([xshift=-3.2cm,yshift=-3.6cm]current page.north east) circle (4.6cm);
    \fill[poppurple!14] ([xshift=2.4cm,yshift=7.6cm]current page.south west) circle (3.8cm);
  \end{tikzpicture}%
  {\poplogo\fontsize{30}{30}\selectfont\color{popink}OnBuch\color{poporange}+}\hfill
  \raisebox{0.6ex}{\poppill{Cours signé \textbullet\ Premium}{popink}{popcream}}\par
  \vspace{1pt}\popsquiggle{poppurple}\par
  \vspace{8pt}
  \begin{tcolorbox}[enhanced,colback=poporange,colframe=popink,boxrule=2.8pt,arc=13pt,
      drop shadow=popink,left=18pt,right=18pt,top=12pt,bottom=13pt,after skip=10pt]
    \poppill{#2 \textbullet\ #3}{popink}{popcream}\hspace{5pt}\poppill{#4}{white}{popink}\par\vspace{8pt}
    {\popdisplay\fontsize{14}{14}\selectfont\color{popink}LEÇON #5}\par\vspace{4pt}
    {\raggedright\hyphenpenalty=10000\exhyphenpenalty=10000\popdisplay\fontsize{25}{27}\selectfont\color{popcream}\MakeUppercase{#1}\par}
    \vspace{8pt}
    {\popxbold\fontsize{9.5}{9.5}\selectfont\color{popink}\MakeUppercase{Durée conseillée : #6}}
  \end{tcolorbox}
  \begin{tcolorbox}[enhanced,colback=white,colframe=popink,boxrule=2.2pt,arc=10pt,
      drop shadow=popink,left=16pt,right=16pt,top=10pt,bottom=10pt,after skip=8pt]
    \poppill{Dans cette leçon}{poppurple}{white}\par\vspace{5pt}
    \popsemi\fontsize{9.3}{12.6}\selectfont\color{popink}#7
  \end{tcolorbox}
  \noindent\begin{minipage}[t]{0.487\linewidth}
    \begin{tcolorbox}[enhanced,colback=popgreen,colframe=popink,boxrule=2.2pt,arc=10pt,
        drop shadow=popink,left=11pt,right=11pt,top=10pt,bottom=10pt,height=3.1cm,valign=center]
      \tcbox[colback=white,colframe=popink,boxrule=1.3pt,arc=5pt,boxsep=0pt,nobeforeafter,
        left=3pt,right=3pt,top=3pt,bottom=3pt,valign=center]{\qrcode[height=1.9cm]{https://play.google.com/store/apps/details?id=com.luvvix.onbuch&hl=fr}}%
      \hspace{8pt}\begin{minipage}[c]{0.5\linewidth}
        {\popdisplay\fontsize{11}{12.5}\selectfont\color{white}\MakeUppercase{Télécharge OnBuch}}\par\vspace{4pt}
        {\raggedright\popsemi\fontsize{8.4}{11}\selectfont\color{white}Gratuit \textbullet\ Google Play\\Tuteur IA, annales, concours, résultats\par}
      \end{minipage}
    \end{tcolorbox}
  \end{minipage}\hfill
  \begin{minipage}[t]{0.487\linewidth}
    \begin{tcolorbox}[enhanced,colback=poppurple,colframe=popink,boxrule=2.2pt,arc=10pt,
        drop shadow=popink,left=11pt,right=11pt,top=10pt,bottom=10pt,height=3.1cm,valign=center]
      \tikz[baseline=-0.7ex]\node[circle,fill=white,draw=popink,line width=1.3pt,minimum size=1.6cm,inner sep=0]{\popdisplay\fontsize{24}{24}\selectfont\color{poppurple}+};%
      \hspace{8pt}\begin{minipage}[c]{0.6\linewidth}
        {\popdisplay\fontsize{11}{12.5}\selectfont\color{white}\MakeUppercase{Passe à OnBuch+}}\par\vspace{4pt}
        {\raggedright\popsemi\fontsize{8.4}{11}\selectfont\color{white}Tous les cours signés, Tuteur IA illimité, fiches PDF — onglet OnBuch+ de l'app\par}
      \end{minipage}
    \end{tcolorbox}
  \end{minipage}
  \vspace{8pt}
  \popcontenu
  \vfill
  \signatureonbuch
  \restoregeometry
  \newpage
}

% Rangée « Ce cours contient » (page de garde)
\newcommand{\poptile}[3]{% #1 couleur #2 gros texte #3 légende
  \begin{tcolorbox}[enhanced,colback=#1,colframe=popink,boxrule=2pt,arc=9pt,shadow={2.5pt}{-2.5pt}{0pt}{popink},
      left=6pt,right=6pt,top=9pt,bottom=9pt,halign=center,valign=center,height=1.75cm,width=\linewidth,nobeforeafter]
    {\popdisplay\fontsize{12.5}{13}\selectfont\color{white}\MakeUppercase{#2}}\par\vspace{3pt}
    {\centering\popsemi\fontsize{7.8}{9.5}\selectfont\color{white}#3\par}
  \end{tcolorbox}}
\newcommand{\popcontenu}{%
  \par\noindent{\popxboldls\fontsize{7.5}{7.5}\selectfont\color{popmuted}CE COURS SIGNÉ CONTIENT}\par\vspace{7pt}
  \noindent\begin{minipage}[t]{0.235\linewidth}\poptile{popblue}{Cours}{complet, illustré, pas à pas}\end{minipage}\hfill
  \begin{minipage}[t]{0.235\linewidth}\poptile{poporange}{Méthodes}{et exercices résolus}\end{minipage}\hfill
  \begin{minipage}[t]{0.235\linewidth}\poptile{poppink}{Exercices}{type examen + corrigés}\end{minipage}\hfill
  \begin{minipage}[t]{0.235\linewidth}\poptile{popgreen}{Fiche bilan}{l'essentiel à retenir}\end{minipage}\par}

% Sommaire
\newtcolorbox{sommaire}{enhanced,colback=white,colframe=popink,boxrule=2.2pt,arc=10pt,
  drop shadow=popink,left=18pt,right=18pt,top=13pt,bottom=13pt,before skip=6pt,after skip=20pt,
  before upper={\poppill{Sommaire}{poppurple}{white}\par\vspace{9pt}\popsemi\fontsize{10.6}{14.5}\selectfont\color{popink}}}

% Signature de fin de cours — poussée tout en bas de la dernière page
\newcommand{\findecours}{%
  \par\vfill\nopagebreak
  \begin{center}\popsquiggle{poporange}\end{center}\vspace{4pt}
  \signatureonbuch
}
\AtBeginDocument{\def\llbracket{\mathopen{[\mkern-3.5mu[}}\def\rrbracket{\mathclose{]\mkern-3.5mu]}}} % intervalles d'entiers (glyphe ⟦ absent de la police)
% Glyphes absents de Fira Math : redéfinis à partir de glyphes disponibles
\AtBeginDocument{%
  \def\setminus{\mathbin{\backslash}}\let\smallsetminus\setminus
  \def\vdots{\mathord{\vbox{\baselineskip3.2pt\lineskiplimit0pt\kern4pt\hbox{.}\hbox{.}\hbox{.}}}}%
  \def\ddots{\mathinner{\mkern1mu\raise6.5pt\hbox{.}\mkern2mu\raise3.6pt\hbox{.}\mkern2mu\raise0.7pt\hbox{.}\mkern1mu}}%
  \def\triangle{\mathord{\scalebox{0.9}{$\symup{\Delta}$}}}%
  \def\nabla{\mathord{\raisebox{\depth}{\scalebox{1}[-1]{$\symup{\Delta}$}}}}%
  \def\checkmark{\popcheck{popdark}}%
  \def\star{\mathbin{\ast}}\def\bigstar{\mathord{\ast}}%
  \def\diamond{\mathbin{\scalebox{0.6}{$\Diamond$}}}%
  \def\bigcup{\mathop{\mathchoice{\raisebox{-0.3em}{\scalebox{1.7}{$\cup$}}}{\scalebox{1.25}{$\cup$}}{\cup}{\cup}}\displaylimits}%
  \def\bigcap{\mathop{\mathchoice{\raisebox{-0.3em}{\scalebox{1.7}{$\cap$}}}{\scalebox{1.25}{$\cap$}}{\cap}{\cap}}\displaylimits}%
  \def\lesseqqgtr{\mathrel{\lessgtr}}\def\bigoplus{\mathop{\oplus}\displaylimits}%
}
\providecommand{\ssi}{si et seulement si} % abréviation fréquente
\AtBeginDocument{\providecommand{\wideparen}{\overparen}} % arc de cercle

%% FIGFIX-BEGIN v4 — correctifs globaux des figures (docs/onbuch-generation/tools/figfix)
\usepackage{adjustbox}\usepackage{etoolbox}
% toute figure TikZ/pgfplots plus large que la ligne est réduite à la largeur disponible
\BeforeBeginEnvironment{tikzpicture}{\begin{adjustbox}{max width=\linewidth}}
\AfterEndEnvironment{tikzpicture}{\end{adjustbox}}
% légendes pgfplots lisibles par-dessus les courbes
\pgfplotsset{every axis/.append style={legend style={fill=white,fill opacity=0.92,text opacity=1,draw=black!15,inner sep=3pt}}}
% étiquettes d'arêtes (syntaxe edge["..."]) avec fond blanc
\tikzset{every edge quotes/.append style={fill=white,inner sep=1.2pt}}
%% FIGFIX-END
\begin{document}

\pagegarde{Produire des écrits utilitaires : le CV}{Français}{Tle Tech}{Français – classe de Terminale technique}{N23}{9 séances (fiche de progression harmonisée)}{Cette leçon apprend à lire, analyser et rédiger un Curriculum Vitae efficace, adapté à une situation d'embauche réelle. Elle mobilise les outils de liaison de la phrase (adverbes de liaison, locutions adverbiales, pronoms relatifs) pour produire un document clair, cohérent et professionnel.
\vspace{4pt}
\begin{multicols}{2}\small
\begin{itemize}
\item À la fin de cette leçon, je sais identifier la structure externe et interne d'un Curriculum Vitae.
\item À la fin de cette leçon, je sais distinguer les rubriques obligatoires et facultatives d'un CV et leur ordre logique.
\item À la fin de cette leçon, je sais repérer et corriger les erreurs fréquentes dans un CV (informations inutiles, fautes, désordre).
\item À la fin de cette leçon, je sais utiliser correctement les adverbes de liaison, les locutions adverbiales et les pronoms relatifs pour lier mes phrases.
\item À la fin de cette leçon, je sais adapter mon CV à une offre d'emploi précise (mise en avant des compétences pertinentes).
\item À la fin de cette leçon, je sais rédiger un CV complet, clair et présenté selon les normes professionnelles.
\end{itemize}\end{multicols}}

\begin{sommaire}
\begin{multicols}{2}
\begin{enumerate}
\item Découverte et réception du CV
\item Structure externe du CV
\item Structure interne du CV : les rubriques
\item Langue : les outils de liaison dans la phrase
\item Rédaction du CV : travail préparatoire
\item Activités d'intégration, compte rendu et remédiation
\item Activité d'intégration
\item Méthodes et exercices résolus
\item Exercices
\item Corrigés des exercices
\item Fiche bilan
\end{enumerate}
\end{multicols}
\end{sommaire}

\begin{objectifs}
\begin{itemize}
\item À la fin de cette leçon, je sais identifier la structure externe et interne d'un Curriculum Vitae.
\item À la fin de cette leçon, je sais distinguer les rubriques obligatoires et facultatives d'un CV et leur ordre logique.
\item À la fin de cette leçon, je sais repérer et corriger les erreurs fréquentes dans un CV (informations inutiles, fautes, désordre).
\item À la fin de cette leçon, je sais utiliser correctement les adverbes de liaison, les locutions adverbiales et les pronoms relatifs pour lier mes phrases.
\item À la fin de cette leçon, je sais adapter mon CV à une offre d'emploi précise (mise en avant des compétences pertinentes).
\item À la fin de cette leçon, je sais rédiger un CV complet, clair et présenté selon les normes professionnelles.
\item À la fin de cette leçon, je sais justifier mes choix de rédaction (ordre des rubriques, vocabulaire, mise en page).
\item À la fin de cette leçon, je sais appliquer ces compétences à des situations concrètes de la vie courante (recherche d'emploi, de stage, de bourse).
\end{itemize}
\end{objectifs}

\begin{prerequis}
\begin{itemize}
\item Les types de textes et les textes fonctionnels (lettre administrative, lettre de motivation) vus en classe de Première
\item La phrase simple et la phrase complexe (classes de 4e et 3e)
\item Les pronoms personnels et possessifs (révision de morphosyntaxe)
\item Le vocabulaire du monde du travail (emploi, stage, recrutement, qualification)
\item La mise en page élémentaire d'un document (titres, paragraphes, alignement)
\end{itemize}
\end{prerequis}

\newpage

\coursec{Découverte et réception du CV}

\courssub{Situation d'accroche et problématique}

À Douala, une entreprise de BTP publie une offre d'emploi de technicien en génie civil. En une semaine, elle reçoit \num{350} candidatures, mais le recruteur ne consacre que 30 secondes à chaque CV avant de décider de le garder ou de l'écarter. Parmi les candidats, deux jeunes techniciens issus du même lycée technique de Yaoundé ont des compétences identiques : pourtant, un seul est convoqué à l'entretien. Qu'est-ce qui a fait la différence ? La forme, la structure et la clarté de son Curriculum Vitae.

\textbf{Problématique :} comment rédiger un CV capable de convaincre un employeur camerounais en moins d'une minute ? Pour répondre à cette question, nous allons d'abord découvrir ce qu'est un CV, à quoi il sert, comment un recruteur le lit, et quels sont ses différents types.

\courssub{Qu'est-ce qu'un Curriculum Vitae ?}

Imaginons que tu postules à un stage dans un atelier de mécanique. Tu ne peux pas raconter toute ta vie au patron : il n'a pas le temps. Il te faut un document court, clair et organisé, qui résume qui tu es, ce que tu as appris et ce que tu sais faire. Ce document, c'est le CV.

\begin{definition}[Curriculum Vitae (CV)]
Le \cle{Curriculum Vitae} (expression latine signifiant \emph{déroulement de la vie}), abrégé \cle{CV}, est un \textbf{document fonctionnel} qui présente de manière \textbf{synthétique et organisée} le parcours scolaire, la formation professionnelle, les compétences et l'expérience d'un candidat à un emploi, à un stage ou à une formation.
\end{definition}

Trois idées sont essentielles dans cette définition :

\begin{itemize}
\item C'est un document \textbf{fonctionnel} : il ne sert pas à raconter, il sert à \emph{agir} (obtenir un emploi).
\item Il est \textbf{synthétique} : il sélectionne et résume l'information, il ne dit pas tout.
\item Il est \textbf{organisé} : l'information est classée en rubriques faciles à repérer.
\end{itemize}

\begin{attention}[Le CV n'est pas une autobiographie]
Erreur fréquente : confondre le CV avec une \textbf{autobiographie détaillée}. Le CV ne raconte pas ta vie du berceau à aujourd'hui ; il ne retient que les informations \emph{utiles pour le poste visé}. Inutile de mentionner ton certificat d'études primaires si tu es en Terminale, ni tes loisirs sans rapport avec l'emploi. Un CV trop long et trop bavard fatigue le recruteur et finit à la corbeille.
\end{attention}

\courssub{Les fonctions du CV}

Pourquoi rédige-t-on un CV ? Le CV remplit trois fonctions complémentaires, qu'il faut toujours avoir à l'esprit en le rédigeant :

\begin{enumerate}
\item \textbf{Informer} l'employeur : le CV donne des renseignements précis et vérifiables (identité, diplômes, expériences, compétences).
\item \textbf{Convaincre} : par sa présentation soignée et la pertinence de son contenu, le CV montre que le candidat est sérieux et correspond au profil recherché.
\item \textbf{Obtenir un entretien d'embauche} : c'est l'objectif final. Le CV ne donne pas le poste ; il ouvre la porte de l'entretien, où le candidat se défendra de vive voix.
\end{enumerate}

\begin{aretenir}[À retenir]
Le CV est une \textbf{porte d'entrée} : sa seule mission est de décrocher un entretien. Tout ce qui ne sert pas cet objectif doit être éliminé.
\end{aretenir}

\begin{savaistu}[30 secondes pour convaincre]
Des enquêtes menées auprès de recruteurs montrent qu'un employeur passe en moyenne \textbf{30 à 40 secondes} sur un CV avant une première sélection. Au Cameroun comme ailleurs, lorsqu'une entreprise reçoit des centaines de dossiers, le premier tri se fait sur la lisibilité et la clarté. Un CV bien structuré, aéré et sans fautes survit à ce premier filtre ; un CV confus est écarté avant même d'être lu.
\end{savaistu}

\courssub{Lecture guidée de deux CV contrastés}

Pour comprendre concrètement ce qui distingue un bon CV d'un mauvais, observons deux candidats fictifs, tous deux titulaires d'un Baccalauréat technique F3 (génie civil), qui répondent à l'offre de l'entreprise de BTP de Douala.

\begin{exemplebox}[Deux CV face à une même offre]
\textbf{CV A — Marlyse N., bien rédigé.} Une page, présentation aérée, rubriques claires (État civil, Formation, Expérience professionnelle, Compétences, Langues, Centres d'intérêt). Les diplômes sont classés du plus récent au plus ancien. Un stage de trois mois chez un entrepreneur de Yaoundé est décrit en deux lignes précises : \emph{« participation au suivi de chantier, lecture de plans, métrés »}. Aucune faute d'orthographe. Adresse électronique simple : \emph{marlyse.n@mail.com}.

\textbf{CV B — Rodrigue E., défectueux.} Trois pages serrées, paragraphes rédigés comme une rédaction : \emph{« Je suis né à Bafoussam où j'ai passé une enfance heureuse\ldots »}. Les diplômes sont mélangés (le BEPC avant le Probatoire), les dates sont absentes, l'expérience de stage est noyée dans un long récit. Plusieurs fautes : \emph{« j'ai effectuer un stage »}. Adresse électronique fantaisiste : \emph{lebossdelaville237@mail.com}. Photo de mauvaise qualité, informations personnelles inutiles (pointure, nombre de frères et sœurs).
\end{exemplebox}

L'observation et la comparaison de ces deux documents permettent de dégager des conclusions nettes, résumées dans le tableau ci-dessous.

\begin{popfigure}
\begin{center}
\rowcolors{2}{popblueL}{white}
\begin{tabularx}{\linewidth}{|l|X|X|}
\hline
\rowcolor{popblue!40} \textbf{Critère} & \textbf{CV A (points forts)} & \textbf{CV B (points faibles)} \\
\hline
Longueur & 1 page, concis & 3 pages, bavard \\
\hline
Forme & Rubriques claires, aéré, lisible & Texte continu, confus \\
\hline
Contenu & Informations classées et datées & Diplômes mélangés, dates absentes \\
\hline
Expérience & Décrite précisément & Noyée dans un récit \\
\hline
Langue & Aucune faute & Fautes de conjugaison \\
\hline
Pertinence & Adapté à l'offre de BTP & Détails inutiles (vie privée) \\
\hline
\end{tabularx}
\end{center}
\legende{Tableau comparatif des deux CV étudiés : le CV A informe et convainc, le CV B fatigue et décourage le recruteur.}
\end{popfigure}

\textbf{Premières conclusions.} À compétences égales, c'est la \textbf{forme} (lisibilité, organisation) et la \textbf{pertinence} du contenu qui font la différence. Le recruteur de Douala a gardé le CV A en quelques secondes et écarté le CV B sans le finir.

\courssub{Comment un recruteur apprécie-t-il un CV ?}

Pour lire et évaluer rapidement un CV — le sien ou celui d'autrui — on utilise une grille d'analyse simple, fondée sur les quatre critères du recruteur : \textbf{lisibilité, pertinence, honnêteté, concision}.

\begin{methode}[Grille d'analyse rapide d'un CV]
\begin{enumerate}
\item \textbf{La forme} : le CV tient-il sur une à deux pages ? Les rubriques sont-elles visibles au premier coup d'œil ? La présentation est-elle aérée et soignée ?
\item \textbf{Le contenu} : les informations essentielles sont-elles présentes (identité, formation, expérience, compétences) ? Sont-elles datées et classées logiquement ?
\item \textbf{L'adéquation à l'offre} : le contenu correspond-il au profil recherché ? Les compétences mentionnées répondent-elles au poste visé ?
\item \textbf{La fiabilité} : les informations sont-elles exactes et vérifiables ? Le CV est-il exempt de fautes d'orthographe et de mensonges ?
\end{enumerate}
\end{methode}

\begin{attention}[L'honnêteté, condition absolue]
Mentir dans un CV (diplôme inventé, expérience exagérée) est une faute grave : le mensonge est presque toujours découvert à l'entretien ou lors de la vérification des pièces, et il disqualifie définitivement le candidat. La concision s'impose aussi : un CV dépasse rarement \textbf{une à deux pages}.
\end{attention}

\begin{exoresolu}[Évaluer un CV avec la grille]
Énoncé. Rodrigue reprend son CV B (voir plus haut) et te demande conseil. En appliquant la grille d'analyse, indique deux défauts à corriger en priorité et justifie ta réponse.
\tcblower
Solution détaillée. 1) \textbf{Forme} : le CV fait trois pages rédigées en paragraphes continus ; il faut le ramener à une page en rubriques courtes, car le recruteur ne consacre que 30 secondes au premier tri. 2) \textbf{Contenu} : les diplômes ne sont ni datés ni classés ; il faut les ordonner du plus récent au plus ancien (Baccalauréat F3, puis Probatoire, puis BEPC) avec les années, car un recruteur vérifie d'abord la formation. Ces deux corrections rendent le CV lisible et crédible, donc capables de survivre au premier tri.
\end{exoresolu}

\courssub{CV et lettre de motivation : deux documents complémentaires}

Le CV ne voyage jamais seul : il est accompagné d'une \textbf{lettre de motivation}. Il ne faut pas confondre leurs rôles.

\begin{itemize}
\item Le \textbf{CV} est un document \textbf{descriptif} : il présente les faits (diplômes, expériences, compétences) sous forme de rubriques brèves.
\item La \textbf{lettre de motivation} est un texte \textbf{argumentatif} : rédigée en phrases, elle explique \emph{pourquoi} le candidat postule et \emph{pourquoi} il correspond au poste.
\end{itemize}

Les deux documents sont \textbf{complémentaires} : le CV dit \emph{ce que} le candidat a fait, la lettre dit \emph{pourquoi} il est le bon choix. Un dossier de candidature complet réunit les deux, sans que l'un répète l'autre.

\courssub{Les différents types de CV}

Selon le parcours du candidat et le poste visé, on choisit l'une des trois présentations suivantes :

\begin{itemize}
\item Le \textbf{CV chronologique} : il présente la formation et l'expérience par ordre de dates, le plus souvent de la plus récente à la plus ancienne (chronologie inversée). C'est le modèle le plus répandu ; il convient aux parcours réguliers.
\item Le \textbf{CV thématique} (ou par compétences) : il regroupe les informations par domaines de compétences (maçonnerie, électricité, bureautique\ldots), indépendamment des dates. Il convient aux candidats polyvalents ou aux parcours interrompus.
\item Le \textbf{CV mixte} : il combine les deux formules, en ouvrant par un résumé des compétences suivi d'une chronologie. C'est le modèle le plus souple.
\end{itemize}

\begin{popfigure}
\begin{center}
\begin{tikzpicture}[
  grow=down,
  level 1/.style={sibling distance=4.6cm, level distance=1.8cm},
  level 2/.style={level distance=1.6cm},
  every node/.style={draw=popblue, fill=popblueL, rounded corners=2pt, align=center, font=\small, text=popdark},
  edge from parent/.style={draw=popblue, thick}
]
\node[fill=popblue!45, font=\small\bfseries] {Les types de CV}
  child { node {Chronologique}
    child { node[fill=popgreenL, draw=popgreen, align=center]{Parcours régulier\\(dates récentes d'abord)} } }
  child { node[align=center]{Thématique\\(par compétences)}
    child { node[fill=popgreenL, draw=popgreen, align=center]{Parcours varié\\ou interrompu} } }
  child { node {Mixte}
    child { node[fill=popgreenL, draw=popgreen, align=center]{Compétences +\\chronologie} } };
\end{tikzpicture}
\end{center}
\legende{Schéma en arbre : les trois types de CV et leurs usages. Le choix dépend du parcours du candidat et du poste visé.}
\end{popfigure}

\begin{exercice}[À toi de jouer]
1. Définis le Curriculum Vitae en une phrase complète, sans recopier la définition du cours.
2. Un ami affirme : « Plus mon CV est long, plus le recruteur sera impressionné. » Corrige cette idée en deux arguments.
3. Un candidat a travaillé dans trois ateliers différents avec des interruptions entre chaque emploi. Quel type de CV lui conseilles-tu ? Justifie.
4. Reprends le CV B de Rodrigue : relève trois informations personnelles inutiles et explique pourquoi il faut les supprimer.
\end{exercice}

\begin{corrige}[Exercice « À toi de jouer »]
1. Le CV est un document fonctionnel qui résume, en rubriques organisées, le parcours, les compétences et l'expérience d'un candidat en vue d'obtenir un entretien.
2. Un recruteur consacre environ 30 secondes au premier tri : un CV long ne sera pas lu. De plus, la concision prouve la capacité à sélectionner l'essentiel, qualité appréciée des employeurs.
3. Le CV thématique (par compétences) : il met en valeur les savoir-faire sans attirer l'attention sur les interruptions de parcours.
4. La pointure, le nombre de frères et sœurs, le récit de l'enfance : ces informations ne servent ni à informer utilement l'employeur ni à convaincre ; elles allongent le CV et nuisent à sa lisibilité.
\end{corrige}

\coursec{Structure externe du CV}

\noindent Après avoir découvert ce qu'est un CV et à quoi il sert dans la section précédente (\emph{Découverte et réception du CV}), nous allons maintenant nous concentrer sur son \textbf{aspect visuel et matériel}. Avant même de lire une seule ligne, le recruteur juge la présentation : un CV soigné, aéré et normé témoigne de votre rigueur, de votre sens de l'organisation et de votre respect des codes professionnels. En Terminale technique, où vous préparez votre insertion (stage, emploi, concours), maîtriser cette structure externe est une compétence transversale indispensable.

\courssub{La présentation matérielle : format, papier, police et titre}

\begin{definition}[Standards matériels du CV professionnel]
Le CV (Curriculum Vitae) est un document normé dont la présentation matérielle obéit à des conventions internationales, adaptées au contexte camerounais :
\begin{itemize}
    \item \textbf{Format :} Papier \textbf{A4} (21 $\times$ 29,7 cm), \textbf{une seule page} (recto seul) pour un jeune diplômé ou étudiant. Une deuxième page n'est tolérée qu'avec une expérience significative (plus de 10 ans).
    \item \textbf{Papier :} Blanc, qualité standard (80--90 g/m²) pour l'impression ; fond blanc pur pour la version numérique (PDF).
    \item \textbf{Police (fonte) :} Une \textbf{seule police sans empattement (sans serif)} pour l'écran (ex. \texttt{Arial}, \texttt{Calibri}, \texttt{Helvetica}) ou \textbf{avec empattement (serif)} pour l'impression (ex. \texttt{Times New Roman}, \texttt{Garamond}). Taille \textbf{11 ou 12 pt} pour le corps de texte ; 14--16 pt pour les titres de rubriques ; 16--18 pt pour le nom du candidat.
    \item \textbf{Couleur :} Noir pour le texte. Une \textbf{seule couleur discrète} (bleu marine, gris anthracite) autorisée pour les titres de rubriques ou les filets de séparation.
\end{itemize}
\end{definition}

\begin{popfigure}
\begin{tikzpicture}[scale=0.75, every node/.style={font=\small}]
    % Page A4 outline
    \draw[popdark, thick] (0,0) rectangle (14.8, 21); % A4 ratio approx 1:1.414, scaled
    % Margins (2cm approx = 1.05 units at this scale)
    \draw[popmuted, dashed, thin] (1.05,1.05) rectangle (13.75, 19.95);
    % Margin labels
    \node[popmuted, rotate=90] at (0.4, 10.5) {Marge gauche 2 cm};
    \node[popmuted] at (7.4, 0.4) {Marge haute 2 cm};
    \node[popmuted, rotate=-90] at (14.4, 10.5) {Marge droite 2 cm};
    \node[popmuted] at (7.4, 20.6) {Marge basse 2 cm};

    % Header Zone
    \draw[popblue, fill=popblueL!50, thick] (1.2, 18.5) rectangle (13.55, 19.7);
    \node[popblue, font=\bfseries] at (7.4, 19.1) {EN-TÊTE : NOM Prénom | Titre du poste visé | « Curriculum Vitae »};

    % Photo Zone (Top Right)
    \draw[poporange, fill=poporangeL!30, thick] (11.8, 16.5) rectangle (13.3, 18.2);
    \node[poporange, align=center] at (12.55, 17.35) {PHOTO\\3,5$\times$4,5 cm};

    % Personal Info Block (Top Left)
    \draw[popgreen, fill=popgreenL!30, thick] (1.2, 16.5) rectangle (11.5, 18.2);
    \node[popgreen, align=left] at (6.35, 17.35) {COORDONNÉES : Adresse, Tél., Email, LinkedIn, Permis, Âge/Date de naissance};

    % Separator line
    \draw[popline, thick] (1.2, 16.2) -- (13.55, 16.2);

    % Section Blocks (Left Column style simulation)
    % Formation
    \draw[poppurple, fill=poppurpleL!20] (1.2, 13.5) rectangle (13.55, 15.9);
    \node[poppurple, font=\bfseries, anchor=west] at (1.4, 15.5) {FORMATION};
    \node[anchor=west] at (1.4, 14.8) {\textbullet\ 2024 -- Baccalauréat Technique Série F4 (Électrotechnique) -- Lycée Technique de Douala};
    \node[anchor=west] at (1.4, 14.2) {\textbullet\ 2021 -- Probatoire Série F4 -- Lycée Technique de Douala};

    % Expérience
    \draw[poppurple, fill=poppurpleL!20] (1.2, 10.5) rectangle (13.55, 13.2);
    \node[poppurple, font=\bfseries, anchor=west] at (1.4, 12.8) {EXPÉRIENCE PROFESSIONNELLE / STAGES};
    \node[anchor=west] at (1.4, 12.1) {\textbullet\ Juil--Août 2023 : Stage maintenance -- ENEO Douala};
    \node[anchor=west] at (1.4, 11.5) {\textbullet\ Juin 2022 : Stage initiation -- SONEL};

    % Compétences
    \draw[poppink, fill=poppinkL!20] (1.2, 7.5) rectangle (13.55, 10.2);
    \node[poppink, font=\bfseries, anchor=west] at (1.4, 9.8) {COMPÉTENCES TECHNIQUES \& LANGUES};
    \node[anchor=west] at (1.4, 9.1) {\textbullet\ Électrotechnique : Câblage, Lecture de plans, Dépannage variateurs};
    \node[anchor=west] at (1.4, 8.5) {\textbullet\ Informatique : Word, Excel, AutoCAD (notions), Python (base)};
    \node[anchor=west] at (1.4, 7.9) {\textbullet\ Langues : Français (maternel), Anglais (B1), Douala (courant)};

    % Centres d'intérêt
    \draw[popgold, fill=popgoldL!20] (1.2, 5.0) rectangle (13.55, 7.2);
    \node[popgold, font=\bfseries, anchor=west] at (1.4, 6.8) {CENTRES D'INTÉRÊT};
    \node[anchor=west] at (1.4, 6.1) {\textbullet\ Football (club), Bricolage électronique, Lecture technique, Bénévolat Croix-Rouge};

    % Signature/Date
    \node[anchor=east, font=\itshape] at (13.3, 3.5) {Fait à Douala, le 15/09/2024};
    \node[anchor=east] at (13.3, 2.5) {Signature};

    % Annotations arrows
    \draw[popdark, <->, thick] (1.2, 19.9) -- node[above, font=\tiny] {Largeur utile $\approx$ 17 cm} (13.55, 19.9);
    \draw[popdark, <->, thick] (0.6, 1.05) -- node[left, font=\tiny, rotate=90] {Hauteur utile $\approx$ 25,7 cm} (0.6, 19.95);
\end{tikzpicture}
\legende{Schéma annoté de la structure externe d'un CV standard (format A4, marges 2 cm, organisation en blocs verticaux). Les zones sont identifiées par couleur : en-tête (bleu), photo (orange), coordonnées (vert), rubriques principales (violet/rose/or).}
\end{popfigure}

\begin{methode}[Les 5 règles d'or pour réussir la mise en page d'un CV]
\begin{enumerate}
    \item \textbf{Respecter les marges :} 2 cm minimum de chaque côté (haut, bas, gauche, droite) pour garantir la lisibilité et l'impression sans rognage.
    \item \textbf{Unifier la police :} Une seule famille de police pour tout le document. Jouer sur les \textbf{graisses} (Gras, Normal, Light) et les \textbf{tailles} (11, 12, 14, 16 pt) pour hiérarchiser, jamais sur le changement de police.
    \item \textbf{Aérer par les blancs :} Laisser un interligne de 1,15 ou 1,5 dans les paragraphes ; sauter une ligne (ou 6--8 pt) entre chaque rubrique ; espacer les items d'une même liste.
    \item \textbf{Aligner rigoureusement :} Choisir un alignement (gauche pour le texte, droite pour les dates) et le tenir sur toute la page. Utiliser les tabulations ou un tableau invisible, jamais la barre d'espace.
    \item \textbf{Exporter en PDF :} Seul le format PDF fige la mise en page (polices, marges, photo) quel que soit l'ordinateur du recruteur. Tester l'impression avant envoi.
\end{enumerate}
\end{methode}

\courssub{L'aération du document : marges, espaces, alignements et puces}

L'\textbf{aération} (ou « espace blanc ») n'est pas du vide perdu : c'est un \textbf{outil de lecture}. L'œil humain repère d'abord les masses sombres (texte) et les masses claires (marges). Un CV « dense » fatigue l'œil en 6 secondes (temps moyen de lecture diagonale) ; un CV « aéré » guide le regard vers l'essentiel.

\begin{exemplebox}[Modèle de pagination recommandé (standard MINESEC / Administration camerounaise)]
\begin{center}
\begin{tabularx}{\linewidth}{|l|X|}\hline
\textbf{Paramètre} & \textbf{Valeur recommandée} \\ \hline
Marges (Haut, Bas, Gauche, Droite) & 2,0 cm (minimum 1,5 cm) \\ \hline
Police principale & Arial 11 pt ou Times New Roman 12 pt \\ \hline
Interligne (corps de texte) & 1,15 (simple) ou 1,5 (très lisible) \\ \hline
Espacement avant/après rubrique & 6 pt / 12 pt (équivaut à une demi-ligne / une ligne) \\ \hline
Alignement texte & Justifié (si largeur colonne > 8 cm) ou Gauche (recommandé) \\ \hline
Alignement dates & Droite (colonne étroite) ou Gauche (aligné sur texte) \\ \hline
Puces & \textbullet\ (rond plein) ou -- (tiret) ou $\rightarrow$ (flèche simple), \textbf{identiques partout} \\ \hline
Trait de séparation rubriques & Filet fin 0,5 pt, couleur gris 50\% ou bleu marine, largeur 100\% ou 80\% \\ \hline
\end{tabularx}
\end{center}
\end{exemplebox}

\begin{popfigure}
\begin{tikzpicture}[scale=0.6, every node/.style={font=\tiny}]
    % --- MAQUETTE AÉRÉE (BONNE) ---
    \begin{scope}[xshift=0cm]
        \draw[popgreen, thick] (0,0) rectangle (7, 10);
        \node[popgreen, font=\bfseries] at (3.5, 10.3) {MAQUETTE AÉRÉE (Recommandée)};
        % Margins visual
        \draw[popgreen!50, dashed] (0.4,0.4) rectangle (6.6, 9.6);
        % Header
        \draw[popblue, fill=popblueL] (0.5, 9.0) rectangle (6.5, 9.6);
        \node at (3.5, 9.3) {NOM Prénom \hfill Titre Poste};
        % Photo
        \draw[poporange, fill=poporangeL] (5.2, 7.8) rectangle (6.3, 9.0);
        \node at (5.75, 8.4) {Photo};
        % Coord
        \draw[popgreen, fill=popgreenL] (0.5, 7.8) rectangle (5.0, 9.0);
        \node[align=left] at (2.75, 8.4) {Adresse \newline Tel \newline Email};
        % Separator
        \draw[popline] (0.5, 7.6) -- (6.5, 7.6);
        % Formation
        \draw[poppurple, fill=poppurpleL] (0.5, 6.0) rectangle (6.5, 7.4);
        \node[align=left, anchor=north west] at (0.6, 7.3) {\textbf{FORMATION}\newline \textbullet\ Bac Tech 2024\newline \textbullet\ Probatoire 2021};
        % Expérience
        \draw[poppurple, fill=poppurpleL] (0.5, 4.0) rectangle (6.5, 5.8);
        \node[align=left, anchor=north west] at (0.6, 5.7) {\textbf{EXPÉRIENCE}\newline \textbullet\ Stage ENEO 2023\newline \textbullet\ Stage SONEL 2022};
        % Compétences
        \draw[poppink, fill=poppinkL] (0.5, 2.0) rectangle (6.5, 3.8);
        \node[align=left, anchor=north west] at (0.6, 3.7) {\textbf{COMPÉTENCES}\newline \textbullet\ Technique, IT, Langues};
        % Intérêts
        \draw[popgold, fill=popgoldL] (0.5, 0.5) rectangle (6.5, 1.8);
        \node[align=left, anchor=north west] at (0.6, 1.7) {\textbf{INTÉRÊTS}\newline \textbullet\ Sport, Bénévolat};
    \end{scope}

    % --- MAQUETTE SURCHARGÉE (MAUVAISE) ---
    \begin{scope}[xshift=8.5cm]
        \draw[popink, thick] (0,0) rectangle (7, 10);
        \node[popink, font=\bfseries] at (3.5, 10.3) {MAQUETTE SURCHARGÉE (À éviter)};
        % Tiny margins
        \draw[popink!50, dashed] (0.15,0.15) rectangle (6.85, 9.85);
        % Header cramped
        \draw[popblue, fill=popblueL] (0.2, 9.5) rectangle (6.8, 9.85);
        \node[font=\tiny] at (3.5, 9.65) {NOM Prénom Titre Poste Adresse Tel Email Tout Collé};
        % No photo space or overlapping
        % Text blocks cramped, tiny font simulated by tight packing
        \node[align=left, anchor=north west, font=\tiny, text width=6.4cm] at (0.2, 9.3) {
            \textbf{FORMATION} Bac Tech 2024 Lycée Technique Douala Probatoire 2021 Lycée Technique Douala
            \textbf{EXPÉRIENCE} Stage ENEO Juillet Août 2023 Maintenance Stage SONEL Juin 2022 Initiation
            \textbf{COMPÉTENCES} Électrotechnique Câblage Plans Variateurs Informatique Word Excel AutoCAD Python Langues Français Anglais Douala
            \textbf{INTÉRÊTS} Football Club Bricolage Lecture Bénévolat Croix-Rouge
        };
    \end{scope}
\end{tikzpicture}
\legende{Comparaison visuelle : à gauche, une mise en page aérée (marges respectées, interlignes, puces, séparation nette des rubriques) ; à droite, une mise en page surchargée (marges rognées, pas d'interligne, texte compact, illisible en lecture diagonale).}
\end{popfigure}

\begin{attention}[Erreur fréquente : Polices multiples, couleurs fantaisistes et « design » amateur]
\begin{itemize}
    \item \textbf{Changer de police selon les rubriques :} Ex. \texttt{Arial} pour le nom, \texttt{Times} pour la formation, \texttt{Comic Sans} pour les loisirs. \textbf{Conséquence :} Aspect brouillon, manque de cohérence, rejet immédiat par les ATS (logiciels de tri automatique).
    \item \textbf{Utiliser des couleurs vives (rose, vert fluo, rouge) :} \textbf{Conséquence :} Illisible en impression noir \& blanc (photocopie administration), perçu comme non professionnel.
    \item \textbf{Mettre un fond d'image, des cadres décoratifs, des ombres portées :} \textbf{Conséquence :} Alourdit le fichier PDF, perturbe la lecture OCR/ATS, distrait le recruteur.
    \item \textbf{Justifier le texte dans une colonne étroite (ex. 5 cm) :} Crée des « rivières » d'espaces blancs horribles à l'intérieur des lignes. \textbf{Solution :} Alignement à gauche (drapeau) pour les colonnes étroites.
\end{itemize}
\textbf{Règle d'or :} La sobriété \textbf{est} l'élégance technique. Le contenu prime sur le contenant, mais le contenant doit être invisible (transparent à la lecture).
\end{attention}

\courssub{Le titre du document et la photo d'identité}

\begin{propriete}[L'en-tête : identification immédiate]
La zone haute (premiers 3--4 cm) doit permettre d'identifier le candidat et l'objet du CV en \textbf{moins de 2 secondes}.
\begin{enumerate}
    \item \textbf{Nom PRÉNOM} : En très gros (16--18 pt), \textbf{Gras}, MAJUSCULES pour le nom, Première lettre pour le prénom (ex. \textsc{MBARGA} Jean-Pierre). Centré ou aligné à gauche.
    \item \textbf{Titre du poste visé / Spécialité :} Juste en dessous, taille 14 pt, couleur discrète (ex. \emph{Technicien de Maintenance Électrique -- Spéc. Automatismes}). Cela remplace l'ancienne mention « Curriculum Vitae » souvent jugée redondante.
    \item \textbf{Mention « Curriculum Vitae » :} Facultative, en petit (10 pt, gris) si on tient à la tradition, alignée à droite ou centrée sous le nom.
\end{enumerate}
\end{propriete}

\begin{definition}[La photo d'identité : normes et insertion]
Au Cameroun, comme dans la plupart des pays francophones d'Afrique, \textbf{la photo est attendue, voire exigée} (concours publics, grandes entreprises). Elle n'est pas obligatoire en France (loi anti-discrimination) mais reste la norme en entreprise technique locale.
\begin{itemize}
    \item \textbf{Format :} 3,5 $\times$ 4,5 cm (standard passeport/visa).
    \item \textbf{Qualité :} Fond uni clair (blanc, bleu clair, gris clair), netteté parfaite, visage dégagé, regard vers l'objectif, expression neutre mais souriante (professionnelle).
    \item \textbf{Récence :} Moins de \textbf{6 à 12 mois}.
    \item \textbf{Position :} En \textbf{haut à droite} (classique) ou \textbf{haut à gauche} (si design moderne centré). Jamais au milieu, jamais en bas.
    \item \textbf{Insertion technique :} Image insérée \emph{dans} le document (pas en objet flottant instable), recadrée, compressée (fichier final < 500 Ko).
\end{itemize}
\end{definition}

\begin{savaistu}[Le CV dans la Fonction Publique et les Concours Professionnels au Cameroun]
Dans le cadre des \textbf{concours professionnels} (accès aux corps de l'administration, ENSET, ENSAI, concours d'entrée à l'ENAM, recrutement SNI, CAMTEL, SONARA, etc.), le CV est une \textbf{pièce obligatoire du dossier de candidature}. Il est souvent exigé \textbf{manuscrit} (sur papier à en-tête fourni ou papier blanc) \textbf{ET} sous format numérique.
\begin{itemize}
    \item La photo d'identité \textbf{agrafée ou collée} en haut à droite est impérative.
    \item La signature manuscrite en bas de page est \textbf{obligatoire} (valide l'authenticité des informations).
    \item La mention « Lu et approuvé, bon pour service » avec la date est parfois demandée.
    \item Le format « Europass » est de plus en plus accepté mais le format « classique » structuré par rubriques (Formation, Expérience, Compétences) reste le standard attendu par les jurys camerounais.
\end{itemize}
\end{savaistu}

\courssub{L'organisation en rubriques : mise en page en colonnes vs blocs}

Le choix de la structure visuelle (architecture de la page) conditionne la rapidité de lecture. Deux modèles dominent :

\begin{center}
\begin{tabularx}{\linewidth}{|l|X|X|}\hline
\rowcolor{popblueL}
\textbf{Critère} & \textbf{Mise en page en BLOCS (Verticale / Linéaire)} & \textbf{Mise en page en COLONNES (Latérale / Sidebar)} \\ \hline
\textbf{Principe} & Rubriques les unes sous les autres, largeur page. & Colonne étroite à gauche (30\%) : Coord, Compétences, Intérêts. Colonne large à droite (70\%) : Formation, Expérience. \\ \hline
\textbf{Avantages} & \textbf{Standard universel.} Lu par tous les ATS. Facile à modifier (Word). Idéal pour peu d'expérience. & \textbf{Moderne, visuel.} Met en valeur les compétences (mots-clés) dès le premier regard. Gain de place vertical. \\ \hline
\textbf{Inconvénients} & Peut sembler « vide » si peu d'expériences. Moins « design ». & \textbf{Risque ATS :} Les logiciels lisent mal les colonnes (mélangent colonne gauche/droite). Difficile à éditer sans casser la mise en page. \\ \hline
\textbf{Recommandation Tle Tech} & \textbf{PRIVILÉGIER CE MODÈLE.} Sûr, professionnel, conforme aux attentes administration/entreprises locales. & À réserver pour portefeuille numérique (LinkedIn, site perso) ou secteurs créatifs/IT \textbf{si} on maîtrise LaTeX/InDesign/Canva et qu'on teste l'ATS. \\ \hline
\end{tabularx}
\end{center}

\begin{aretenir}[Synthèse : Check-list de la structure externe (à valider avant envoi)]
\begin{itemize}
    \item [$\square$] Format A4, 1 page, PDF, marges 2 cm.
    \item [$\square$] Police unique (Arial 11 ou Times 12), noir, une couleur max pour titres.
    \item [$\square$] En-tête : NOM Prénom + Titre poste visé (visible immédiatement).
    \item [$\square$] Photo pro récente (3,5$\times$4,5 cm), coin haut droit/gauche, fond neutre.
    \item [$\square$] Coordonnées complètes (Adresse, Tél., Email pro, LinkedIn) regroupées.
    \item [$\square$] Rubriques distinctes, titrées en Gras 14 pt, séparées par filet fin ou espace blanc.
    \item [$\square$] Alignement cohérent (dates à droite ou gauche, texte justifié/gauche).
    \item [$\square$] Puces uniformes, interligne 1,15, espacement rubriques 6--12 pt.
    \item [$\square$] Signature + Date (bas de page, droite) si contexte administratif/camerounais.
    \item [$\square$] Fichier nommé : \texttt{NOM\_Prenom\_CV\_PosteVisé.pdf} (pas \texttt{CV\_final\_v2.pdf}).
\end{itemize}
\end{aretenir}

\courssub{La signature et la date : usage contextuel}

\begin{experience}[Observation : La signature manuscrite, marque d'authenticité au Cameroun]
\textbf{Contexte :} Dépôt de dossier physique (concours MINFOPRA, recrutement direct, stage conventionné).
\textbf{Constat :} Un CV non signé est souvent considéré comme « incomplet » ou « non engagé » par les secrétariats et jurys.
\textbf{Pratique recommandée :}
\begin{enumerate}
    \item Imprimer le CV sur papier qualité.
    \item Signer \textbf{manuscrite} en bas à droite (prénom + nom lisible, pas un gribouillage).
    \item Ajouter la mention manuscrite : \emph{« Fait à [Ville], le [Date du jour] »} au-dessus de la signature.
    \item Scanner l'original signé pour la version numérique (PDF) si envoi par email.
\end{enumerate}
\textbf{Note :} Pour une candidature 100\% dématérialisée (plateforme emploi, email direct DRH privée), une signature \textbf{numérique} (image de la signature scannée insérée en bas du CV) est acceptée et professionnelle. La date reste indispensable (fraîcheur du document).
\end{experience}

\begin{exoresolu}[Analyse critique de deux mises en page externes]
\textbf{Énoncé :} Voici deux extraits de mise en page (simulés). Identifiez 5 défauts majeurs dans le CV A et 3 points forts dans le CV B.

\textbf{CV A (Extrait) :}
\begin{verbatim}
Marges : 0,5 cm partout. Police : Titre "Comic Sans MS 18pt Rose", Corps "Times New Roman 10pt", Rubriques "Arial 14pt Bleu", Compétences "Calibri 11pt Vert".
Interligne : Simple (1.0) partout, pas d'espace entre rubriques.
Photo : Photo de vacances recadrée (fond plage), 5x5 cm, milieu de page.
En-tête : "CURRICULUM VITAE" seul en gros au centre. Nom petit en bas à gauche.
Alignement : Dates dispersées (certaines gauche, certaines centre, certaines droite). Puces : *, -, >, • mélangées.
Fichier : "Mon_CV_definitif_vraiment_final.docx" envoyé par mail.
\end{verbatim}

\textbf{CV B (Extrait) :}
\begin{verbatim}
Marges : 2 cm. Police : Arial 11 pt unique. Titres rubriques : Arial 14 pt Gras Bleu Marine.
Interligne : 1.15. Espace 8pt entre rubriques. Filet gris 0.5pt sous chaque titre.
Photo : Pro, fond blanc, 3.5x4.5 cm, coin haut droit.
En-tête : "MBARGA Jean-Pierre" (16pt Gras) / "Technicien Maintenance Industrielle" (13pt).
Alignement : Dates alignées à droite (colonne virtuelle). Puces : "•" uniquement.
Fichier : "MBARGA_Jean-Pierre_CV_Technicien_Maintenance.pdf" (350 Ko).
\end{verbatim}

\tcblower
\textbf{Solution détaillée :}

\textbf{CV A -- 5 défauts majeurs (éliminatoires) :}
\begin{enumerate}
    \item \textbf{Marges trop fines (0,5 cm) :} Risque de rognage à l'impression, sensation d'étouffement, non-conforme aux standards admin.
    \item \textbf{4 polices différentes + couleurs fantaisistes :} \texttt{Comic Sans} (interdit pro), \texttt{Times}, \texttt{Arial}, \texttt{Calibri}. Couleurs rose/bleu/vert. $\rightarrow$ Incohérence visuelle totale, illisible en N\&B, rejetté par ATS.
    \item \textbf{Photo non professionnelle :} Fond plage, format carré, position centrale. $\rightarrow$ Décrédibilise le candidat (manque de sérieux).
    \item \textbf{En-tête inefficace :} « Curriculum Vitae » ne dit pas \emph{qui} ni \emph{quoi}. Nom caché.
    \item \textbf{Alignement chaotique + puces mixtes + format .docx :} Instabilité de la mise en page selon la version de Word du destinataire. Nom de fichier non professionnel.
\end{enumerate}

\textbf{CV B -- 3 points forts (modèle à suivre) :}
\begin{enumerate}
    \item \textbf{Cohérence graphique parfaite :} Police unique (Arial), hiérarchie par graisse/taille/couleur unique (Bleu Marine), marges standard, interligne lisible.
    \item \textbf{En-tête fonctionnelle :} Nom gros + Titre poste visé = identification immédiate. Photo pro bien placée.
    \item \textbf{Rigueur technique :} Alignement dates (colonne virtuelle), puce unique, fichier PDF léger nommé correctement. Prêt pour ATS et impression.
\end{enumerate}
\end{exoresolu}

\begin{exercice}[Application : Audit de votre propre CV]
À partir de la check-list de l'\texttt{aretenir} ci-dessus, prenez votre CV actuel (ou créez-en un brouillon) et cochez chaque critère.
\begin{enumerate}
    \item Combien de critères sont validés ? Combien restent à corriger ?
    \item Identifiez les 3 corrections les plus urgentes à apporter.
    \item Renommez votre fichier selon la convention \texttt{NOM\_Prenom\_CV\_PosteVisé.pdf}.
\end{enumerate}
\textit{Travail à faire en autonomie, à présenter lors de la prochaine séance pour relecture par les pairs.}
\end{exercice}

\noindent Cette maîtrise de la structure externe constitue le \textbf{socle visuel} de votre candidature. La section suivante (\emph{Structure interne du CV : les rubriques}) détaillera le contenu précis de chaque zone identifiée ici (Formation, Expérience, Compétences, etc.) et la rédaction des items qui les composent.

\coursec{Structure interne du CV : les rubriques}

Dans la section précédente, nous avons étudié la \cle{structure externe} du CV : sa présentation matérielle (marges régulières, police sobre et lisible, longueur d'une page, mise en page aérée, photo facultative). Mais un beau contenant ne suffit pas : c'est le \emph{contenu} qui convainc un recruteur. Ce contenu n'est pas jeté en vrac sur la page : il est organisé en blocs thématiques que l'on appelle des \cle{rubriques}. C'est l'objet de cette section : identifier chaque rubrique, savoir ce qu'elle doit contenir, ce qu'elle ne doit pas contenir, et dans quel ordre les disposer selon son profil.

\courssub{Notion de rubrique et ordre antichronologique}

\begin{definition}[Rubrique]
Une \cle{rubrique} est une partie clairement délimitée du CV, introduite par un titre visible (en gras, éventuellement en couleur), qui regroupe des informations de même nature : état civil, formation, expérience professionnelle, compétences, centres d'intérêt, références. Chaque rubrique répond à une question précise que se pose le recruteur : « Qui est ce candidat ? Que sait-il faire ? Où l'a-t-il appris ? Où l'a-t-il déjà fait ? »
\end{definition}

\begin{definition}[Ordre antichronologique]
L'\cle{ordre antichronologique} (ou chronologique inverse) consiste à présenter les éléments \textbf{du plus récent au plus ancien}. Dans un CV, on cite donc d'abord le dernier diplôme obtenu ou le dernier poste occupé, puis on remonte le temps. Cet ordre est la norme car le recruteur s'intéresse d'abord à ce que le candidat est et fait \emph{aujourd'hui}.
\end{definition}

\begin{attention}[Piège de rédaction]
Ne mélangez jamais les deux ordres dans une même rubrique : si vous commencez par le diplôme le plus récent, continuez ainsi jusqu'au plus ancien. Une chronologie désordonnée donne l'impression d'un candidat négligent. Veillez aussi à l'\textbf{homogénéité de la présentation} : mêmes tirets, mêmes formats de dates (par exemple toujours « juillet-septembre 2025 », jamais « 07/2025 » ailleurs).
\end{attention}

\courssub{L'état civil et les informations personnelles}

C'est la première rubrique, placée en haut du CV. Elle permet au recruteur d'identifier le candidat et de le \textbf{contacter facilement}. Elle comprend :

\begin{itemize}
\item le \textbf{nom} (en majuscules) et le \textbf{prénom} ;
\item la \textbf{date et le lieu de naissance} ;
\item la \textbf{nationalité} ;
\item les \textbf{contacts} : numéro de téléphone, adresse e-mail, adresse postale ou quartier de résidence ;
\item la \textbf{situation matrimoniale} : célibataire, marié(e)... Cette mention est \emph{facultative} ; on peut l'omettre sans inconvénient.
\end{itemize}

\begin{exemplebox}[État civil correctement rédigé]
\textbf{MBALLA NGOUESSA Jean-Paul}\\
Né le 12 mars 2005 à Bafoussam — Nationalité : camerounaise\\
Tél. : 6 90 00 00 00 — E-mail : jpmballa@gmail.com\\
Yaoundé, quartier Nkolbisson
\end{exemplebox}

\begin{attention}[Informations à proscrire]
Une erreur fréquente consiste à mentionner des informations \textbf{discriminatoires ou inutiles} : la religion, l'ethnie ou la région d'origine, le nombre d'enfants, l'appartenance politique, le groupe sanguin. Ces éléments ne concernent pas la compétence professionnelle et peuvent exposer le candidat à des préjugés. Un bon CV ne contient que ce qui sert la candidature. De même, évitez les adresses e-mail fantaisistes du type « beaugosse2000@... » : créez une adresse sobre, formée de votre nom et prénom.
\end{attention}

\courssub{L'objectif professionnel}

Placé juste sous l'état civil, l'\cle{objectif professionnel} est une rubrique \textbf{complémentaire mais utile} : en une ligne, le candidat annonce le poste ou le type d'emploi recherché. Elle montre au recruteur que le CV n'est pas envoyé au hasard et qu'il correspond à l'offre d'emploi visée.

\begin{exemplebox}[Objectif professionnel]
« Technicien de maintenance industrielle — jeune diplômé du Baccalauréat technique, série F3, à la recherche d'un premier emploi dans une entreprise de production. »
\end{exemplebox}

\courssub{La formation et les diplômes}

Cette rubrique liste les diplômes et formations suivies, dans l'\textbf{ordre antichronologique}. Pour chaque élément, on précise : l'\textbf{année d'obtention}, l'\textbf{intitulé exact du diplôme}, l'\textbf{établissement} et la \textbf{ville}, et éventuellement la \textbf{mention} ou un résultat valorisant.

\begin{exemplebox}[Rubrique « Formation »]
\textbf{Formation}\\
2025 : Baccalauréat technique, série F4 (Génie civil), Lycée technique de Douala-Bonabéri — mention Bien.\\
2023 : Probatoire technique F4, Lycée technique de Douala-Bonabéri.\\
2022 : BEPC, Collège Saint-Joseph de Deïdo.
\end{exemplebox}

\courssub{L'expérience professionnelle}

C'est souvent la rubrique la plus regardée par le recruteur. Elle regroupe les \textbf{postes occupés}, mais aussi les \textbf{stages}, les travaux saisonniers et même les activités associatives à responsabilité. Pour le jeune diplômé qui n'a pas encore travaillé, les \cle{stages} scolaires et académiques jouent ce rôle : il ne faut jamais laisser cette rubrique vide.

Chaque expérience se présente sur une ligne (ou deux) selon la formule : \textbf{poste + entreprise + lieu + dates + mission principale}.

\begin{methode}[Rédiger une expérience professionnelle en une ligne efficace]
\begin{enumerate}
\item Commencer par l'\textbf{intitulé du poste} ou de la fonction (ex. : « Stagiaire technicien »).
\item Nommer l'\textbf{entreprise} et le \textbf{lieu} (ex. : « SABC, Douala »).
\item Indiquer les \textbf{dates ou la durée} (ex. : « juillet-septembre 2025 »).
\item Formuler la \textbf{mission principale} avec un verbe d'action à l'infinitif ou au participe (ex. : « maintenance des lignes de production »).
\item Vérifier que la ligne tient en une phrase courte, sans « je ».
\end{enumerate}
\end{methode}

\begin{exemplebox}[Expérience professionnelle modèle]
\textbf{Expérience professionnelle}\\
Juillet-septembre 2025 : Stagiaire technicien, SABC, Douala — maintenance des lignes de production et contrôle qualité.\\
Juillet 2024 : Stage d'observation, Atelier mécanique Eto'o, Bafoussam — révision de moteurs et accueil de la clientèle.
\end{exemplebox}

\courssub{Les compétences}

Cette rubrique distingue deux familles de compétences :

\begin{itemize}
\item les \textbf{compétences techniques} : maîtrise de l'outil informatique (Word, Excel, logiciels de dessin technique), permis de conduire, gestes et techniques propres au métier (soudure, câblage, saisie comptable, conduite d'engins...) ;
\item les \textbf{compétences linguistiques} : français, anglais (préciser le niveau : lu, écrit, parlé ; courant, scolaire), et les langues nationales pratiquées.
\end{itemize}

\begin{savaistu}[Le bilinguisme, un atout majeur]
Au Cameroun, pays officiellement bilingue, la maîtrise du \textbf{français et de l'anglais} est un atout décisif sur le marché de l'emploi : les entreprises nationales et internationales recherchent des techniciens capables de lire une documentation technique en anglais et de rédiger un rapport en français. Mentionner honnêtement son niveau dans les deux langues peut faire la différence entre deux candidatures équivalentes.
\end{savaistu}

\courssub{Les centres d'intérêt et les références}

Les \textbf{centres d'intérêt} (ou activités extra-professionnelles) humanisent le CV : sport pratiqué, lecture, engagement associatif, musique. La règle est la \cle{pertinence} et la \cle{sobriété} : deux ou trois éléments suffisent, de préférence ceux qui révèlent une qualité utile au poste (le football collectif témoigne de l'esprit d'équipe, la présidence d'un club de l'esprit d'initiative).

Les \textbf{références} désignent des personnes (anciens chefs, enseignants) que le recruteur peut contacter pour vérifier les informations. Cette rubrique est \emph{facultative} : on peut se contenter de la formule consacrée « Références disponibles sur demande ».

\begin{center}
\begin{tabularx}{\linewidth}{|l|X|X|}\hline
\rowcolor{popblueL} \textbf{Rubriques obligatoires} & \textbf{Rubriques facultatives} & \textbf{À proscrire} \\ \hline
État civil (nom, contacts) & Situation matrimoniale & Religion, ethnie, région d'origine \\ \hline
Formation / diplômes & Objectif professionnel & Nombre d'enfants, appartenance politique \\ \hline
Expérience professionnelle (stages inclus) & Centres d'intérêt & Faux diplômes, dates inventées \\ \hline
Compétences (techniques et linguistiques) & Références & Adresses e-mail fantaisistes \\ \hline
\end{tabularx}
\end{center}

\courssub{L'ordre logique des rubriques selon le profil}

Toutes les rubriques ne se valent pas selon le candidat. Le principe est simple : \textbf{placer en premier ce qui constitue votre principal atout}.

\begin{itemize}
\item Le \textbf{jeune diplômé} (cas typique de l'élève de Terminale technique) n'a pas encore d'expérience : sa force, c'est sa \textbf{formation}. Il placera donc la rubrique « Formation » immédiatement après l'état civil, avant l'expérience (stages).
\item Le \textbf{candidat expérimenté} met en avant son parcours professionnel : la rubrique « Expérience professionnelle » vient en premier, la formation passe après.
\end{itemize}

\begin{popfigure}
\begin{center}
\begin{tikzpicture}[node distance=0.35cm,
bloc/.style={draw=popblue, fill=popblueL, rounded corners=2pt, text width=9.2cm, align=left, inner sep=4pt, font=\small},
fleche/.style={-{Stealth[length=2.5mm]}, thick, poporange}]
\node[bloc] (b1) {\textbf{1. État civil et contacts} (nom, téléphone, e-mail)} ;
\node[bloc, below=of b1] (b2) {\textbf{2. Objectif professionnel} (une ligne)} ;
\node[bloc, below=of b2] (b3) {\textbf{3. Formation / diplômes} (ordre antichronologique) \emph{— l'atout du jeune diplômé}} ;
\node[bloc, below=of b3] (b4) {\textbf{4. Expérience professionnelle} (stages, petits boulots)} ;
\node[bloc, below=of b4] (b5) {\textbf{5. Compétences} (techniques et linguistiques)} ;
\node[bloc, below=of b5] (b6) {\textbf{6. Centres d'intérêt} (2 ou 3 éléments sobres)} ;
\node[bloc, below=of b6] (b7) {\textbf{7. Références} (« disponibles sur demande »)} ;
\draw[fleche] (b1) -- (b2) ; \draw[fleche] (b2) -- (b3) ;
\draw[fleche] (b3) -- (b4) ; \draw[fleche] (b4) -- (b5) ;
\draw[fleche] (b5) -- (b6) ; \draw[fleche] (b6) -- (b7) ;
\end{tikzpicture}
\end{center}
\legende{Ordre des rubriques pour un jeune diplômé technique camerounais : la formation passe avant l'expérience. Pour un candidat expérimenté, les blocs 3 et 4 s'inversent.}
\end{popfigure}

\begin{exoresolu}[Mettre en ordre les rubriques d'un CV]
Énoncé : Aïcha, 19 ans, vient d'obtenir son Baccalauréat technique F2 (Électronique) au Lycée technique de Yaoundé. Elle a effectué un stage de deux mois à Camtel et maîtrise Word, Excel et l'anglais scolaire. Elle a rédigé son CV dans cet ordre : centres d'intérêt, expérience, état civil, compétences, formation. Corrigez l'ordre de ses rubriques et justifiez.
\tcblower
Solution détaillée : Aïcha est une \textbf{jeune diplômée} : son principal atout est sa formation, qui doit donc précéder l'expérience. L'état civil se place toujours en tête, et les centres d'intérêt à la fin. L'ordre correct est : \textbf{1.} État civil et contacts ; \textbf{2.} (éventuellement) objectif professionnel ; \textbf{3.} Formation (Bac F2, 2025, Lycée technique de Yaoundé) ; \textbf{4.} Expérience professionnelle (stage de deux mois à Camtel, avec dates et mission) ; \textbf{5.} Compétences (Word, Excel ; anglais scolaire ; français courant) ; \textbf{6.} Centres d'intérêt ; \textbf{7.} Références disponibles sur demande.
\end{exoresolu}

\begin{aretenir}[L'essentiel]
Un CV est composé de \cle{rubriques} : état civil, formation, expérience professionnelle, compétences, centres d'intérêt, références. Les diplômes et les expériences se présentent dans l'\cle{ordre antichronologique} (du plus récent au plus ancien). Chaque expérience se rédige selon la formule « poste + entreprise + lieu + dates + mission ». Le jeune diplômé place la formation en premier ; le candidat expérimenté place l'expérience en premier. On proscrit toute information discriminatoire (religion, ethnie, nombre d'enfants).
\end{aretenir}

\begin{exercice}[À vous de jouer]
Classez les informations suivantes dans les bonnes rubriques, puis rédigez la ligne d'expérience correspondante : « permis de conduire catégorie B » ; « stage à la SABC de Douala, juillet-septembre 2025, maintenance des lignes de production » ; « né le 5 janvier 2006 à Bafang » ; « capitaine de l'équipe de handball du lycée » ; « anglais : niveau scolaire ».
\end{exercice}

\begin{corrige}[Exercice]
État civil : né le 5 janvier 2006 à Bafang. — Expérience professionnelle : « Juillet-septembre 2025 : Stagiaire technicien, SABC, Douala — maintenance des lignes de production. » — Compétences : permis de conduire catégorie B ; anglais : niveau scolaire. — Centres d'intérêt : handball (capitaine de l'équipe du lycée), ce qui témoigne d'un esprit d'équipe et de responsabilité.
\end{corrige}

\coursec{Langue : les outils de liaison dans la phrase}

Dans la section précédente, nous avons appris à organiser les rubriques du CV : état civil, formation, expériences professionnelles, compétences, centres d'intérêt. Mais un CV, comme tout écrit utilitaire (lettre de motivation, demande d'emploi, compte rendu), ne se réduit pas à une liste d'informations juxtaposées. Pour que le lecteur — l'employeur — comprenne votre parcours d'un seul coup d'œil, les phrases doivent être \cle{reliées} entre elles. C'est le rôle des \cle{outils de liaison} : adverbes de liaison, locutions adverbiales et pronoms relatifs. Cette section de langue vous donne les moyens grammaticaux de produire un texte fluide, cohérent et professionnel.

\courssub{Pourquoi relier les phrases ? La cohérence textuelle}

Lisons ces deux versions d'un même extrait de lettre de motivation :

\begin{exemplebox}[Deux versions d'un même parcours]
\textbf{Version A (phrases juxtaposées) :} « J'ai obtenu mon Probatoire technique en 2024. J'ai effectué un stage à la SODECOTON. J'ai appris la maintenance industrielle. J'ai obtenu le permis C. Je sollicite un emploi dans votre entreprise. »

\textbf{Version B (phrases liées) :} « Après l'obtention de mon Probatoire technique en 2024, j'ai effectué un stage à la SODECOTON, \cle{où} j'ai appris la maintenance industrielle ; \cle{de plus}, j'ai obtenu le permis C. \cle{C'est pourquoi} je sollicite aujourd'hui un emploi dans votre entreprise. »
\end{exemplebox}

L'observation est immédiate : la version A ressemble à un inventaire télégraphique ; la version B raconte un \emph{parcours}, avec une progression logique. Les informations sont les mêmes, mais les \cle{connecteurs} (mots et groupes de mots de liaison) établissent des rapports : la chronologie (\emph{après}), le lieu (\emph{où}), l'addition (\emph{de plus}), la conséquence (\emph{c'est pourquoi}).

\begin{aretenir}[L'essentiel]
La \cle{cohérence textuelle} repose sur des liaisons claires entre les phrases et entre les propositions. Un texte bien lié est plus facile à lire, plus convaincant et plus professionnel. À l'examen comme dans la vie courante, la qualité des liaisons est un critère d'évaluation de l'expression écrite.
\end{aretenir}

\courssub{Les adverbes de liaison}

\begin{definition}[Adverbe de liaison]
Un \cle{adverbe de liaison} est un mot invariable qui relie deux phrases ou deux propositions en exprimant un rapport logique : chronologie, addition, opposition, cause ou conséquence. Exemples : \emph{d'abord, ensuite, enfin, cependant, néanmoins, ainsi, donc, par conséquent}.
\end{definition}

Chaque adverbe de liaison exprime un rapport précis. Les choisir au hasard détruit le sens ; il faut donc connaître leurs \emph{nuances}.

\textbf{1. La chronologie} (ordre des événements) : \emph{d'abord} annonce le premier point ; \emph{ensuite}, \emph{puis} enchaînent ; \emph{enfin} conclut la série. Dans un CV ou une lettre, ils organisent le parcours : « \emph{D'abord} formé au lycée technique de Douala, j'ai \emph{ensuite} travaillé deux ans comme aide-électricien ; j'ai \emph{enfin} créé mon propre atelier. »

\textbf{2. L'addition} : \emph{aussi} et \emph{également} signalent qu'une information s'ajoute à la précédente ; \emph{ainsi que} introduit un élément qui s'ajoute ou qui illustre (« titulaire du permis B, \emph{ainsi que} du permis C »).

\textbf{3. L'opposition} : \emph{cependant} et \emph{néanmoins} introduisent une restriction. Nuance : \emph{cependant} marque un contraste simple ; \emph{néanmoins} signifie « malgré cela » et suppose une concession. « Mon expérience est courte ; \emph{cependant}, elle est variée. » / « Je n'ai pas encore de diplôme universitaire ; je possède \emph{néanmoins} une solide pratique du métier. »

\textbf{4. La conséquence} : \emph{donc} est courant mais un peu familier dans un écrit formel ; \emph{par conséquent} et \emph{ainsi} (au sens de « c'est pourquoi ») conviennent mieux aux écrits utilitaires. « J'ai suivi une formation en comptabilité ; \emph{par conséquent}, je maîtrise les logiciels de gestion. »

\begin{attention}[Piège fréquent]
N'empilez pas plusieurs connecteurs de même valeur dans un même texte (« donc... donc... donc... »). Variez : \emph{donc, par conséquent, c'est pourquoi, ainsi}. Et n'utilisez jamais un connecteur d'opposition là où vous exprimez une conséquence : « J'ai le permis C, \emph{cependant} je peux conduire un camion » est fautif ; il faut « \emph{donc} je peux conduire un camion ».
\end{attention}

\courssub{Les locutions adverbiales de liaison}

\begin{definition}[Locution adverbiale de liaison]
Une \cle{locution adverbiale} est un groupe de mots invariable qui joue le même rôle qu'un adverbe de liaison : \emph{en effet, en outre, de plus, par ailleurs, en revanche, c'est pourquoi, grâce à, à la suite de}.
\end{definition}

Précisons les emplois des plus utiles :

\begin{itemize}
\item \emph{en effet} : introduit une \textbf{explication ou une justification} de ce qui précède. « Je maîtrise la soudure à l'arc ; \emph{en effet}, j'ai pratiqué ce procédé pendant tout mon stage. »
\item \emph{de plus} / \emph{en outre} / \emph{par ailleurs} : introduisent une \textbf{addition}. Nuance : \emph{par ailleurs} signale souvent un aspect différent, un changement d'angle. « Je parle couramment l'anglais ; \emph{par ailleurs}, je possède de bonnes notions d'informatique. »
\item \emph{en revanche} : introduit une \textbf{opposition} entre deux éléments contrastés. « Ma scolarité a été irrégulière ; \emph{en revanche}, mon expérience professionnelle est riche. »
\item \emph{c'est pourquoi} : introduit une \textbf{conséquence}. « Votre entreprise recherche un technicien polyvalent ; \emph{c'est pourquoi} je me permets de vous adresser ma candidature. »
\item \emph{grâce à} : exprime une \textbf{cause favorable}. « \emph{Grâce à} ma formation en froid et climatisation, je peux entretenir vos installations. »
\item \emph{à la suite de} : exprime une \textbf{cause ou un enchaînement} dans le temps. « \emph{À la suite de} mon stage, l'entreprise m'a proposé un contrat. »
\end{itemize}

\begin{methode}[Choisir le bon connecteur]
Pour choisir un connecteur, suivez ces étapes :
\begin{enumerate}
\item \textbf{Identifiez le rapport logique} entre les deux idées : s'agit-il d'ajouter, d'opposer, d'expliquer, de conclure, d'ordonner dans le temps ?
\item \textbf{Interrogez-vous} : la deuxième phrase ajoute-t-elle une information (addition) ? Contredit-elle la première (opposition) ? L'explique-t-elle (cause) ? En découle-t-elle (conséquence) ?
\item \textbf{Choisissez dans la bonne colonne} du tableau de classement ci-dessous.
\item \textbf{Vérifiez la ponctuation} : le connecteur est généralement suivi d'une virgule, et les deux phrases sont séparées par un point-virgule ou un point.
\item \textbf{Relisez à voix haute} : si le rapport logique n'est pas évident pour un lecteur, le connecteur est mal choisi.
\end{enumerate}
\end{methode}

\begin{popfigure}
\begin{center}
\begin{tabularx}{\linewidth}{|l|X|X|}
\hline
\rowcolor{popblueL} \textbf{Rapport logique} & \textbf{Adverbes de liaison} & \textbf{Locutions adverbiales} \\
\hline
Addition & aussi, également & de plus, en outre, par ailleurs, ainsi que \\
\hline
Opposition & cependant, néanmoins, pourtant & en revanche, par contre, malgré cela \\
\hline
Cause / explication & --- & en effet, grâce à, à cause de, à la suite de \\
\hline
Conséquence & donc, ainsi, par conséquent, alors & c'est pourquoi, de ce fait \\
\hline
Chronologie & d'abord, ensuite, puis, enfin & au début, dans un premier temps, pour finir \\
\hline
\end{tabularx}
\end{center}
\legende{Tableau de classement des principaux connecteurs selon le rapport logique exprimé. À consulter avant chaque rédaction d'écrit utilitaire.}
\end{popfigure}

\courssub{Les pronoms relatifs simples}

Jusqu'ici, les connecteurs relient des \emph{phrases}. Le \cle{pronom relatif}, lui, relie deux \emph{propositions à l'intérieur d'une même phrase}, en reprenant un nom déjà mentionné (son \cle{antécédent}). Il évite les répétitions et allège le texte — une qualité précieuse dans un CV où chaque ligne compte.

\begin{definition}[Pronom relatif]
Le \cle{pronom relatif} est un mot qui remplace un nom (l'\cle{antécédent}) et introduit une \cle{proposition subordonnée relative} qui complète ce nom. Les pronoms relatifs simples sont : \emph{qui, que, quoi, dont, où}.
\end{definition}

\begin{popfigure}
\begin{center}
\begin{tikzpicture}[
  box/.style={draw=popblue, thick, rounded corners, fill=popblueL, align=center, minimum height=0.9cm, inner sep=6pt},
  relbox/.style={draw=poporange, thick, rounded corners, fill=poporangeL, align=center, minimum height=0.9cm, inner sep=6pt},
  fleche/.style={-{Stealth[length=3mm]}, thick, popdark}
]
\node[box, align=center] (ant) at (0,0){\textbf{Antécédent}\\ « l'entreprise »};
\node[relbox, align=center] (pr) at (5.2,0){\textbf{Pronom relatif}\\ « \emph{qui} »};
\node[box, fill=poppurpleL, draw=poppurple, align=center] (prop) at (10.4,0){\textbf{Proposition relative}\\ « m'a recruté »};
\draw[fleche] (pr) to[bend right=25] node[midway, below, font=\small\itshape]{reprend et remplace} (ant);
\draw[fleche, poporange] (ant) -- (pr);
\draw[fleche, poppurple] (pr) -- (prop);
\node[align=center, font=\small, popmuted] at (5.2,-2.1) {Phrase complète : « J'ai travaillé dans une entreprise \emph{qui} m'a recrutée. »};
\end{tikzpicture}
\end{center}
\legende{Fonction du pronom relatif : il reprend l'antécédent (ici « l'entreprise ») et introduit la proposition subordonnée relative qui le précise.}
\end{popfigure}

\textbf{Emplois des pronoms relatifs simples :}

\begin{itemize}
\item \emph{\textbf{qui}} : \textbf{sujet} du verbe de la subordonnée. « J'ai suivi une formation \emph{qui} dure deux ans. » (\emph{qui} = la formation, sujet de « dure »).
\item \emph{\textbf{que}} : \textbf{complément d'objet direct} (COD). « Le diplôme \emph{que} j'ai obtenu est reconnu par l'État. » (\emph{que} = le diplôme, COD de « j'ai obtenu »). Attention : \emph{que} s'élide en \emph{qu'} devant voyelle.
\item \emph{\textbf{quoi}} : après une \textbf{préposition}, pour les choses, surtout dans des tournures figées. « Voici \emph{à quoi} je me suis consacré pendant mon stage. »
\item \emph{\textbf{dont}} : complément introduit par \textbf{de} (voir la propriété ci-dessous).
\item \emph{\textbf{où}} : complément de \textbf{lieu} ou de \textbf{temps}. « L'atelier \emph{où} j'ai travaillé » ; « l'année \emph{où} j'ai passé mon Bac ».
\end{itemize}

\begin{propriete}[Le pronom relatif « dont »]
Le pronom relatif \emph{\cle{dont}} remplace un complément introduit par la préposition \emph{de} : complément du nom (« le responsable \emph{de} ce service » $\rightarrow$ « le service \emph{dont} je suis responsable »), complément d'adjectif (« fier \emph{de} ce résultat » $\rightarrow$ « le résultat \emph{dont} je suis fier ») ou complément d'un verbe construit avec \emph{de} (« parler \emph{de} ce projet » $\rightarrow$ « le projet \emph{dont} j'ai parlé »). Cette construction est exigible dans les épreuves de langue du Baccalauréat technique.
\end{propriete}

\begin{exemplebox}[« dont » en action]
« J'ai dirigé une équipe. Je suis fier de cette équipe. » $\rightarrow$ « J'ai dirigé une équipe \emph{dont} je suis fier. »

« J'ai utilisé une machine. J'ai besoin de cette machine. » $\rightarrow$ « La machine \emph{dont} j'ai besoin est une fraiseuse numérique. »
\end{exemplebox}

\begin{attention}[Erreur fréquente : « où » et « ou »]
Ne confondez jamais \emph{\textbf{où}} (avec accent), pronom relatif ou adverbe interrogatif indiquant le lieu ou le temps, et \emph{\textbf{ou}} (sans accent), conjonction de coordination qui exprime l'alternative. Test infaillible : si l'on peut remplacer par « \emph{ou bien} », on écrit \emph{ou} sans accent. « L'entreprise \emph{où} j'ai fait mon stage » (on ne peut pas dire « ou bien ») ; « Vous pouvez postuler par courrier \emph{ou} par courriel » (= « ou bien »). Une faute sur ce point dans un CV est très mal perçue par un employeur.
\end{attention}

\courssub{Les pronoms relatifs composés}

Lorsque l'antécédent est introduit par une préposition autre que \emph{de} (\emph{à, sur, dans, pour, avec...}), on utilise les \cle{pronoms relatifs composés} : \emph{lequel, laquelle, lesquels, lesquelles}. Ils s'\textbf{accordent en genre et en nombre} avec l'antécédent.

\begin{center}
\begin{tabularx}{\linewidth}{|l|X|X|X|}
\hline
\rowcolor{popblueL} \textbf{Préposition} & \textbf{Masculin singulier} & \textbf{Féminin singulier} & \textbf{Pluriel} \\
\hline
à & auquel & à laquelle & auxquels / auxquelles \\
\hline
de & duquel & de laquelle & desquels / desquelles \\
\hline
sur, dans, pour, avec... & sur lequel, dans lequel... & sur laquelle, dans laquelle... & sur lesquels, dans lesquelles... \\
\hline
\end{tabularx}
\end{center}

\begin{exemplebox}[Pronoms relatifs composés dans un parcours professionnel]
« Le poste \emph{auquel} je postule exige de la rigueur. » (postuler \emph{à} un poste)

« Les entreprises \emph{pour lesquelles} j'ai travaillé peuvent fournir des références. » (travailler \emph{pour} des entreprises)

« Le logiciel \emph{avec lequel} je réalise mes devis est Excel. » (réaliser \emph{avec} un logiciel)
\end{exemplebox}

\begin{attention}[Accord et contraction]
Deux difficultés : (1) le pronom composé s'accorde avec l'antécédent : « la formation \emph{à laquelle} » et non « \emph{auquel} » ; (2) avec \emph{à} et \emph{de}, les formes contractées sont obligatoires : \emph{auquel, auxquels, duquel, desquels} — jamais « à lequel » ni « de lequel ». Rappel : pour un complément avec \emph{de}, on préfère le plus souvent \emph{dont} : « l'entreprise \emph{dont} je parle » est plus naturel que « l'entreprise \emph{de laquelle} je parle ».
\end{attention}

\courssub{Application au CV et aux écrits utilitaires}

Pourquoi cette leçon de grammaire est-elle placée dans une unité consacrée au CV ? Parce que les outils de liaison y remplissent trois fonctions concrètes :

\begin{enumerate}
\item \textbf{Éviter les répétitions.} Grâce aux pronoms relatifs, un même nom n'est jamais répété : « J'ai suivi un stage de trois mois à la SODECOTON, stage \emph{au cours duquel} j'ai découvert la maintenance industrielle » plutôt que « J'ai suivi un stage... Pendant ce stage, j'ai découvert... ».
\item \textbf{Relier les informations.} Les adverbes et locutions montrent la logique du parcours : \emph{d'abord... ensuite... enfin} pour la chronologie des diplômes ; \emph{de plus, en outre} pour cumuler les compétences ; \emph{c'est pourquoi} pour conclure une lettre de motivation.
\item \textbf{Fluidifier le parcours présenté.} Un texte bien lié donne l'image d'un candidat organisé et clair — qualité que tout employeur recherche.
\end{enumerate}

\begin{exoresolu}[Transformer des phrases juxtaposées en phrases liées]
\textbf{Énoncé.} Réécrivez le texte suivant en reliant les phrases à l'aide de connecteurs et de pronoms relatifs appropriés : « J'ai obtenu le BEPC en 2021. Je me suis inscrit au lycée technique de Bafoussam. J'ai préparé un Baccalauréat F3. J'ai effectué un stage à la CAMWATER. J'ai appris la lecture des plans de réseau. Je n'ai pas d'expérience longue. Je suis disponible immédiatement. Je sollicite le poste de technicien. »
\tcblower
\textbf{Solution détaillée.} On identifie d'abord les rapports logiques : chronologie (BEPC $\rightarrow$ lycée $\rightarrow$ Bac), lieu du stage, contenu du stage, opposition (peu d'expérience / disponibilité), conséquence finale.

« Après l'obtention du BEPC en 2021, je me suis inscrit au lycée technique de Bafoussam, \emph{où} j'ai préparé un Baccalauréat F3. \emph{Ensuite}, j'ai effectué un stage à la CAMWATER, stage \emph{au cours duquel} j'ai appris la lecture des plans de réseau. Je n'ai pas une longue expérience ; \emph{en revanche}, je suis disponible immédiatement. \emph{C'est pourquoi} je sollicite le poste de technicien dans votre service. »

Vérification : \emph{où} reprend « lycée » (lieu) ; \emph{au cours duquel} reprend « stage » (masculin singulier, préposition \emph{au cours de}) ; \emph{en revanche} marque bien l'opposition ; \emph{c'est pourquoi} introduit la conséquence. Le texte est fluide et sans répétition.
\end{exoresolu}

\begin{exercice}[Entraînement à la réécriture]
\textbf{Exercice 1.} Reliez chaque paire de phrases à l'aide du pronom relatif qui convient (\emph{qui, que, dont, où, auquel, pour laquelle}) :
\begin{enumerate}
\item J'ai obtenu un diplôme. Ce diplôme est très recherché par les entreprises.
\item L'entreprise m'a embauché. J'avais envoyé mon CV à cette entreprise.
\item Voici l'atelier. J'ai appris la soudure dans cet atelier.
\item J'ai suivi une formation. J'ai besoin de cette formation pour ce poste.
\item Le poste est vacant. Je postule à ce poste.
\item J'ai travaillé pour une société. Cette société est spécialisée en génie civil.
\end{enumerate}
\textbf{Exercice 2.} Complétez avec un connecteur adapté au rapport logique indiqué : « J'ai une formation en électrotechnique ; \dots (addition), je possède le permis B. Je n'ai jamais travaillé dans une grande entreprise ; \dots (opposition), mes stages m'ont donné une réelle pratique. Votre société recrute des techniciens débutants ; \dots (conséquence), je vous adresse ma candidature. »
\end{exercice}

\begin{corrige}[Exercice 1 et 2]
\textbf{Exercice 1.} 1. « J'ai obtenu un diplôme \emph{qui} est très recherché » (\emph{qui}, sujet). 2. « L'entreprise \emph{à laquelle} j'avais envoyé mon CV m'a embauché » (envoyer \emph{à}). 3. « Voici l'atelier \emph{où} j'ai appris la soudure » (lieu). 4. « La formation \emph{dont} j'ai besoin... » (avoir besoin \emph{de}). 5. « Le poste \emph{auquel} je postule est vacant » (postuler \emph{à}). 6. « La société \emph{pour laquelle} j'ai travaillé est spécialisée en génie civil » (travailler \emph{pour}).

\textbf{Exercice 2.} « J'ai une formation en électrotechnique ; \emph{de plus} (ou \emph{en outre}), je possède le permis B. Je n'ai jamais travaillé dans une grande entreprise ; \emph{cependant} (ou \emph{en revanche}), mes stages m'ont donné une réelle pratique. Votre société recrute des techniciens débutants ; \emph{c'est pourquoi} (ou \emph{par conséquent}), je vous adresse ma candidature. »
\end{corrige}

\begin{aretenir}[Bilan de la section]
Les outils de liaison garantissent la cohérence d'un écrit utilitaire. Les \cle{adverbes de liaison} (\emph{d'abord, ensuite, cependant, donc, par conséquent}) et les \cle{locutions adverbiales} (\emph{en effet, de plus, en revanche, c'est pourquoi, grâce à}) relient les phrases selon un rapport logique précis : chronologie, addition, opposition, cause ou conséquence. Les \cle{pronoms relatifs} simples (\emph{qui, que, quoi, dont, où}) et composés (\emph{lequel, auquel, duquel...}) relient les propositions en reprenant un antécédent, ce qui évite les répétitions. Maîtriser ces outils, c'est donner à son CV et à ses lettres la fluidité d'un parcours professionnel clair et convaincant.
\end{aretenir}

\coursec{Rédaction du CV : travail préparatoire}

Après avoir maîtrisé les \cle{outils de liaison} (adverbes, locutions adverbiales, pronoms relatifs) qui assurent la fluidité et la cohésion des phrases, nous abordons maintenant la phase concrète de \textbf{construction} du Curriculum Vitae. Posséder les briques linguistiques ne suffit pas ; il faut un plan d'architecte et une méthode rigoureuse pour assembler son parcours en un document persuasif, adapté à une offre précise. C'est l'objet de ce \emph{travail préparatoire} : passer de l'inventaire brut à la sélection stratégique.

\courssub{La démarche en six étapes : de l'analyse de l'offre au document final}

Rédiger un CV efficace ne s'improvise pas. C'est un \cle{processus itératif} qui commence bien avant la mise en page. Voici la méthode standardisée, applicable à toute situation d'embauche.

\begin{methode}[La démarche de rédaction du CV en 6 étapes]
\textbf{Étape 1 : Analyser l'offre d'emploi (ou la situation d'embauche).}
Décrypter l'annonce : identifier le \cle{poste}, le \cle{profil exigé} (diplômes, années d'expérience), les \cle{compétences techniques} (savoir-faire) et \cle{compétences transverses} (savoir-être), les \cle{mots-clés} récurrents. Noter les exigences « indispensables » vs « souhaitables ».

\textbf{Étape 2 : Faire son inventaire personnel (le « stock »).}
Lister \emph{exhaustivement} : diplômes/formation, expériences professionnelles (stages, jobs d'été, alternance), compétences techniques (logiciels, machines, langages), langues (niveau CECRL), outils bureautiques, permis, distinctions, centres d'intérêt pertinents. Ne rien filtrer pour l'instant.

\textbf{Étape 3 : Sélectionner et hiérarchiser l'information pertinente.}
Croiser l'inventaire (Étape 2) avec l'analyse de l'offre (Étape 1). Ne garder que ce qui \textbf{répond au besoin} de l'employeur. Hiérarchiser : le plus pertinent pour \emph{ce} poste en premier. Supprimer ou résumer le superflu (anciens jobs sans lien, détails excessifs sur des formations anciennes).

\textbf{Étape 4 : Choisir le type de CV adapté au profil.}
\begin{itemize}
    \item \textbf{Chronologique (antéchronologique) :} Parfait pour un parcours linéaire, sans trous, en progression dans un même secteur. Le plus classique.
    \item \textbf{Thématique (par compétences) :} Idéal pour une reconversion, des trous dans le parcours, ou une expérience très dispersée. Met l'accent sur \emph{ce qu'on sait faire} plutôt que \emph{quand/où}.
    \item \textbf{Mixte (hybride) :} Profil « Compétences » en haut + Expériences détaillées en bas. Le plus utilisé aujourd'hui car il rassure (dates) et vend (compétences).
\end{itemize}

\textbf{Étape 5 : Rédiger le brouillon rubrique par rubrique.}
Utiliser des \cle{formulations brèves} (style télégraphique, pas de phrases complètes), des \cle{verbes d'action} à l'infinitif ou au participe passé (gérer, concevoir, installer, réparer, organiser, former, développer, piloter...), des \cle{résultats chiffrés} si possible. Appliquer les outils de liaison vus en section 4 pour lier les items si phrases il y a (ex: \emph{« Responsable maintenance ; à ce titre, j'ai supervisé... »}).

\textbf{Étape 6 : Relire, corriger, mettre en page (mise en forme finale).}
Vérifier : \textbf{orthographe/grammaire} (zéro faute), \textbf{cohérence des dates} (pas de chevauchements inexpliqués), \textbf{outils de liaison} (cohérence interne), \textbf{format} (PDF, police lisible, 1 page idéalement, 2 max si expérience senior), \textbf{coordonnées} à jour. Faire relire par un tiers.
\end{methode}

\begin{popfigure}
\begin{tikzpicture}[
    node distance=1.2cm and 1.5cm,
    every node/.style={font=\small, align=center},
    stepbox/.style={rectangle, rounded corners, draw=popblue, thick, fill=popblueL, minimum width=3.5cm, minimum height=1.8cm},
    arrow/.style={-Stealth, thick, popdark},
    decision/.style={diamond, draw=poporange, thick, fill=poporangeL, aspect=2, minimum width=3cm}
]
% Nœuds
\node[stepbox, align=center] (e1){\textbf{Étape 1}\\\textbf{Analyser l'offre}\\(Poste, Profil, Mots-clés)};
\node[stepbox, below=of e1, align=center] (e2){\textbf{Étape 2}\\\textbf{Inventaire personnel}\\(Diplômes, Expériences, Compétences, Langues)};
\node[stepbox, below=of e2, align=center] (e3){\textbf{Étape 3}\\\textbf{Sélection \& Hiérarchisation}\\(Croiser Offre $\leftrightarrow$ Inventaire)};
\node[decision, below=of e3, align=center] (choix){Choix du type de CV\\Chrono / Thématique / Mixte};
\node[stepbox, below=of choix, align=center] (e5){\textbf{Étape 5}\\\textbf{Rédaction du brouillon}\\(Verbes d'action, Résultats, Liaisons)};
\node[stepbox, below=of e5, align=center] (e6){\textbf{Étape 6}\\\textbf{Relecture \& Finalisation}\\(Orthographe, Dates, Mise en page, PDF)};
% Flèches
\draw[arrow] (e1) -- (e2);
\draw[arrow] (e2) -- (e3);
\draw[arrow] (e3) -- (choix);
\draw[arrow] (choix) -- (e5);
\draw[arrow] (e5) -- (e6);
% Boucle de rétroaction (optionnel, style léger)
\draw[arrow, popmuted, dashed] (e6.east) -- ++(1.5,0) |- (e3.east) node[pos=0.25, right, font=\tiny, popmuted] {Si inadapté : recommencer};
\end{tikzpicture}
\legende{Organigramme : les 6 étapes de la démarche de rédaction d'un CV. Le processus est itératif : une relecture finale peut imposer de revenir à la sélection (Étape 3).}
\end{popfigure}

\courssub{Étape 1 : Analyser l'offre d'emploi — La clé de l'adaptation}

C'est l'étape la plus \textbf{négligée} et pourtant la plus \textbf{déterminante}. Un CV n'est pas une autobiographie ; c'est une \cle{réponse ciblée} à un besoin explicite.

\begin{exemplebox}[Analyse d'une offre : Technicien de Maintenance Industrielle]
\textbf{Extrait d'annonce :} « Entreprise agroalimentaire recherche Technicien Maintenance H/F. Bac Pro MSMA / BTS MS ou équivalent. 2 ans exp. exigés sur ligne de conditionnement. Maîtrise GMAO (DIMO Maint), lecture plans hydrauliques/pneumatiques. Autonome, réactif, esprit d'équipe. Permis B. Poste en 3x8. »

\textbf{Analyse structurée (grille de lecture) :}
\begin{center}
\begin{tabularx}{\linewidth}{|l|X|X|}\hline
\rowcolor{popblueL}
\textbf{Critère} & \textbf{Exigence de l'offre (Mots-clés)} & \textbf{Preuve à apporter dans le CV} \\ \hline
Diplôme & Bac Pro MSMA / BTS MS & \textbf{Formation} : Mettre en avant (titre exact, année, mention). \\ \hline
Expérience & 2 ans sur ligne conditionnement & \textbf{Expérience} : Détailler missions conditionnement (cadence, format). \\ \hline
Technique & GMAO DIMO Maint, Plans Hydra/Pneu & \textbf{Compétences Techniques} : Liste explicite \cle{DIMO Maint}, \cle{Schémas ISO}. \\ \hline
Savoir-être & Autonome, réactif, équipe & \textbf{Réalisations/Soft skills} : Exemple concret (ex: « Réduit temps d'arrêt 15\% »). \\ \hline
Contraintes & Permis B, 3x8 & \textbf{Infos pratiques / Disponibilité} : Préciser « Permis B + Véhicule », « Dispo 3x8 ». \\ \hline
\end{tabularx}
\end{center}
\par\medskip\textit{Tableau croisé : Offre d'emploi (exigences) $\leftrightarrow$ Éléments du CV à valoriser. Cette grille sert de feuille de route pour l'Étape 3.}
\end{exemplebox}

\begin{attention}[Erreur fatale : Le CV « passe-partout »]
\textbf{Envoyer le même CV à toutes les offres sans adaptation.}
Un recruteur passe en moyenne \textbf{6 à 30 secondes} sur un CV au premier tri. Si les mots-clés de \emph{son} annonce (ex: « GMAO DIMO Maint », « 3x8 », « Conditionnement ») n'apparaissent pas \textbf{visiblement} dans votre CV (rubrique Compétences, titres d'expériences), il classe le dossier sans lire la suite.
\par\medskip
\textbf{La parade :} Créez une \cle{version maître} (votre inventaire complet Étape 2) et \textbf{dupliquez-la} pour chaque candidature en ne gardant que les items validés par le tableau croisé (Étape 3). Changez l'ordre des rubriques, le titre du poste visé, les verbes d'action.
\end{attention}

\courssub{Étape 2 : L'inventaire personnel — Constituer sa « base de données »}

Avant de trier, il faut tout sortir. Cet inventaire est un document de travail \textbf{privé}, jamais envoyé tel quel. Il sert de réservoir pour toutes vos futures candidatures.

\begin{experience}[Atelier : Mon Inventaire Personnel (30 min)]
\textbf{Matériel :} Feuille A4 paysage ou tableur (colonnes : Rubrique, Détail, Dates, Résultats/Mots-clés).
\textbf{Protocole :}
\begin{enumerate}
    \item \textbf{Formation :} Du plus récent au plus ancien. Intitulé exact, établissement, année, mention, options, mémoire/projet de fin d'études.
    \item \textbf{Expériences pro :} Stages, alternance, jobs étudiants, bénévolat structurant. Pour \textbf{chaque ligne} : Structure, Dates (Mois/Année), Intitulé poste, \textbf{3 à 5 missions principales} (verbes d'action), \textbf{1 à 2 résultats chiffrés} (ex: « Gestion stock 500 refs », « Encadrement 3 stagiaires », « Maintenance 15 machines »).
    \item \textbf{Compétences Techniques (Hard Skills) :} Logiciels (CAO, ERP, GMAO, Bureautique), Machines/Outillages, Normes (ISO, HACCP), Langages (Python, Ladder), Permis/Engins (CACES).
    \item \textbf{Langues :} Niveau CECRL (A1 à C2) ou scores (TOEIC 780). Ne pas surestimer.
    \item \textbf{Soft Skills / Centres d'intérêt :} Sport collectif (esprit équipe), Musique (discipline), Association (responsabilité), Veille techno. Éviter : « Lecture, Cinéma, Voyage » (trop vague).
\end{enumerate}
\textbf{Observation :} Cet inventaire révèle souvent des compétences oubliées (gestion de planning en job d'été, gestion budget association) qui deviennent des atouts pour un poste d'assistant ou technicien encadrant.
\end{experience}

\courssub{Étape 3 : Sélection et hiérarchisation — Le tri sélectif}

C'est l'étape de la \cle{stratégie}. On ne met pas tout. On met ce qui \emph{vend} pour \emph{ce} poste.

\begin{propriete}[Règle de pertinence : Le ratio Signal/Bruit]
Toute information qui ne renforce pas votre candidature pour \emph{ce} poste spécifique est du \textbf{bruit} qui noie le \textbf{signal}.
\begin{itemize}
    \item \textbf{Jeune diplômé (peu d'exp) :} Formation en \textbf{1er} (détails : projets, stages, options). Expériences même courtes en 2ème. Compétences techniques issues de l'école.
    \item \textbf{Expérimenté :} Expériences en \textbf{1er} (les 3 dernières détaillées, les anciennes résumées). Formation condensée.
    \item \textbf{Reconversion / Trous :} CV \textbf{Thématique ou Mixte}. Rubrique « Compétences Clés » ou « Réalisations Majeures » en tout premier. Expériences réduites à lignes dates/titres.
\end{itemize}
\end{propriete}

\begin{exoresolu}[Hiérarchisation pour un poste de « Technicien Bureau d'Études Électrique »]
\textbf{Profil candidat :} BTS Électrotechnique (2022). 1 an stage BE (See Electrical). 2 ans Opérateur câblage (usine). 1 an Job été Vente informatique. Anglais B1. Permis B.
\textbf{Offre :} Schémas See Electrical, Nomenclatures, Notes de calcul, Normes NF C 15-100. Anglais technique lu. Déplacements occasionnels.

\textbf{Choix de structure (Mixte) et Ordre des rubriques :}
\begin{enumerate}
    \item \textbf{Titre :} Technicien Bureau d'Études Électrique (Junior)
    \item \textbf{Compétences Clés (Thématique) :} \cle{See Electrical (Expert)}, \cle{Nomenclatures / BOM}, \cle{NF C 15-100}, \cle{Notes de calcul sectionnement}, \cle{Anglais technique (Lu/Écrit)}.
    \item \textbf{Expériences Professionnelles (Antéchrono) :}
        \begin{itemize}
            \item 2023-2024 : \textbf{Stagiaire puis Technicien BE} (Entreprise X) — \emph{Détail 4 missions See Electrical + 1 résultat : « 50 schémas/mois, 0 retour atelier »}.
            \item 2022-2023 : \textbf{Opérateur Câblage} (Usine Y) — \emph{1 ligne : « Lecture schémas, câblage armoires, respect normes. Compréhension terrain du BE »}. (Valorise la double compétence).
        \end{itemize}
    \item \textbf{Formation :} BTS Électrotechnique (2022) — Projet fin d'études : « Automatisme ligne conditionnement » (pertinent).
    \item \textbf{Job Été / Divers :} Vente Informatique (2021) — \emph{1 ligne : « Relation client, conseil technique »}. (Soft skill utile).
    \item \textbf{Langues / Permis / Centres d'intérêt :} Anglais B1, Permis B, Vélo (endurance).
\end{enumerate}
\textbf{Justification :} L'expérience « Vente » est reléguée en bas, l'expérience « Câblage » est reformulée pour montrer qu'elle nourrit la compétence BE. La rubrique Compétences techniques est en tête pour matcher les mots-clés de l'annonce immédiatement.
\tcblower
\textbf{Solution commentée :} Le candidat a appliqué la grille de l'Étape 1. Il a masqué l'aspect « junior » en mettant la compétence technique (See Electrical) avant la chronologie. Il a transformé l'expérience « câblage » (a priori hors sujet BE) en atout « connaissance terrain ».
\end{exoresolu}

\courssub{Étape 4 : Choisir le type de CV — Adapter le contenant au contenu}

Le choix du format n'est pas esthétique, il est \textbf{rhétorique}.

\begin{center}
\begin{tabularx}{\linewidth}{|l|X|X|X|}\hline
\rowcolor{popblueL}
\textbf{Type de CV} & \textbf{Profil Cible} & \textbf{Avantage Principal} & \textbf{Risque / Vigilance} \\ \hline
\textbf{Chronologique (Antéchrono)} & Parcours linéaire, montée en charge, même secteur. & Rassurant : lisibilité temporelle, fidélité. & Met en lumière les trous, la stagnation, la dispersion. \\ \hline
\textbf{Thématique (Fonctionnel)} & Reconversion, trous longs, multi-activités, freelance. & Masque la chronologie, met en avant le « Savoir-faire ». & Déstabilise le recruteur (où/quand ?). Exige une rubrique « Parcours » claire en bas. \\ \hline
\textbf{Mixte (Hybride)} & \textbf{Majorité des cas (Junior à Senior).} & \textbf{Le meilleur des deux :} Accroche par compétences + Preuve par l'expérience datée. & Plus long à construire. Risque de redondance (compétences listées 2 fois). \\ \hline
\end{tabularx}
\end{center}

\courssub{Étape 5 : Rédaction du brouillon — Le style « CV »}

Le style du CV est \textbf{nominal, télégraphique, orienté résultat}. Pas de « Je », pas de phrases complètes (sauf accroche/profil).

\begin{exemplebox}[Verbes d'action à privilégier — Banque de formulation]
Remplacez les verbes faibles (fait, aidé, participé, géré) par des verbes \textbf{précis, mesurables, techniques} :
\begin{center}
\begin{tabularx}{\linewidth}{|l|X|}\hline
\rowcolor{popblueL}
\textbf{Domaine / Action} & \textbf{Verbes d'action forts (Infinitif)} \\ \hline
\textbf{Gestion / Pilotage} & \cle{Piloter}, \cle{Superviser}, \cle{Coordonner}, \cle{Planifier}, \cle{Budgétiser}, \cle{Optimiser}, \cle{Rationaliser} \\ \hline
\textbf{Technique / Production} & \cle{Concevoir}, \cle{Développer}, \cle{Programmer}, \cle{Usiner}, \cle{Câbler}, \cle{Installer}, \cle{Mettre en service}, \cle{Qualifier}, \cle{Valider} \\ \hline
\textbf{Maintenance / Dépannage} & \cle{Diagnostiquer}, \cle{Réparer}, \cle{Prévenir}, \cle{Fiabiliser}, \cle{Modifier}, \cle{Retrofiter}, \cle{Documenter} \\ \hline
\textbf{Étude / Analyse} & \cle{Analyser}, \cle{Modéliser}, \cle{Simuler}, \cle{Dimensionner}, \cle{Chiffrer}, \cle{Rédiger (cahier des charges)}, \cle{Veiller} \\ \hline
\textbf{Relationnel / Commercial} & \cle{Négocier}, \cle{Fidéliser}, \cle{Prospecter}, \cle{Conseiller}, \cle{Former}, \cle{Animer}, \cle{Médier} \\ \hline
\textbf{Organisation / Qualité} & \cle{Structurer}, \cle{Standardiser}, \cle{Normaliser}, \cle{Auditer}, \cle{Certifier}, \cle{Améliorer (Démarche 5S, Kaizen)} \\ \hline
\end{tabularx}
\end{center}
\textbf{Astuce :} Associez \textbf{Verbe + Objet + Contexte + Résultat Chiffré}.
\emph{Exemple :} « \textbf{Optimisé} la maintenance préventive (objet) \textbf{sur presse 200T} (contexte) $\rightarrow$ \textbf{Réduction 20\% temps d'arrêt} (résultat). »
\end{exemplebox}

\textbf{Application des outils de liaison (Rappel Section 4).}\par
Même en style télégraphique, la cohérence interne exige des liaisons logiques.
\begin{itemize}
    \item \textbf{Adverbes de liaison :} \cle{Par ailleurs}, \cle{En outre}, \cle{En conséquence}, \cle{Aussi} (pour lister des réalisations dans une même expérience).
    \item \textbf{Locutions adverbiales :} \cle{Dans ce cadre}, \cle{À ce titre}, \cle{En lien avec}, \cle{En vue de} (pour contextualiser une mission).
    \item \textbf{Pronoms relatifs (si phrases complètes en accroche) :} \cle{Dont}, \cle{Où}, \cle{Lequel} (ex: « Projet de fin d'études \textbf{dont} le sujet portait sur l'optimisation énergétique... »).
\end{itemize}

\courssub{Étape 6 : Relecture, correction et finalisation — La qualité professionnelle}

Un CV avec une faute d'orthographe, une date incohérente ou une police illisible est \textbf{éliminatoire} dans 80\% des cas pour un poste technique (rigueur attendue).

\begin{methode}[Check-list de relecture finale (Ordre d'exécution)]
\begin{enumerate}
    \item \textbf{Orthographe / Grammaire :} Correcteur automatique \textbf{+} relecture à voix haute \textbf{+} relecture par un tiers (professeur, pair, pro). Attention aux accords (ex: \emph{« stages effectués »}, \emph{« missions réalisées »}).
    \item \textbf{Cohérence temporelle :} Vérifier les chevauchements (Stage + Job été = OK si dates précises mois/année). Expliquer les trous $> 3$ mois (formation, recherche active, projet perso, congé parental) \emph{dans la lettre de motivation}, pas dans le CV.
    \item \textbf{Outils de liaison :} Vérifier la cohérence logique entre les puces d'une même expérience (\emph{« Diagnostiquer pannes ; \textbf{par conséquent} proposer solutions »}).
    \item \textbf{Mise en page :} Police unique (Calibri, Arial, Roboto, Open Sans), taille 10-11. Interlignes 1,15. Marges 1,5-2 cm. \textbf{Alignements parfaits} (tabulations, pas d'espaces). \textbf{En-tête :} Nom, Téléphone, Email pro, LinkedIn (URL courte), Ville, Permis. \textbf{Pas de photo} obligatoire (sauf pays/métier l'exigeant), si photo : pro, fond neutre, sourire.
    \item \textbf{Format de sortie :} \textbf{PDF uniquement} (nom : \texttt{Nom\_Prenom\_PosteVisé.pdf}). Tester l'ouverture sur mobile et PC. Fichier $< 500$ Ko.
    \item \textbf{RGPD / Données sensibles :} Pas de date de naissance, nationalité, situation familiale, photo d'identité (sauf demande), numéro de sécurité sociale.
\end{enumerate}
\end{methode}

\courssub{Adapter le CV à une situation d'embauche : L'exercice du « Sur-mesure »}

Au Probatoire / Baccalauréat technique, l'épreuve peut fournir une \textbf{annonce fictive} et demander de produire le CV correspondant. La réussite repose sur la \cle{transposition explicite} des exigences de l'annonce dans les rubriques du CV.

\begin{exoresolu}[Exercice d'intégration type Bac : Rédaction CV adapté]
\textbf{Sujet :} « Vous êtes titulaire d'un Baccalauréat Technicien en Électrotechnique (2024). Vous avez effectué un stage de 8 semaines chez \texttt{ENERCAM} (maintenance préventive/curative postes HT/BT) et un job d'été 2023 chez \texttt{ELECTRO-DISTRI} (préparation commandes, conseil client). Vous maîtrisez le logiciel \texttt{See Electrical}, la lecture de plans unifilaires, les normes NF C 15-100. Anglais niveau A2 (lu). Permis B. Vous postulez à l'annonce ci-dessous. \textbf{Rédigez votre CV.} »

\textbf{Annonce :} « \texttt{SOCAELEC} recrute \textbf{Technicien Maintenance Itinérant H/F}. Bac Pro/Bac Techno Électrotechnique. Déplacements régionaux (Permis B obligatoire). Maintenance préventive/curative équipements industriels. Lecture plans, GMAO appréciée. Autonomie, rigueur, relationnel client. Débutant accepté. »

\textbf{Production du CV (Structure Mixte recommandée) :}

\begin{center}
\begin{tabularx}{\linewidth}{|X|}\hline
\textbf{NGONO Jean-Pierre} \\
\textbf{Technicien de Maintenance Électrotechnique (Junior)} \\
Douala, Bonapriso | +237 6XX XX XX XX | jean-pierre.ngono@email.com | LinkedIn: linkedin.com/in/jp-ngono \\
Permis B + Véhicule personnel | Disponibilité immédiate | Mobilité régionale \\ \hline
\end{tabularx}
\end{center}

\vspace{0.3cm}
\begin{center}
\begin{tabularx}{\linewidth}{|l|X|}\hline
\textbf{COMPÉTENCES TECHNIQUES MAJEURES} & \textbf{Électrotechnique :} Maintenance préventive/curative HT/BT, Lecture plans unifilaires/multifilaires, Normes NF C 15-100, Schémas \cle{See Electrical (Notions)}, Mesures/Tests (Mégohmmètre, Contrôleur d'isolement). \\
& \textbf{Outils / Méthodes :} GMAO (notions), Démarche 5S, Rédaction rapports d'intervention, Consignation/LOTO. \\
& \textbf{Bureautique :} Pack Office (Excel : TCD, RechercheV), Outlook. \\ \hline
\textbf{EXPÉRIENCES PROFESSIONNELLES} & \textbf{03/2024 -- 05/2024 (8 sem.) : Stagiaire Technicien Maintenance -- ENERCAM, Douala} \\
& \begin{itemize}[leftmargin=*, label=--]
    \item \textbf{Réalisé} la maintenance préventive sur postes de transformation 30/0,4 kV (contrôle isolement, serrage, thermographie) \textbf{selon gammes 5S}.
    \item \textbf{Assisté} les techniciens itinérants en dépannage curatif BT (recherche défauts terre, court-circuit) \textbf{en autonomie partielle}.
    \item \textbf{Exploité} la GMAO pour clôturer les OT (Ordres de Travail) et remonter les anomalies \textbf{(+15\% traçabilité)}.
    \item \textbf{Lu} et \textbf{interprété} plans unifilaires et schémas de câblage pour localiser organes de coupure.
\end{itemize} \\
& \textbf{07/2023 -- 08/2023 (2 mois) : Équipier Magasin -- ELECTRO-DISTRI, Douala} \\
& \begin{itemize}[leftmargin=*, label=--]
    \item \textbf{Préparé} commandes clients (matériel électrique) : lecture bons de commande, picking, conditionnement.
    \item \textbf{Conseillé} clients professionnels sur choix disjoncteurs, câbles, coffrets \textbf{(connaissance nomenclature)}.
    \item \textbf{Géré} réceptions fournisseurs : contrôle quantitatif/qualitatif, étiquetage, rangement stock.
\end{itemize} \\ \hline
\textbf{FORMATION} & \textbf{2024 : Baccalauréat Technicien Électrotechnique} -- Lycée Technique de Douala (Mention Assez Bien). \\
& \emph{Projet de fin d'études :} « Automatisation poste de relevage eaux usées » (Automate Siemens S7-1200, TIA Portal, See Electrical). \\ \hline
\textbf{LANGUES / DIVERS} & \textbf{Français :} Langue maternelle. \textbf{Anglais :} Niveau A2 (Lu technique / Écrit basique) -- \emph{En cours de perfectionnement}. \\
& \textbf{Centres d'intérêt :} Football (collectif, endurance), Bricolage électronique (réparation petits appareils), Bénévolat Croix-Rouge (secourisme). \\ \hline
\end{tabularx}
\end{center}

\tcblower
\textbf{Analyse de l'adaptation (Justification de la démarche) :}
\begin{enumerate}
    \item \textbf{Titre ciblé :} Reprend l'intitulé exact de l'annonce + « Junior » (honnêteté).
    \item \textbf{En-tête :} « Permis B + Véhicule » et « Mobilité régionale » mis en \textbf{très haut} car critère éliminatoire explicite (« Déplacements régionaux », « Permis B obligatoire »).
    \item \textbf{Rubrique Compétences en 1er (Mixte) :} Le candidat est débutant. La formation est récente mais l'expérience (stage) est le point fort. La rubrique Compétences valide \textbf{immédiatement} les mots-clés : \cle{Maintenance HT/BT}, \cle{Plans unifilaires}, \cle{NF C 15-100}, \cle{See Electrical}, \cle{GMAO}, \cle{5S}.
    \item \textbf{Stage détaillé (8 semaines = expérience significative) :} Verbes d'action (\cle{Réalisé}, \cle{Assisté}, \cle{Exploité}, \cle{Lu/Interprété}). Résultat chiffré (\cle{+15\% traçabilité}). Vocabulaire métier (\cle{OT}, \cle{Consignation}, \cle{Thermographie}).
    \item \textbf{Job d'été valorisé :} Pas « vendeur » mais « Équipier Magasin matériel électrique ». Missions tournées vers : \cle{Connaissance nomenclature}, \cle{Conseil technique}, \cle{Réception/Stock} (rigueur, relation client demandés dans l'annonce).
    \item \textbf{Projet fin d'études :} Montre la capacité projet + See Electrical + Automatisme (polyvalence).
    \item \textbf{Anglais A2 assumé :} « Lu technique » (utile pour notices), « En cours de perfectionnement » (volonté).
\end{enumerate}
\end{exoresolu}

\courssub{Justifier sa démarche : L'oralisation de la stratégie}

Au-delà de l'écrit, le programme exige de \textbf{justifier ses choix} (oralement ou par écrit). C'est la preuve de la \cle{métacognition} : vous savez \emph{pourquoi} vous avez fait ce CV ainsi.

\begin{aretenir}[Grille d'auto-évaluation / Justification type]
Pour chaque choix fort, soyez prêt à répondre : \textbf{« Pourquoi ceci ? Pourquoi pas cela ? »}
\begin{itemize}
    \item \textbf{Ordre des rubriques :} « J'ai mis Compétences avant Expériences car l'annonce listait 6 compétences techniques précises ; je voulais qu'elles soient vues en 5 secondes. »
    \item \textbf{Choix du type Mixte :} « Mon stage est mon atout principal, mais je n'ai qu'une ligne d'expérience. Le Mixte me permet de valider les mots-clés (Compétences) puis de prouver par le stage (Expériences). »
    \item \textbf{Formulation « Réalisé maintenance préventive... selon gammes 5S » :} « L'annonce demande "rigueur". "Selon gammes 5S" prouve la rigueur méthodologique, pas seulement l'action. »
    \item \textbf{Valorisation job d'été :} « J'ai gardé ce job car l'annonce demande "relationnel client". Conseiller des pros sur du matériel électrique prouve ce relationnel technique. »
    \item \textbf{Mention « Permis B + Véhicule » en en-tête :} « Critère obligatoire + Itinérance. Si je le mets en bas, le recruteur ne le voit pas au premier tri. »
\end{itemize}
\end{aretenir}

\begin{savaistu}[Le CV au Baccalauréat Technique : Enjeu d'examen]
À l'épreuve de Français du \textbf{Baccalauréat Technique (Probatoire / Bac)}, la production d'écrits utilitaires est une composante notée. Le sujet peut prendre la forme :
\begin{itemize}
    \item \textbf{Production guidée :} « À partir des éléments biographiques fournis et de l'annonce ci-jointe, rédigez le CV du candidat. » (Noté sur : structure, adaptation, vocabulaire technique, correction, mise en page).
    \item \textbf{Production libre / Intégration :} « Vous souhaitez postuler à... Rédigez votre CV et justifiez vos choix de rédaction. » (Noté sur : pertinence, cohérence, argumentation méta-langagière).
    \item \textbf{Contraction/Résumé :} « Résumez le parcours suivant en une rubrique "Expériences" de CV (80 mots max). » (Noté sur : sélection, concision, verbes d'action).
\end{itemize}
\textbf{Conseil stratégique :} Entraînez-vous à rédiger des CV \textbf{sur feuille blanche} en 30-40 minutes chrono (analyse annonce + brouillon + copie propre). La gestion du temps est cruciale.
\end{savaistu}

\begin{exercice}[Entraînement : De l'annonce au CV adapté]
\textbf{Consigne :} À partir de l'annonce ci-dessous et du profil « Élève ingénieur / BTS option Froid/Clim (2024) », effectuez les \textbf{Étapes 1, 3 et 5} (Analyse, Sélection, Rédaction brouillon rubrique Compétences + 1 Expérience).
\begin{quote}
\textbf{Annonce :} « \texttt{FRIGO-SERVICES} (Yaoundé) recrute \textbf{Technicien Frigoriste Itinérant}. Bac Pro/BTS Froid/Clim. Expérience stage/alternance sur installations commerciales (centrales, vitrines, chambres froides). Fluides : R404A, R134a, CO2 (transcritique). Lecture plans frigorifiques, schémas électriques, régulation (Danfoss, Carel). Diagnostic pannes, recharge, étanchéité. Permis B. Déplacements régionaux. Rigueur, autonomie, relation client. »
\end{quote}
\textbf{Profil candidat (Inventaire) :}
\begin{itemize}
    \item BTS Froid/Clim 2024 (Mention Bien). Projet : « Centrales CO2 transcritique supermarché ».
    \item Stage 12 sem. (2023) : \texttt{COLD-CAM} (Douala). Maintenance préventive/curative supermarchés (Centrales R404A, Vitrines R134a). GMAO \texttt{Maintair}. Détection fuites (gaz traceur). Relation client (directeur magasin).
    \item Job été 2022 : \texttt{CLIM-EXPRESS} (Yaoundé). Pose/SAV clim split/console (R32). Électricité (tableau, raccordement).
    \item Compétences : Plans (AutoCAD, See Electrical), Régulation (Danfoss AK, Carel), Fluides (R404A, R134a, R32, Notions CO2), Électrotechnique, Permis B, Anglais B1.
\end{itemize}
\textbf{Travail demandé :}
\begin{enumerate}
    \item Remplir le tableau croisé \textbf{Offre $\leftrightarrow$ CV} (comme Exemple p. 2).
    \item Choisir le type de CV (Chrono / Thématique / Mixte) et justifier en 2 lignes.
    \item Rédiger la rubrique \textbf{COMPÉTENCES TECHNIQUES} (style télégraphique, mots-clés annonce) et l'\textbf{EXPÉRIENCE STAGE} (4 puces : Verbe + Objet + Contexte + Résultat).
\end{enumerate}
\end{exercice}

\begin{corrige}[Exercice n°1 : Entraînement adaptation CV]
\textbf{1. Tableau croisé (extrait) :}
\begin{center}
\begin{tabularx}{\linewidth}{|l|X|X|}\hline
\rowcolor{popblueL}
\textbf{Exigence Annonce} & \textbf{Mot-clé à matcher} & \textbf{Preuve Profil (Inventaire)} \\ \hline
Fluides R404A, R134a, CO2 & \cle{R404A, R134a, CO2 (transcritique)} & Stage (R404A, R134a) + Projet BTS (CO2 transcritique) + Job (R32) \\ \hline
Régulation Danfoss, Carel & \cle{Danfoss AK, Carel} & Compétences techniques listées \\ \hline
Diagnostic, étanchéité, recharge & \cle{Diagnostic pannes, Étanchéité (gaz traceur), Recharge} & Stage : Détection fuites, maintenance curative \\ \hline
Permis B, Itinérance, Client & \cle{Permis B + Véhicule}, \cle{Relation client} & Permis B + Stage (relation directeur magasin) \\ \hline
\end{tabularx}
\end{center}

\textbf{2. Choix du type : \textbf{Mixte}.}
\textit{Justification :} Profil junior avec un stage long (12 sem.) très pertinent (supermarchés, fluides annonce) + Projet fin d'études fort (CO2). Le Mixte permet de valider la liste technique exhaustive (Compétences) puis de prouver l'opérationnalité terrain (Expérience Stage détaillée).

\textbf{3. Rédaction brouillon :}

\textbf{COMPÉTENCES TECHNIQUES MAJEURES}
\begin{itemize}
    \item \textbf{Fluides frigorigènes :} Manipulation/Recharge \cle{R404A, R134a, R32} ; Notions \cle{CO2 transcritique} (Projet BTS) ; Certif. aptitude manipulation fluides (Cat I) en cours.
    \item \textbf{Installations commerciales :} Centrales frigorifiques (basse/moyenne temp.), Vitrines, Chambres froides, Groupes logés.
    \item \textbf{Régulation / Automatisme :} Paramétrage/Diagnostic \cle{Danfoss AK (AK-CC, AK-PC)}, \cle{Carel (pCO, uPC)} ; Lecture schémas électriques / plans frigorifiques (AutoCAD, See Electrical).
    \item \textbf{Maintenance / Diagnostic :} Maintenance préventive (gammes, GMAO \cle{Maintair}) ; Recherche fuites \cle{gaz traceur / détecteur électronique} ; Diagnostic pannes (électrique, frigorifique, régulation) ; Étanchéité / Mise sous vide / Recharge.
    \item \textbf{Électrotechnique / Sécurité :} Raccordement armoires, lecture schémas puissance/commande, Consignation LOTO, Habilitation électrique (BR/BC visée).
\end{itemize}

\textbf{EXPÉRIENCES PROFESSIONNELLES}
\textbf{02/2023 -- 07/2023 (12 sem.) : Stagiaire Technicien Frigoriste Itinérant -- COLD-CAM, Douala}
\begin{itemize}
    \item \textbf{Assuré} la maintenance préventive hebdomadaire de \textbf{12 sites supermarchés} (centrales R404A, vitrines R134a) \textbf{sous GMAO Maintair} $\rightarrow$ \textbf{Taux de respect planning 98\%}.
    \item \textbf{Diagnostiqué et réparé} pannes critiques (compresseurs, vannes électroniques, régulation Danfoss/Carel) \textbf{en autonomie} $\rightarrow$ \textbf{Réduction temps d'arrêt vitrines 30\%}.
    \item \textbf{Réalisé} recherches de fuites \textbf{au gaz traceur (H2/N2)} et détecteur électronique ; \textbf{effectué} mises sous vide, recharges fluides \textbf{selon procédures sécurité/environnement}.
    \item \textbf{Assuré l'interface client} (Directeurs magasins) : compte-rendu interventions, conseil préventif, gestion réclamations $\rightarrow$ \textbf{Satisfaction client 4,8/5}.
\end{itemize}
\end{corrige}

\coursec{Activités d'intégration, compte rendu et remédiation}

Dans la section précédente, tu as préparé ta rédaction : analyse de l'offre d'emploi, fiche d'identité professionnelle, tri des informations, plan des rubriques. Le moment est venu de passer à la \cle{production} complète, de soumettre ton texte au regard des autres, puis de corriger ce qui doit l'être. C'est toute la logique des \cle{activités d'intégration} : réinvestir dans une tâche réelle et complète tout ce que l'unité t'a appris (structure externe, structure interne, outils de liaison), puis vérifier, grâce au \cle{compte rendu} et à la \cle{remédiation}, que les acquis sont solides.

\courssub{Production individuelle : rédiger son propre CV}

\courssub{Consigne de production}

À partir de l'offre d'emploi ci-dessous (contexte camerounais), rédige ton CV complet en une page maximum.

\begin{exemplebox}[Offre d'emploi (fictive) — contexte camerounais]
\textbf{SOCIÉTÉ CAMTEL-SERVICES, Douala — recrute un(e) technicien(ne) de maintenance réseaux.}\par
Profil exigé : titulaire d'un Baccalauréat technique (F3, G ou équivalent) ou d'un BTS en télécommunications ; maîtrise de l'outil informatique ; sens de l'organisation ; permis B apprécié ; bonne expression écrite et orale en français, l'anglais est un atout.\par
Dossier à déposer : CV + lettre de motivation, avant le 30 novembre, au siège de l'entreprise, quartier Bonanjo.
\end{exemplebox}

\begin{methode}[Étapes de la production individuelle]
\begin{enumerate}
\item \textbf{Relire l'offre} et souligner les exigences (diplôme, compétences, qualités) : ton CV doit y répondre point par point.
\item \textbf{Reprendre ta fiche préparatoire} (section 5) : état civil, formation, expériences, compétences, centres d'intérêt.
\item \textbf{Rédiger rubrique par rubrique}, de haut en bas : en-tête, formation (de la plus récente à la plus ancienne), expérience professionnelle, compétences, divers.
\item \textbf{Soigner les liaisons} : utilise des pronoms relatifs (\emph{qui, que, où, dont}) et des adverbes de liaison (\emph{d'abord, ensuite, enfin, de plus}) pour éviter les phrases télégraphiques maladroites.
\item \textbf{Mettre en page} : une page, marges régulières, titres de rubriques visibles, alignements propres.
\item \textbf{Relire à voix haute} au moins deux fois, puis faire relire par un camarade.
\end{enumerate}
\end{methode}

\begin{attention}[L'erreur qui coûte cher : négliger la relecture finale]
Une seule faute d'orthographe dans un CV peut \textbf{éliminer une candidature}, car le recruteur y voit un manque de rigueur. Les fautes les plus fréquentes chez les candidats : \emph{« j'ai travailler »} (au lieu de \emph{travaillé}), \emph{« expérience professionel »} (deux \emph{l}), \emph{« maîtrise de l'informatique »} écrit \emph{« maitrise »} sans accent circonflexe, oublis d'accords (\emph{« les stages effectuée »}). Relis toujours ton CV \textbf{après une pause}, et fais-le relire par quelqu'un d'autre : l'œil de l'auteur ne voit plus ses propres erreurs.
\end{attention}

\courssub{Évaluation croisée : la grille d'analyse}

Une fois rédigés, les CV s'échangent au sein de la classe. Chaque élève annote le CV d'un camarade à l'aide de la grille construite collectivement en classe. Cette \cle{évaluation croisée} a deux vertus : elle oblige à lire un CV avec l'œil d'un recruteur, et elle fait découvrir ses propres défauts en les repérant chez autrui.

\begin{methode}[Grille d'auto-évaluation et d'évaluation croisée d'un CV]
Pour chaque critère, coche la case correspondante : \textbf{O} = oui, \textbf{P} = partiellement, \textbf{N} = non. Un CV est réussi s'il obtient au moins 8 « O » sur 10.
\begin{enumerate}
\item L'en-tête est complet (nom, prénom, adresse, téléphone, courriel).
\item La photo, si elle existe, est correcte et le CV tient sur une page.
\item Les rubriques sont clairement séparées et titrées.
\item La formation est présentée de la plus récente à la plus ancienne.
\item Les expériences sont datées et décrites par des verbes d'action.
\item Les compétences répondent aux exigences de l'offre.
\item Les pronoms relatifs et adverbes de liaison sont correctement employés.
\item L'orthographe et la grammaire sont irréprochables.
\item La mise en page est aérée, alignée, lisible.
\item Le CV est honnête : aucune information inventée ou exagérée.
\end{enumerate}
\end{methode}

\begin{popfigure}
\begin{center}
\rowcolors{2}{popblueL}{white}
\begin{tabularx}{\linewidth}{|c|X|c|c|c|}
\hline
\rowcolor{popblue}\textcolor{white}{\textbf{N\textdegree}} & \textcolor{white}{\textbf{Critère d'évaluation}} & \textcolor{white}{\textbf{O}} & \textcolor{white}{\textbf{P}} & \textcolor{white}{\textbf{N}} \\
\hline
1 & En-tête complet et exact & $\square$ & $\square$ & $\square$ \\
\hline
2 & Une page, photo correcte & $\square$ & $\square$ & $\square$ \\
\hline
3 & Rubriques titrées et ordonnées & $\square$ & $\square$ & $\square$ \\
\hline
4 & Formation en ordre antichronologique & $\square$ & $\square$ & $\square$ \\
\hline
5 & Expériences datées, verbes d'action & $\square$ & $\square$ & $\square$ \\
\hline
6 & Compétences adaptées à l'offre & $\square$ & $\square$ & $\square$ \\
\hline
7 & Liaisons correctes (relatifs, adverbes) & $\square$ & $\square$ & $\square$ \\
\hline
8 & Orthographe et grammaire soignées & $\square$ & $\square$ & $\square$ \\
\hline
9 & Mise en page aérée et lisible & $\square$ & $\square$ & $\square$ \\
\hline
10 & Honnêteté des informations & $\square$ & $\square$ & $\square$ \\
\hline
\end{tabularx}
\end{center}
\legende{Grille d'évaluation d'un CV en 10 critères : O = oui, P = partiellement, N = non. À utiliser pour l'auto-évaluation et l'évaluation croisée.}
\end{popfigure}

\courssub{Compte rendu collectif}

Après l'évaluation croisée, le professeur choisit trois ou quatre productions représentatives (un CV très réussi, un CV moyen, un CV faible) qui sont lues ou projetées devant la classe. Chaque auteur \textbf{justifie sa démarche} : pourquoi cet ordre de rubriques ? pourquoi avoir mis en avant tel stage ? pourquoi tel adjectif ? La classe commente ensuite à l'aide de la grille.

\begin{aretenir}[Ce que le compte rendu doit faire apparaître]
\begin{itemize}
\item Un bon CV \textbf{répond à l'offre} : il ne raconte pas toute ta vie, il sélectionne ce qui intéresse le recruteur.
\item La \textbf{forme compte autant que le fond} : un contenu excellent mal présenté passe inaperçu.
\item \textbf{Justifier ses choix} fait partie du savoir-faire évalué : au Baccalauréat technique, on te demandera souvent d'expliquer ta démarche, pas seulement de produire.
\end{itemize}
\end{aretenir}

\courssub{Correction collective d'un CV défectueux}

Le professeur distribue ou projette le CV ci-dessous, volontairement truffé d'erreurs. La classe l'analyse collectivement, critère par critère.

\begin{exoresolu}[Identifier et corriger les erreurs d'un CV défectueux]
Énoncé. Voici des extraits du CV de « MBALLA Jean », candidat au poste de technicien de maintenance. Relevez toutes les erreurs (structure, langue, contenu) et proposez une correction.\par
\emph{« CV — mballa jean, né le 12/03/2005 a Yaoundé. Expérience professionnelle : 2024 : stage a la SABC ; 2023 : vendeur au marché. Formation : BEPC en 2021, Probatoire technique en 2024, Baccalauréat technique en 2025. J'ai travailler dans un atelier qui réparait les machines, dont j'ai appris beaucoup de choses. Compétances : informatique. Je suis le meilleur candidat que vous pouvez recruter. »}
\tcblower
Solution détaillée. Erreurs relevées :\par
1. \textbf{En-tête incomplet} : pas d'adresse, ni téléphone, ni courriel ; le nom s'écrit avec majuscules : \emph{MBALLA Jean}.\par
2. \textbf{Ordre chronologique fautif} : la formation doit aller de la plus récente à la plus ancienne : Baccalauréat (2025), Probatoire (2024), BEPC (2021).\par
3. \textbf{Orthographe et grammaire} : \emph{« a Yaoundé »} $\rightarrow$ \emph{à Yaoundé} ; \emph{« stage a la SABC »} $\rightarrow$ \emph{stage à la SABC} ; \emph{« j'ai travailler »} $\rightarrow$ \emph{j'ai travaillé} ; \emph{« compétances »} $\rightarrow$ \emph{compétences}.\par
4. \textbf{Pronom relatif mal employé} : \emph{« un atelier ... dont j'ai appris beaucoup de choses »} est incorrect : on dit \emph{« un atelier où j'ai appris beaucoup de choses »} (lieu $\rightarrow$ \emph{où}) ou \emph{« un atelier dans lequel... »}.\par
5. \textbf{Phrase maladroite et prétentieuse} : \emph{« Je suis le meilleur candidat que vous pouvez recruter »} $\rightarrow$ tourner sobrement : \emph{« Motivé et rigoureux, je saurai m'intégrer rapidement à votre équipe. »} De plus, \emph{« que vous pouvez »} $\rightarrow$ \emph{« que vous puissiez »} (subjonctif après superlatif) — mais la suppression totale est préférable.\par
6. \textbf{Rubrique « compétences » pauvre} : préciser (\emph{traitement de texte, tableur, câblage réseau}) et ajouter langues et permis.
\end{exoresolu}

\courssub{Remédiation ciblée}

La remédiation reprend, en petits groupes, les difficultés observées pendant l'évaluation croisée et la correction collective.

\begin{itemize}
\item \textbf{Difficulté 1 — Structure externe.} Exercice : remettre en page un CV « en vrac » (toutes les informations données en désordre) en une page, avec rubriques titrées et alignements.
\item \textbf{Difficulté 2 — Ordre des rubriques.} Exercice : classer cinq diplômes et trois expériences dans l'ordre antichronologique.
\item \textbf{Difficulté 3 — Pronoms relatifs.} Exercice : compléter par \emph{qui, que, où, dont} : « L'entreprise \ldots{} j'ai effectué mon stage est située à Douala. Le poste \ldots{} je sollicite exige le permis B. Le chef d'atelier, \ldots{} j'ai beaucoup appris, m'a recommandé. Les outils \ldots{} je me servais étaient modernes. » (Réponses : \emph{où ; que ; dont ; dont}.)
\item \textbf{Difficulté 4 — Orthographe.} Dictée ciblée sur le vocabulaire du CV : \emph{expérience professionnelle, compétences, maîtrise, disponibilité, recrutement, candidature}.
\item \textbf{Difficulté 5 — Mise en page.} Atelier : transformer un paragraphe continu en rubriques à puces alignées.
\end{itemize}

\begin{exercice}[Remédiation — à faire individuellement]
1. Réécris cette phrase en deux versions, avec un pronom relatif puis avec une locution adverbiale de liaison : \emph{« J'ai fait un stage à la CAMTEL. Ce stage m'a formé aux réseaux. »}\par
2. Corrige : \emph{« j'ai effectuer un stage dont je suis fier »}.\par
3. Rédige la rubrique « Formation » d'un élève qui a obtenu le BEPC en 2022, le Probatoire F3 en 2025 et prépare le Baccalauréat F3 en 2026.
\end{exercice}

\begin{corrige}[Exercice de remédiation]
1. Avec pronom relatif : \emph{« J'ai fait un stage à la CAMTEL, qui m'a formé aux réseaux. »} Avec locution de liaison : \emph{« J'ai fait un stage à la CAMTEL ; par la suite, j'ai été formé aux réseaux. »}\par
2. \emph{« J'ai effectué un stage dont je suis fier. »} (infinitif $\rightarrow$ participe passé après \emph{avoir} ; \emph{dont} est ici correct car on est fier \emph{de} quelque chose.)\par
3. \textbf{Formation} : 2026 : Baccalauréat technique F3 en préparation, Lycée technique de \ldots{} ; 2025 : Probatoire technique F3 ; 2022 : BEPC. (Ordre antichronologique respecté.)
\end{corrige}

\courssub{Bilan des acquis : auto-évaluation}

\begin{exemplebox}[Barème type d'évaluation d'un CV à l'examen]
\begin{center}
\begin{tabularx}{\linewidth}{|X|c|}
\hline
\rowcolor{popblue}\textcolor{white}{\textbf{Critère}} & \textcolor{white}{\textbf{Pondération}} \\
\hline
Contenu (informations pertinentes, adaptées à l'offre) & 40 \% \\
\hline
Cohérence (ordre des rubriques, logique, liaisons) & 30 \% \\
\hline
Présentation (mise en page, lisibilité, une page) & 20 \% \\
\hline
Langue (orthographe, grammaire, vocabulaire) & 10 \% \\
\hline
\end{tabularx}
\end{center}
Ce barème montre qu'un CV beau mais vide ne rapporte rien : le contenu et la cohérence pèsent 70 \% de la note.
\end{exemplebox}

Pour clore la leçon, chaque élève coche honnêtement la liste des objectifs :

\begin{itemize}
\item[$\square$] Je sais identifier la structure externe d'un CV (en-tête, rubriques, mise en page).
\item[$\square$] Je sais construire les rubriques internes (formation, expérience, compétences, divers) dans le bon ordre.
\item[$\square$] J'emploie correctement les pronoms relatifs et les adverbes de liaison.
\item[$\square$] Je sais adapter un CV à une offre d'embauche précise.
\item[$\square$] Je sais évaluer et corriger un CV à l'aide d'une grille.
\end{itemize}

Tout objectif non validé donne lieu à une reprise individuelle avec le professeur avant l'évaluation sommative.

\courssub{Prolongement (complément utile au Bac) : CV, lettre de motivation et entretien}

\begin{attention}[Complément utile au Baccalauréat technique]
Ce prolongement dépasse légèrement le strict objet de l'unité, mais il est \textbf{fortement conseillé} : les épreuves de production d'écrits peuvent articuler plusieurs textes fonctionnels.
\end{attention}

Le CV ne voyage jamais seul : il s'accompagne d'une \cle{lettre de motivation}, qui reprend en quelques paragraphes argumentés les points forts du CV (le schéma classique : \emph{vous — moi — nous} : l'entreprise, mes atouts, notre collaboration possible). Cohérence obligatoire : dates, diplômes et expériences doivent être \textbf{strictement identiques} dans les deux documents. Enfin, le CV prépare l'\cle{entretien d'embauche} : tout ce que tu écris peut être questionné oralement — n'écris donc rien que tu ne puisses expliquer et prouver.

\begin{savaistu}[Le CV, passeport des concours camerounais]
De nombreux concours d'entrée dans les grandes écoles camerounaises — ENAM, IFCPA, écoles de génie (ENSP), écoles de santé — exigent un CV dans le dossier de candidature. Un CV clair, honnête et sans faute y est souvent le premier document lu par le jury : savoir le rédiger en Terminale technique, c'est déjà préparer son insertion professionnelle et ses concours.
\end{savaistu}

\begin{popfigure}
\begin{center}
\begin{tikzpicture}[font=\small,
box/.style={draw=popblue,fill=popblueL,rounded corners=2pt,align=center,minimum width=3.1cm,minimum height=0.9cm},
arr/.style={-{Stealth[length=2.5mm]},thick,popdark}]
\node[box, align=center] (cv) at (0,0){\textbf{CV}\\ les faits};
\node[box, align=center] (lm) at (4.6,0){\textbf{Lettre de motivation}\\ l'argumentation};
\node[box, align=center] (en) at (9.2,0){\textbf{Entretien}\\ la parole};
\draw[arr] (cv) -- (lm) node[midway,above,popdark]{cohérence};
\draw[arr] (lm) -- (en) node[midway,above,popdark]{preuves};
\draw[arr,poporange,dashed] (en.south) .. controls (4.6,-1.4) .. (cv.south) node[midway,below,poporange]{tout ce qui est écrit peut être questionné};
\end{tikzpicture}
\end{center}
\legende{Articulation des trois étapes d'une candidature : le CV présente les faits, la lettre argumente, l'entretien vérifie oralement.}
\end{popfigure}

\begin{popfigure}
\begin{center}
\begin{tikzpicture}[font=\footnotesize]
\draw[draw=popdark,fill=popcream,rounded corners=3pt] (0,0) rectangle (10.6,13.2);
\node[anchor=north west,font=\bfseries\large,popblue] at (0.4,12.9) {ETO'A Marie Chantal};
\node[anchor=north west] at (0.4,12.3) {Née le 05/04/2006 à Bafoussam — Tél. : 6 90 00 00 00 — marie.etoa@mail.cm};
\draw[popblue,thick] (0.4,11.9) -- (10.2,11.9);
\node[anchor=north west,font=\bfseries,popblue] at (0.4,11.6) {FORMATION};
\node[anchor=north west] at (0.4,11.1) {2026 : Baccalauréat technique F3 en préparation, Lycée technique de Bafoussam};
\node[anchor=north west] at (0.4,10.7) {2025 : Probatoire technique F3 \hfill 2022 : BEPC};
\draw[popline] (0.4,10.35) -- (10.2,10.35);
\node[anchor=north west,font=\bfseries,popblue] at (0.4,10.05) {EXPÉRIENCE PROFESSIONNELLE};
\node[anchor=north west] at (0.4,9.6) {2025 : stage de maintenance réseaux à la CAMTEL, où j'ai appris le câblage};
\node[anchor=north west] at (0.4,9.2) {2024 : dépannage informatique dans un cybercafé, ce qui m'a formée au contact client};
\draw[popline] (0.4,8.85) -- (10.2,8.85);
\node[anchor=north west,font=\bfseries,popblue] at (0.4,8.55) {COMPÉTENCES};
\node[anchor=north west] at (0.4,8.1) {Câblage réseau ; traitement de texte et tableur ; anglais (niveau scolaire) ; permis B};
\draw[popline] (0.4,7.75) -- (10.2,7.75);
\node[anchor=north west,font=\bfseries,popblue] at (0.4,7.45) {DIVERS};
\node[anchor=north west] at (0.4,7.0) {Sport (handball), lecture, membre du club informatique du lycée};
% annotations
\node[draw=poporange,fill=poporangeL,rounded corners=2pt,align=left,anchor=west] (a1) at (11.0,12.5) {1. En-tête complet};
\draw[-{Stealth},poporange] (a1.west) -- (10.4,12.6);
\node[draw=poporange,fill=poporangeL,rounded corners=2pt,align=left,anchor=west] (a2) at (11.0,10.9) {2. Ordre antichronologique};
\draw[-{Stealth},poporange] (a2.west) -- (10.4,10.9);
\node[draw=poporange,fill=poporangeL,rounded corners=2pt,align=left,anchor=west] (a3) at (11.0,9.4) {3. Relatifs \emph{où}, \emph{ce qui}};
\draw[-{Stealth},poporange] (a3.west) -- (10.4,9.4);
\node[draw=poporange,fill=poporangeL,rounded corners=2pt,align=left,anchor=west] (a4) at (11.0,8.1) {4. Compétences ciblées sur l'offre};
\draw[-{Stealth},poporange] (a4.west) -- (10.4,8.1);
\end{tikzpicture}
\end{center}
\legende{Modèle de CV corrigé et annoté (référence de la classe) : 1 en-tête complet ; 2 formation de la plus récente à la plus ancienne ; 3 expériences datées avec pronoms relatifs corrects ; 4 compétences adaptées à l'offre CAMTEL-SERVICES.}
\end{popfigure}

\begin{aretenir}[L'essentiel de la section]
L'activité d'intégration suit un cycle complet : \textbf{produire} (rédiger son CV en réponse à une offre), \textbf{évaluer} (grille à 10 critères, échange croisé), \textbf{rendre compte} (présenter et justifier sa démarche), \textbf{remédier} (reprendre structure, orthographe, relatifs, mise en page), puis \textbf{bilan} (auto-évaluation des objectifs). Ce cycle est exactement celui de l'épreuve de production d'écrits au Baccalauréat technique : rédiger, relire, corriger — et pouvoir justifier chacun de ses choix.
\end{aretenir}

\newpage

\coursec{Activité d'intégration}

\courssub{Objectifs de l'activité}

À l'issue de cette activité d'intégration, l'élève sera capable de :
\begin{itemize}
\item analyser une offre d'emploi authentique pour en dégager le \cle{profil exigé} (diplôme, compétences, expérience, qualités personnelles) ;
\item établir l'\cle{inventaire de son parcours} (réel ou à partir d'une fiche-personnage) et sélectionner les informations pertinentes ;
\item rédiger un \cle{CV complet} respectant la structure externe (mise en page, rubriques, ordre antichronologique) et la structure interne (état civil, formation, expérience professionnelle, compétences, centres d'intérêt) ;
\item employer correctement au moins trois \cle{outils de liaison} (adverbes de liaison, locutions adverbiales, pronoms relatifs) dans les parties rédigées du CV ;
\item évaluer un CV à l'aide d'une \cle{grille d'évaluation} construite collectivement et \cle{justifier} un choix de recrutement ;
\item présenter et défendre oralement sa démarche de rédaction devant un jury simulé.
\end{itemize}

\courssub{Situation}

\textbf{« Décroche ton poste ! » — Simulation de recrutement en classe.}

Le lycée technique de la localité organise, en partenariat avec le Club des jeunes professionnels, une journée de simulation de recrutement intitulée « Décroche ton poste ! ». Trois entreprises camerounaises (réelles ou fictives) ont publié les offres d'emploi suivantes.

\begin{exemplebox}[Offre n° 1 — Technicien de maintenance (SODECOTON, Garoua)]
La SODECOTON (Société de développement du coton du Cameroun) recrute un(e) technicien(ne) de maintenance industrielle pour son usine de Garoua. Profil exigé : Baccalauréat technique (série F) ou BTS en maintenance industrielle ; au moins \textbf{1 an d'expérience} (stages acceptés) ; maîtrise de la lecture de schémas électriques et mécaniques ; sens de l'organisation et disponibilité. Dossier : CV + lettre de motivation, à déposer avant le 30 juin.
\end{exemplebox}

\begin{exemplebox}[Offre n° 2 — Dessinateur-projeteur (cabinet d'architecture, Douala)]
Un cabinet d'architecture de Bonanjo (Douala) recrute un(e) dessinateur(trice)-projeteur. Profil exigé : Baccalauréat technique (série F4 — Génie civil) ou DUT ; maîtrise du logiciel AutoCAD ; rigueur, créativité et capacité à travailler en équipe ; une première expérience en cabinet serait un atout. Dossier : CV + lettre de motivation.
\end{exemplebox}

\begin{exemplebox}[Offre n° 3 — Comptable (PME, Yaoundé)]
Une PME de négoce de la commune de Yaoundé III recrute un(e) comptable. Profil exigé : Baccalauréat technique (série G — Techniques comptables) ou BTS en comptabilité ; maîtrise du tableur Excel et d'un logiciel de comptabilité ; honnêteté, discrétion, sens du détail ; débutant accepté. Dossier : CV + lettre de motivation.
\end{exemplebox}

\textbf{Règles du jeu.} Chaque élève choisit \textbf{une} offre. Ceux qui le souhaitent rédigent leur CV à partir de leur parcours réel ; les autres tirent au sort une \cle{fiche-personnage} (exemple : « Marlyse Abena, 19 ans, Baccalauréat G obtenu en 2025 au lycée technique de Yaoundé, stage de 2 mois à la SNI, membre du club informatique, parlant français et anglais »). Les CV rédigés sont déposés anonymement dans la \cle{boîte à candidatures}. Un jury de 6 élèves (2 par offre), jouant le rôle de recruteurs, évalue chaque CV à l'aide de la grille construite en classe (sur 20 points), sélectionne les \textbf{trois meilleurs CV par offre} et justifie ses choix devant la classe. Les auteurs sélectionnés présentent leur démarche en 3 minutes. La séance se clôt par un compte rendu collectif et une remédiation.

\courssub{Consignes}

\begin{enumerate}
\item \textbf{Analyse de l'offre.} Souligne dans l'offre choisie les éléments du profil exigé et classe-les dans un tableau à trois colonnes : diplôme/formation, compétences techniques, qualités personnelles.
\item \textbf{Inventaire du parcours.} À partir de ton parcours réel ou de ta fiche-personnage, liste toutes les informations disponibles (état civil, études, stages, compétences, loisirs), puis barre celles qui ne sont pas utiles pour l'offre. Justifie deux de tes choix.
\item \textbf{Plan du CV.} Présente le squelette de ton CV : titre, ordre des rubriques, ordre des informations dans chaque rubrique (antichronologique). Explique pourquoi tu places une rubrique avant une autre.
\item \textbf{Rédaction.} Rédige ton CV complet sur une page. Rédige en phrases les rubriques « Formation » ou « Expérience » en employant \textbf{au moins trois outils de liaison} différents (par exemple : \emph{d'abord, ensuite, en outre, c'est pourquoi, qui, que, où}).
\item \textbf{Relecture.} Vérifie ton CV à l'aide de la grille d'évaluation : présentation soignée, rubriques complètes, ordre respecté, liaisons correctes, adaptation à l'offre, orthographe.
\item \textbf{Évaluation croisée.} En tant que membre du jury, attribue une note sur 20 à chaque CV de ta pile et rédige une phrase de justification pour le CV que tu retiens.
\item \textbf{Défense orale.} Si ton CV est sélectionné, présente en 3 minutes : l'offre choisie, les informations retenues et écartées, les outils de liaison utilisés, les difficultés rencontrées.
\item \textbf{Compte rendu.} Note les trois erreurs les plus fréquentes relevées par le jury et propose pour chacune une correction.
\end{enumerate}

\courssub{Ressources à mobiliser}

\begin{aretenir}[Ce qu'il faut se rappeler]
\begin{itemize}
\item \textbf{Structure externe} : un CV tient sur \cle{une page}, est aéré, sans ratures, avec un titre (le poste visé), des rubriques nettement séparées, une photo facultative.
\item \textbf{Structure interne (rubriques)} : \cle{état civil} (nom, date et lieu de naissance, nationalité, contacts), \cle{formation} (diplômes du plus récent au plus ancien), \cle{expérience professionnelle} (stages, emplois, ordre antichronologique), \cle{compétences} (langues, informatique), \cle{centres d'intérêt}.
\item \textbf{Outils de liaison} : adverbes de liaison (\emph{d'abord, ensuite, enfin, cependant}), locutions adverbiales (\emph{en outre, de plus, c'est pourquoi, par la suite}), pronoms relatifs (\emph{qui, que, où, dont}).
\item \textbf{Adaptation à l'offre} : ne retenir que les informations qui répondent au profil exigé ; mettre en avant l'expérience ou la compétence la plus demandée.
\item \textbf{Justification} : toute décision (choix d'une rubrique, ordre, information écartée) doit pouvoir être expliquée par une phrase : « J'ai choisi... parce que... ».
\end{itemize}
\end{aretenir}

\begin{center}
\begin{tabularx}{\linewidth}{|l|X|c|}
\hline
\rowcolor{popblueL} \textbf{Critère} & \textbf{Ce que le jury vérifie} & \textbf{Barème} \\
\hline
Présentation & Page unique, clarté, soin, titre du poste & 4 pts \\
\hline
Rubriques & Les 5 rubriques sont présentes et complètes & 4 pts \\
\hline
Ordre & Ordre antichronologique respecté & 3 pts \\
\hline
Adaptation & Le CV répond au profil exigé par l'offre & 4 pts \\
\hline
Langue & Au moins 3 outils de liaison corrects, orthographe & 5 pts \\
\hline
\textbf{Total} & & \textbf{20 pts} \\
\hline
\end{tabularx}
\end{center}

\courssub{Démarche et corrigé commenté}

\begin{exoresolu}[Corrigé commenté : le CV de Marlyse pour l'offre n° 3 (comptable, Yaoundé)]
Énoncé : à partir de la fiche-personnage de Marlyse Abena, rédiger un CV adapté à l'offre de comptable dans la PME de Yaoundé III et justifier la démarche.
\tcblower
\textbf{Étape 1 — Analyse de l'offre.} Le profil exigé se décompose ainsi :
\begin{itemize}
\item formation : Baccalauréat G ou BTS comptabilité ;
\item compétences : Excel, logiciel de comptabilité ;
\item qualités : honnêteté, discrétion, sens du détail ; débutant accepté.
\end{itemize}
\textbf{Étape 2 — Inventaire et tri.} Marlyse possède : Baccalauréat G (2025), BEPC (2022), stage de 2 mois à la SNI (service comptabilité), club informatique, bilinguisme. Toutes ces informations sont utiles ; on écarte en revanche des détails sans rapport avec le poste (par exemple « aime la danse » peut rester en centres d'intérêt, mais « championne de handball scolaire » n'apporte rien ici). Justification : un CV doit être \emph{sélectif} ; seules les informations qui servent la candidature sont retenues.
\par
\textbf{Étape 3 — Plan.} Titre : « Comptable débutante ». Ordre des rubriques : état civil, formation (le diplôme est le point fort d'un débutant), expérience, compétences, centres d'intérêt.
\par
\textbf{Étape 4 — Rédaction (extrait des rubriques rédigées).}
\par
\emph{Formation :} « J'ai d'abord obtenu le BEPC en 2022 au collège bilingue d'Essos ; ensuite, j'ai préparé le Baccalauréat technique série G au lycée technique de Yaoundé, où j'ai été reçue en 2025 avec mention assez bien. »
\par
\emph{Expérience :} « En outre, j'ai effectué un stage de deux mois à la SNI, dans le service comptabilité, stage qui m'a permis de pratiquer la saisie des factures sur Excel. C'est pourquoi je me sens prête à tenir la comptabilité d'une PME. »
\par
Les outils de liaison employés sont : \emph{d'abord} et \emph{ensuite} (adverbes de liaison), \emph{en outre} et \emph{c'est pourquoi} (locutions adverbiales), \emph{où} et \emph{qui} (pronoms relatifs) : la consigne des trois outils est donc largement remplie.
\par
\textbf{Étape 5 — Justification devant le jury.} « J'ai placé la formation avant l'expérience parce que mon diplôme est mon principal atout, mon expérience étant courte. J'ai conservé le club informatique parce qu'il prouve ma maîtrise de l'outil informatique exigée par l'offre. »
\par
\textbf{Étape 6 — Évaluation.} Selon la grille : présentation 4/4, rubriques 4/4, ordre 3/3, adaptation 4/4, langue 4/5 (une faute d'accord) : total \textbf{19/20}. Le jury retient ce CV « parce qu'il répond point par point au profil exigé et que les liaisons rendent la lecture fluide ».
\end{exoresolu}

\begin{attention}[Erreurs fréquentes relevées pendant la simulation]
\begin{itemize}
\item Mélanger les dates (diplôme ancien placé avant le récent) : toujours vérifier l'ordre \cle{antichronologique}.
\item Donner des informations inutiles (taille, religion, nom des parents) : le CV n'est pas un acte d'état civil.
\item Oublier le titre du CV : le recruteur doit voir immédiatement le poste visé.
\item Juxtaposer les phrases sans liaison : employer \emph{d'abord, ensuite, en outre, c'est pourquoi} et les pronoms relatifs pour enchaîner.
\item Recopier l'offre sans adapter : le CV doit répondre au profil, pas le répéter.
\end{itemize}
\end{attention}

\courssub{Bilan de l'activité}

Cette activité d'intégration a permis de mobiliser simultanément tous les savoirs de la leçon dans une situation concrète et motivante. L'\cle{analyse de l'offre} a montré qu'un bon CV commence avant la rédaction : il faut d'abord identifier le profil exigé. Le \cle{tri des informations} a fait comprendre qu'un CV est un document sélectif et non exhaustif. La rédaction a consolidé la maîtrise de la \cle{structure externe} (page unique, rubriques, ordre antichronologique) et de la \cle{structure interne} (contenu de chaque rubrique), tandis que la consigne sur les liaisons a réinvesti les acquis de morphosyntaxe (adverbes de liaison, locutions adverbiales, pronoms relatifs) dans une production réelle. Enfin, l'évaluation par les pairs et la défense orale ont développé deux savoir-faire essentiels du programme : appliquer les notions à des situations de la vie courante et \cle{justifier une démarche}. La remédiation collective, à partir des erreurs relevées dans la boîte à candidatures, a permis à chaque élève de corriger son CV avant de le conserver dans son portfolio de fin de cycle, prêt à servir lors d'une véritable recherche d'emploi ou de stage.

\newpage

\coursec{Méthodes et exercices résolus}

\coursec{Produire des écrits utilitaires : le CV}

\courssub{Découverte et réception du CV}

\begin{methode}[Analyser un CV : repérer la nature, la destination et l'organisation]
    \textbf{Objectif :} Comprendre ce qu'est un CV, à qui il s'adresse et comment il est structuré globalement.
    \begin{enumerate}
        \item \textbf{Identifier la nature du document :} C'est un \emph{texte fonctionnel} (ou utilitaire) à visée professionnelle. Il ne raconte pas une histoire (narratif) mais présente des faits (informatif).
        \item \textbf{Déterminer la destination :} À qui s'adresse-t-il ? Au \emph{recruteur} (employeur, DRH, chef de service). Il doit être lu en \textbf{moins de 30 secondes} (lecture diagonale).
        \item \textbf{Repérer les deux structures :}
              \begin{itemize}
                  \item \cle{Structure externe} : Mise en page, police, espaces, mise en forme (gras, puces, alignement). C'est l'aspect visuel immédiat.
                  \item \cle{Structure interne} : Le contenu organisé en \emph{rubriques} standardisées (Identité, Formation, Expériences, Compétences, Centres d'intérêt).
              \end{itemize}
        \item \textbf{Vérifier les 3C du CV efficace :} \textbf{C}oncis (1 page max), \textbf{C}lair (lisible, aéré), \textbf{C}ohérent (adapté au poste visé).
        \item \textbf{Reflexe Bac :} Face à un sujet d'examen, lire le CV support en se demandant : « Quel poste vise ce candidat ? Les rubriques sont-elles dans l'ordre logique ? L'essentiel saute-t-il aux yeux ? »
    \end{enumerate}
\end{methode}

\begin{exoresolu}[Analyse d'un CV type : détection des forces et faiblesses]
    \textbf{Énoncé :} Vous êtes membre d'un jury de recrutement pour un poste de \emph{Technicien en maintenance industrielle} (Bac Pro / BT Maintenance). On vous soumet le CV ci-dessous (texte brut). Identifiez 3 points forts et 3 points faibles structurels. Proposez une réorganisation des rubriques.

    \textit{CV de M. Ngono Paul}
    Né le 12/05/2000 à Douala. Célibataire. Permis B.
    Expériences : Stage 2 mois (2022) : Garage Auto Modern, Douala (vidanges, freins). Job d'été 2021 : Manutentionnaire, Socapalm. Stage 1 mois (2020) : Électricité bâtiment, père.
    Formation : 2023 : BT Maintenance des Systèmes Automatisés, Lycée Technique de Douala. 2020 : Probatoire Série F4. 2018 : BEPC.
    Compétences : Maintenance préventive/curative automatismes, Lecture plans schémas, Soudure TIG/MIG, Programmation API (Siemens S7-1200 bases), Français (langue maternelle), Anglais (lu, écrit, parlé), Bureautique (Word, Excel).
    Centres d'intérêt : Football (capitaine équipe quartier), Musique (DJ événements), Lecture romans policiers.

    \tcblower
    \textbf{Solution détaillée :}

    \textbf{1. Points forts (Structure interne / Contenu) :}
    \begin{itemize}
        \item \textbf{Compétences techniques ciblées :} La rubrique « Compétences » liste des savoir-faire précis attendus pour le poste (API Siemens, lecture plans, soudure TIG/MIG). C'est le \cle{cœur de métier}.
        \item \textbf{Diplôme pertinent :} Le BT Maintenance des Systèmes Automatisés (2023) correspond exactement au niveau requis (Bac Pro / BT).
        \item \textbf{Permis B + Mobilité :} Indiqué en en-tête, c'est un atout majeur pour la maintenance (déplacements sites clients).
    \end{itemize}

    \textbf{2. Points faibles (Structure externe / Organisation) :}
    \begin{itemize}
        \item \textbf{Ordre chronologique inversé non respecté :} Les expériences et formations sont mélangées ou dans le désordre (ex: stage 2020 avant 2021). Le recruteur veut voir le \textbf{plus récent en premier}.
        \item \textbf{Présentation en « bloc » unique (texte continu) :} Absence de puces, de gras sur les intitulés de poste/entreprise, de séparation visuelle des rubriques. Lecture diagonale impossible. \textbf{Échec des 3C (Clarté).}
        \item \textbf{Expériences non valorisées :} « Stage électricité bâtiment, père » manque de cadre formel (nom entreprise ?). « Manutentionnaire » : pas de lien fait avec la maintenance (port de charges ? respect consignes sécurité ?). Verbes d'action absents.
    \end{itemize}

    \textbf{3. Proposition de réorganisation (Structure standard « anti-chronologique ») :}
    \begin{enumerate}
        \item \textbf{EN-TÊTE :} NOM Prénom, Adresse, Tél, Email, LinkedIn (si pro), Permis B, Âge/Ville.
        \item \textbf{TITRE / ACCROCHE :} \emph{Technicien Maintenance Industrielle (Automatismes) — BT MSMA — Permis B — Mobile}.
        \item \textbf{FORMATION (du plus récent au plus ancien) :} BT MSMA (2023) $\rightarrow$ Probatoire F4 (2020) $\rightarrow$ BEPC (2018). Préciser options, projets de fin d'études.
        \item \textbf{EXPÉRIENCES PROFESSIONNELLES (du plus récent au plus ancien) :}
              \begin{itemize}
                  \item 2022 (2 mois) : \textbf{Stagiaire Maintenance} — Garage Auto Modern, Douala. \textit{Missions :} Maintenance préventive/curative véhicules, diagnostic valise, gestion stock pièces.
                  \item 2021 (2 mois) : \textbf{Agent de Maintenance / Manutention} — Socapalm. \textit{Missions :} Surveillance convoyeurs, graissage, alerte pannes, port EPI.
                  \item 2020 (1 mois) : \textbf{Stagiaire Électricité} — Entreprise [Nom réel si possible], Douala. \textit{Missions :} Tirage câbles, pose appareillage, lecture schémas unifilaires.
              \end{itemize}
        \item \textbf{COMPÉTENCES TECHNIQUES :} Regrouper par domaines (Automatismes, Électricité, Mécanique, Informatique/CAO, Langues).
        \item \textbf{CENTRES D'INTÉRÊT :} Sélectionner ceux montrant des \cle{soft skills} (Capitaine foot = leadership/esprit d'équipe ; DJ = gestion événementiel/technique son).
    \end{enumerate}

    \textbf{Conclusion :} Un CV doit être une \emph{publicité} pour le candidat. La structure externe (mise en page aérée, puces, gras) guide l'œil vers la structure interne (compétences/postes) qui prouve l'adéquation poste/profil.
\end{exoresolu}

\courssub{Structure externe du CV}

\begin{methode}[Maîtriser les codes de la mise en page professionnelle]
    \textbf{Objectif :} Produire un document visuellement impeccable, lisible en 30 secondes, compatible ATS (logiciels de tri automatique).
    \begin{enumerate}
        \item \textbf{Format \& Police :} Format \textbf{A4}, 1 page recto (sauf profil senior). Police \textbf{sans empattement} (sans serif) : \emph{Calibri, Arial, Helvetica, Roboto, Open Sans}. Taille : \textbf{10 à 11 pt} pour le corps, \textbf{13--14 pt} pour les titres de rubriques, \textbf{16--18 pt} pour le Nom.
        \item \textbf{Marges \& Espacement :} Marges \textbf{1,5 à 2 cm}. Interligne \textbf{1,15}. Espace \textbf{6--8 pt} après chaque rubrique. \textbf{Alignement justifié} ou \textbf{aligné à gauche} (jamais centré pour le corps).
        \item \textbf{Hiérarchisation visuelle (le « guidage de l'œil ») :}
              \begin{itemize}
                  \item \textbf{Gras} sur : Intitulés de poste, Noms d'entreprises, Diplômes, Mots-clés techniques.
                  \item \textbf{Puces (bullets)} pour les listes (missions, compétences). Symbole simple : \textbullet\ ou --\ ou $\triangleright$.
                  \item \textbf{Séparateurs} : Filets fins horizontaux (0,5 pt) ou changement de couleur de fond léger (gris très clair \texttt{popgrayL}) pour délimiter les blocs.
              \end{itemize}
        \item \textbf{Couleur :} \textbf{Une seule couleur d'accent} (bleu marine \texttt{popblue}, anthracite \texttt{popdark}) pour les titres de rubriques ou la barre latérale. Fond blanc. \textit{Interdit :} polices fantaisie, photos de vacances, cadres lourds, cliparts.
        \item \textbf{Photo :} Au Cameroun, \textbf{photo d'identité professionnelle} (fond neutre, tenue correcte, sourire léger) en haut à droite ou gauche (3,5 $\times$ 4,5 cm). Optionnelle mais recommandée dans le privé.
        \item \textbf{Export :} Toujours \textbf{PDF} (nommé \texttt{Nom\_Prenom\_CV\_Poste.pdf}). Vérifier l'affichage sur mobile.
    \end{enumerate}
\end{methode}

\begin{exoresolu}[Correction de mise en page : du « brouillon » au « pro »]
    \textbf{Énoncé :} Voici un extrait de CV mal mis en page (police Times New Roman 12, justifié serré, pas de puces, tout en noir, marges 1 cm). Appliquez la méthode pour le transformer. Montrez le code LaTeX (environnement \texttt{tabularx} ou description) qui produit une version propre.

    \textit{Extrait brut :}
    FORMATION 2023 BT Maintenance Systemes Automatises Lycée Technique de Douala Option Productique 2020 Probatoire F4 Lycée de New-Bell 2018 BEPC College Libermann
    EXPERIENCES 2022 Stage Garage Auto Modern Douala Vidanges Freins Diagnostic 2021 Manutentionnaire Socapalm Port de charges Securite

    \tcblower
    \textbf{Solution détaillée :}

    \textbf{1. Application de la méthode :}
    \begin{itemize}
        \item Police : \texttt{\textbackslash setsansfont\{Calibri\}} (XeLaTeX) $\rightarrow$ \textbf{Sans serif}.
        \item Structure : Séparation nette \textbf{FORMATION} / \textbf{EXPÉRIENCES}.
        \item Ordre : Antéchronologique (2023 en haut).
        \item Puces : Ajout de \textbullet\ pour chaque ligne d'action.
        \item Gras : Sur diplômes, établissements, postes, entreprises.
        \item Espacement : \texttt{\textbackslash smallskip} entre entrées.
    \end{itemize}

    \textbf{2. Code LaTeX proposée (corps de document, compilable avec le préambule OnBuch+) :}
    \begin{verbatim}
    % Dans le préambule :     
    % \definecolor{cvblue}{HTML}{1B3A5C} % Bleu marine pro

    \begin{center}
        \textbf{\LARGE NGONO Paul} \\ \vspace{2mm}
        \textcolor{cvblue}{\rule{\linewidth}{0.8pt}} \\ \vspace{3mm}
    \end{center}

    \textbf{\textcolor{cvblue}{FORMATION}} \hfill \textit{Antéchronologique}
    \begin{itemize}
        \item \textbf{2023 -- BT Maintenance des Systèmes Automatisés} (Option Productique) \\
              \textit{Lycée Technique de Douala} -- Projet fin d'études : « Automatisation ligne convoyage »
        \item \textbf{2020 -- Probatoire Série F4 (Électrotechnique)} \\
              \textit{Lycée de New-Bell} -- Mention Assez Bien
        \item \textbf{2018 -- BEPC} \\
              \textit{Collège Libermann}
    \end{itemize}

    \textbf{\textcolor{cvblue}{EXPÉRIENCES PROFESSIONNELLES}} \hfill \textit{Antéchronologique}
    \begin{itemize}
        \item \textbf{03/2022 -- 05/2022 : Stagiaire Technicien Maintenance} \\
              \textit{Garage Auto Modern, Douala} \\
              \textbullet\ Maintenance préventive/curative véhicules légers (vidange, freins, distribution). \\
              \textbullet\ Diagnostic électronique (valise multimarque). \\
              \textbullet\ Gestion du stock pièces de rechange (inventaire mensuel).
        \item \textbf{07/2021 -- 08/2021 : Agent de Maintenance / Manutentionnaire} \\
              \textit{SOCAPALM, Dibombari} \\
              \textbullet\ Surveillance et graissage convoyeurs chaîne de production huile de palme. \\
              \textbullet\ Détection pannes niveau 1, alerte équipe maintenance. \\
              \textbullet\ Respect strict consignes sécurité (EPI, consignation).
    \end{itemize}
    \end{verbatim}

    \textbf{3. Résultat visuel attendu :}
    \begin{center}
    \begin{tabularx}{\linewidth}{|X|}\hline
    \rowcolor{popblueL} \textbf{FORMATION} \\
    \textbf{2023 -- BT Maintenance Systèmes Automatisés} (Option Productique) \\
    \textit{Lycée Technique de Douala} -- Projet : Automatisation ligne convoyage \\ \hline
    \rowcolor{popblueL} \textbf{EXPÉRIENCES PROFESSIONNELLES} \\
    \textbf{03/2022 -- 05/2022 : Stagiaire Technicien Maintenance} \\
    \textit{Garage Auto Modern, Douala} \\
    \textbullet\ Maintenance préventive/curative, Diagnostic valise, Gestion stock \\ \hline
    \end{tabularx}
    \end{center}

    \textbf{Conclusion :} La structure externe (police, puces, gras, espacement, couleur unique) transforme une liste illisible en un document professionnel scannable en 10 secondes.
\end{exoresolu}

\courssub{Structure interne du CV : les rubriques}

\begin{methode}[Construire le squelette : rubriques obligatoires, recommandées, optionnelles]
    \textbf{Objectif :} Savoir ordonner et remplir chaque bloc d'information selon sa pertinence par rapport au poste visé.
    \begin{enumerate}
        \item \textbf{Rubriques OBLIGATOIRES (le « socle ») :}
              \begin{enumerate}
                  \item \textbf{En-tête / Identité :} Nom, Prénom, Adresse (Ville suffit), Téléphone, Email pro (\texttt{prenom.nom@email.com}), LinkedIn (URL personnalisée), Permis(s), Âge/Date de naissance (usage Cameroun), Photo pro.
                  \item \textbf{Titre / Accroche (Tagline) :} \emph{1 ligne} sous le nom. Ex: \textit{Technicien Maintenance Industrielle -- Automatismes API -- Permis B -- Disponible immédiatement}. Pas de « Je recherche un poste... ».
                  \item \textbf{Formation :} Antéchronologique. Diplôme, Année, Établissement, Ville, \textbf{Option/Spécialité}, \textbf{Mention}, \textbf{Projet/Mémoire} (si pertinent).
                  \item \textbf{Expériences Professionnelles :} Antéchronologique. Dates (Mois/Année), \textbf{Intitulé exact du poste}, \textbf{Entreprise}, Ville. \textbf{3 à 5 puces} : Verbe d'action + Contexte + Résultat/Chiffre.
                  \item \textbf{Compétences Techniques (Hard Skills) :} Regroupées par familles (ex: Automatismes, Électricité, Mécanique, Informatique/CAO, Langues). Niveau : \emph{Maîtrise / Autonome / Notions}.
              \end{enumerate}
        \item \textbf{Rubriques RECOMMANDÉES (différenciantes) :}
              \begin{itemize}
                  \item \textbf{Projets / Réalisations :} Projet fin d'études, mini-projets, hackathons, maintenance personnelle (ex: « Rénovation complète installation électrique maison »).
                  \item \textbf{Certifications / Habilitations :} \cle{Habilitation électrique (BR, BC, B2V)}, CACES, SST (Sauveteur Secouriste du Travail), ISO 9001, Permis grue/chariot. \textbf{Très valorisés en technique.}
              \end{itemize}
        \item \textbf{Rubriques OPTIONNELLES (si place et pertinence) :}
              \begin{itemize}
                  \item \textbf{Langues :} Niveau CECRL (A1--C2) ou « Lu, écrit, parlé / Courant / Notions ». Ne pas mentir.
                  \item \textbf{Informatique / Bureautique / CAO :} Pack Office, SolidWorks, AutoCAD, Eagle, TIA Portal, Python/C basique.
                  \item \textbf{Centres d'intérêt / Vie associative :} Uniquement si \cle{soft skills} transférables (Capitaine = leadership, Trésorier = rigueur/gestion, Bénévolat = engagement). Éviter : « Lecture, Cinéma, Voyage » (trop vague).
              \end{itemize}
        \item \textbf{Ordre stratégique :} \textbf{Jeune diplômé} $\rightarrow$ Formation avant Expériences. \textbf{Expérimenté} $\rightarrow$ Expériences avant Formation.
        \item \textbf{Reflexe Bac :} Ne jamais inventer de rubrique farfelue. Respecter les standards du secteur (Industrie / BTP / Tertiaire).
    \end{enumerate}
\end{methode}

\begin{exoresolu}[Compléter un CV squelette pour une offre précise]
    \textbf{Énoncé :} Offre d'emploi : \emph{« Entreprise de BTP recherche Technicien Bureau d'Études Électricité (BT/BTS). Missions : Études schémas (AutoCAD/SEE Electrical), dimensionnement câbles, chiffrage, suivi chantier. Habilitation B2V exigée. Permis B. »}
    Complétez le squelette ci-dessous pour M. \textbf{Kamga Eric}, 24 ans, BT Électrotechnique 2022, Habilitation B2V (2023), 1 an stage BET, 6 mois chantier, AutoCAD/SEE Electrical, Excel. Permis B. Capitaine handball.

    \textit{Squelette :}
    \begin{verbatim}
    EN-TÊTE : [À remplir]
    TITRE : [À remplir]
    FORMATION : [À remplir]
    EXPÉRIENCES : [À remplir]
    COMPÉTENCES TECHNIQUES : [À remplir]
    CERTIFICATIONS : [À remplir]
    CENTRES D'INTÉRÊT : [À remplir]
    \end{verbatim}

    \tcblower
    \textbf{Solution détaillée :}

    \textbf{EN-TÊTE :}
    \textbf{KAMGA Eric} -- Douala, Akwa -- +237 6XX XX XX XX -- eric.kamga@email.com -- linkedin.com/in/eric-kamga -- \textbf{Permis B} -- 24 ans -- \textit{Photo pro}

    \textbf{TITRE :}
    \textit{Technicien Bureau d'Études Électricité -- BT Électrotechnique -- Habilitation B2V -- AutoCAD / SEE Electrical -- Permis B}

    \textbf{FORMATION (Antéchronologique) :}
    \begin{itemize}
        \item \textbf{2022 -- BT Électrotechnique} (Option Courants Forts) -- \textit{Lycée Technique de Douala} -- Mention Bien. \textit{Projet : Dimensionnement installation électrique immeuble R+3.}
        \item \textbf{2019 -- Probatoire F3 (Électrotechnique)} -- \textit{Lycée Technique de Bonabéri} -- Mention Assez Bien.
        \item \textbf{2017 -- BEPC} -- \textit{Collège Polyvalent de Bonamoussadi}.
    \end{itemize}

    \textbf{EXPÉRIENCES PROFESSIONNELLES (Antéchronologique) :}
    \begin{itemize}
        \item \textbf{09/2023 -- 03/2024 (6 mois) : Technicien BET Électricité (Stage long / Contrat pro)} \\
              \textit{ETUDElec Ingénierie, Douala} \\
              \textbullet\ Réalisation schémas unifilaires et de câblage sous \textbf{SEE Electrical} et \textbf{AutoCAD}. \\
              \textbullet\ Dimensionnement câbles et protections (norme NF C 15-100) pour 5 projets tertiaires. \\
              \textbullet\ Participation au chiffrage (métrage, consultation fournisseurs, devis). \\
              \textbullet\ Visites de chantier pour reprise de côtes et suivi conformité exécution.
        \item \textbf{03/2023 -- 08/2023 (6 mois) : Stagiaire Électricien de Chantier} \\
              \textit{Entreprise BTP « Bâtir Plus », Douala} \\
              \textbullet\ Tirage câbles, pose appareillage, raccordement tableaux BT. \\
              \textbullet\ Lecture plans d'exécution, report modifications (DOE). \\
              \textbullet\ Application stricte habilitation \textbf{B2V} (consignation, vérification absence tension).
    \end{itemize}

    \textbf{COMPÉTENCES TECHNIQUES :}
    \begin{itemize}
        \item \textbf{Logiciels BET :} \textbf{SEE Electrical} (Expert), \textbf{AutoCAD 2D} (Avancé), \textbf{Excel} (Avancé : TCD, RechercheV, macros bases).
        \item \textbf{Électricité BT :} Dimensionnement (NF C 15-100), Schémas unifilaires/multifilaires, Chiffrage/Devis, Lecture plans architectes.
        \item \textbf{Chantier :} Pose chemins de câbles, Raccordement TGBT, Mise en service, Respect normes sécurité.
        \item \textbf{Langues :} Français (Maternel), Anglais (B1 -- Lu, écrit, parlé technique).
    \end{itemize}

    \textbf{CERTIFICATIONS / HABILITATIONS :}
    \begin{itemize}
        \item \textbf{Habilitation Électrique B2V / BR / BC} (Recyclage 2023 -- Valide).
        \item \textbf{SST -- Sauveteur Secouriste du Travail} (2022 -- Valide).
        \item \textbf{Permis B} (2020).
    \end{itemize}

    \textbf{CENTRES D'INTÉRÊT :}
    \begin{itemize}
        \item \textbf{Handball (Capitaine équipe régionale)} : Leadership, gestion de groupe, stratégie, endurance.
        \item \textbf{Bricolage / Domotique personnelle} : Installation éclairage connecté, programmation scénarios (veille technologique pratique).
    \end{itemize}

    \textbf{Conclusion :} Chaque rubrique répond à un critère de l'offre (B2V, AutoCAD/SEE, Chiffrage, Chantier, Permis). L'ordre Formation $\rightarrow$ Expériences convient au jeune diplômé. Les centres d'intérêt prouvent des \cle{soft skills} utiles (capitaine = encadrement).
\end{exoresolu}

\courssub{Langue : les outils de liaison dans la phrase}

\begin{methode}[Fluidifier et professionnaliser la rédaction : connecteurs logiques et relatifs]
    \textbf{Objectif :} Éviter le style « télégraphique » (suite de phrases courtes) ou le style « pâté » (phrases trop longues). Utiliser les outils de liaison pour montrer la logique (cause, conséquence, but, opposition, chronologie) et alléger la syntaxe.
    \begin{enumerate}
        \item \textbf{Les adverbes de liaison (mots invariables) :} Placés en début de phrase ou après le sujet. Indiquent la relation logique entre deux phrases.
              \begin{center}
              \begin{tabularx}{\linewidth}{|l|X|l|}\hline
              \rowcolor{popblueL} \textbf{Relation} & \textbf{Adverbes / Locutions} & \textbf{Exemple CV} \\ \hline
              \textbf{Ajout / Énumération} & \cle{De plus}, \cle{En outre}, \cle{Par ailleurs}, \cle{Aussi}, \cle{Également} & \textit{J'ai géré le stock. \textbf{De plus}, j'ai formé 2 stagiaires.} \\ \hline
              \textbf{Conséquence / Résultat} & \cle{Par conséquent}, \cle{Aussi}, \cle{Donc}, \cle{En conséquence}, \cle{Dès lors} & \textit{Maintenance préventive rigoureuse. \textbf{Par conséquent}, 0 panne critique sur l'année.} \\ \hline
              \textbf{Opposition / Concession} & \cle{Cependant}, \cle{Néanmoins}, \cle{Pourtant}, \cle{En revanche}, \cle{Toutefois} & \textit{Stage court (1 mois). \textbf{Néanmoins}, acquisition solides bases API.} \\ \hline
              \textbf{Chronologie / Suite} & \cle{Ensuite}, \cle{Par la suite}, \cle{Ulérieurement}, \cle{Préalablement} & \textit{Diagnostic panne. \textbf{Ensuite}, remplacement composant défectueux.} \\ \hline
              \textbf{But} & \cle{Afin de}, \cle{Dans le but de}, \cle{Pour} (locutions) & \textit{Analyse schémas \textbf{afin de} localiser le défaut.} \\ \hline
              \end{tabularx}
              \end{center}
        \item \textbf{Les pronoms relatifs (substitution / enchaînement) :} Évitent la répétition du nom (antécédent). Obligatoires pour faire des phrases complexes fluides.
              \begin{itemize}
                  \item \cle{qui} (sujet) : \textit{Le technicien \textbf{qui} a réparé la machine est parti.}
                  \item \cle{que} (COD) : \textit{La machine \textbf{que} j'ai réparée fonctionne.} (Accord !)
                  \item \cle{dont} (objet prépositionnel \textit{de}) : \textit{L'entreprise \textbf{dont} je dépends...} / \textit{Le projet \textbf{dont} j'ai la charge...}
                  \item \cle{où} (lieu, temps) : \textit{L'usine \textbf{où} j'ai travaillé...} / \textit{L'année \textbf{où} j'ai obtenu mon BT...}
                  \item \cle{lequel, laquelle, lesquels, lesquelles} (après préposition \textit{à, pour, avec, sans, dans, sur...}) : \textit{L'outil \textbf{avec lequel} je travaille...}
              \end{itemize}
        \item \textbf{Application au CV (règle d'or) :} Dans les \textbf{puces d'expériences}, commencer par un \textbf{verbe d'action à l'infinitif} (ou participe passé si phrase nominale) + \textbf{compléments} + \textbf{connecteur de résultat/contexte si besoin}.
              \begin{itemize}
                  \item \textit{Mauvais :} « J'ai fait la maintenance. C'était bien. J'ai utilisé une valise. »
                  \item \textit{Bien (infinitif + connecteurs) :} \textbf{Assurer} la maintenance préventive/curative \textbf{afin de} garantir la disponibilité des équipements ; \textbf{diagnostiquer} les pannes \textbf{grâce à} l'outil valise multimarque \textbf{et} analyser les défauts \textbf{dont} l'origine est électronique.
              \end{itemize}
    \end{enumerate}
\end{methode}

\begin{exoresolu}[Réécriture professionnelle de puces d'expérience]
    \textbf{Énoncé :} Réécrivez les 3 descriptions d'expérience suivantes en style « CV professionnel » (puces, verbes d'action à l'infinitif, connecteurs logiques, pronoms relatifs pour fluidifier). Intégrez un résultat chiffré si possible.

    1. \textit{« J'ai fait un stage au garage Auto Modern. J'ai fait des vidanges. J'ai aussi changé des freins. J'ai utilisé la valise diag. C'était bien. »}
    2. \textit{« J'ai travaillé à la Socapalm. J'ai porté des charges. C'était dur. Mais j'ai appris la sécurité. »}
    3. \textit{« J'ai fait de l'électricité chez mon père. J'ai tiré des câbles. J'ai posé des interrupteurs. »}

    \tcblower
    \textbf{Solution détaillée :}

    \textbf{Principes appliqués :} Verbe infinitif (standard CV francophone) ou Participe passé (style « réalisé »). Suppression « Je ». Intégration connecteurs (\textit{afin de, grâce à, dont, par conséquent}). Quantification.

    \textbf{1. Garage Auto Modern (Stage Maintenance) :}
    \begin{itemize}
        \item \textbf{Réaliser} les opérations de maintenance préventive (vidanges, filtres, niveaux) \textbf{et} curative (freins, distribution, suspension) \textbf{conformément aux} carnets constructeur.
        \item \textbf{Diagnostiquer} les pannes électroniques \textbf{grâce à} l'outil de diagnostic multimarque (valise) \textbf{afin d'identifier} les codes défauts \textbf{dont} l'origine est capteur/actionneur.
        \item \textbf{Participer} à la réception client et à l'explication des interventions \textbf{afin de} garantir la satisfaction client (\textbf{+15\%} avis positifs atelier pendant le stage).
    \end{itemize}

    \textbf{2. SOCAPALM (Agent Maintenance/Manutention) :}
    \begin{itemize}
        \item \textbf{Assurer} la manutention de charges lourdes (palettes, fûts) \textbf{en respectant strictement} les règles de gestuelle/posture \textbf{et} le port des EPI.
        \item \textbf{Effectuer} la surveillance de premier niveau des convoyeurs (graissage, tension courroies, alignement) \textbf{afin de} prévenir les arrêts ligne.
        \item \textbf{Signaler} toute anomalie (bruit, vibration, fuite) \textbf{auprès de} l'équipe maintenance \textbf{par conséquent} réduction du temps d'arrêt moyen (MTTR) \textbf{de 20\%} sur le poste.
    \end{itemize}

    \textbf{3. Électricité Bâtiment (Stage/Expérience personnelle) :}
    \begin{itemize}
        \item \textbf{Exécuter} l'installation électrique courant fort/faible en rénovation \textbf{selon} la norme NF C 15-100 : tirage câbles (gaines, tubes), pose appareillage (interrupteurs, prises, tableaux divisionnaires).
        \item \textbf{Lire} et \textbf{interpréter} les schémas unifilaires \textbf{dont} dépend le câblage des circuits éclairage/force motrice.
        \item \textbf{Effectuer} les tests de continuité, isolement et mise en service \textbf{afin de} valider la conformité \textbf{avant} passage Consuel.
    \end{itemize}

    \textbf{Analyse linguistique (Justification Bac) :}
    \begin{itemize}
        \item \textbf{Infinitif :} Standard pour les puces CV (impersonnel, orienté action). \textit{Réaliser, Diagnostiquer, Assurer, Effectuer, Exécuter, Lire.}
        \item \textbf{Connecteurs de but :} \textit{afin de, pour, dans le but de} (introduisent l'objectif professionnel).
        \item \textbf{Connecteurs de moyen :} \textit{grâce à, au moyen de, en utilisant}.
        \item \textbf{Pronoms relatifs :} \textit{dont} (origine \textbf{de} laquelle / \textbf{duquel}) $\rightarrow$ « codes défauts \textbf{dont} l'origine... » ; \textit{dont} (dépendre \textbf{de}) $\rightarrow$ « schémas \textbf{dont} dépend le câblage ».
        \item \textbf{Connecteur conséquence :} \textit{par conséquent} (marque le résultat prouvé).
    \end{itemize}

    \textbf{Conclusion :} La maîtrise des outils de liaison transforme une liste de tâches banales en une démonstration de \cle{compétences transférables} (diagnostic, norme, sécurité, résultat).
\end{exoresolu}

\courssub{Rédaction du CV : travail préparatoire}

\begin{methode}[La méthode « Inventaire -- Sélection -- Adaptation » (ISA) avant d'écrire]
    \textbf{Objectif :} Ne pas écrire son CV « au feeling ». Construire une \cle{banque de données personnelle} réutilisable et cibler l'offre.
    \begin{enumerate}
        \item \textbf{Étape 1 : INVENTAIRE (Le « Dépôt brut ») -- Faire le vide complet.}
              \begin{itemize}
                  \item Lister \textbf{Tout} : Diplômes (dates, mentions, options, projets, notes clés), Expériences (stages, jobs, bénévolat, projets persos) : Dates, Structure, Missions détaillées, Résultats, Outils, Équipe, Contexte.
                  \item Lister \textbf{Toutes} compétences : Techniques (logiciels, machines, normes, langages, habilitations), Langues (niveau CECRL, scores), Transverses (Permis, Premiers secours, Bureautique, Gestion de projet).
                  \item Lister \textbf{Centres d'intérêt} : Rôle, responsabilité, fréquence, niveau (compétition, loisir), compétences associées (trésorier = Excel/gestion ; capitaine = leadership).
              \end{itemize}
              \textit{Support :} Fichier Excel / Notion / Word « Master CV » (3--4 pages), mis à jour chaque semestre.
        \item \textbf{Étape 2 : ANALYSE DE L'OFFRE (Le « Filtre ») -- Décoder l'ADN du poste.}
              \begin{itemize}
                  \item Surligner : \textbf{Mots-clés techniques} (logiciels, machines, normes, langages), \textbf{Soft skills} demandés (autonomie, rigueur, travail équipe, relation client), \textbf{Diplômes/Niveau}, \textbf{Expérience requise}, \textbf{Permis/Déplacements}, \textbf{Habilitations}.
                  \item Identifier le \textbf{vocabulaire de l'entreprise} (ex: « Maintenance préventive » vs « Maintenance planifiée » ; « Chiffrage » vs « Devis »).
              \end{itemize}
        \item \textbf{Étape 3 : SÉLECTION \& ADAPTATION (Le « Sur-mesure ») -- Matching Offre/Profil.}
              \begin{itemize}
                  \item \textbf{Copier-Coller ciblé :} Depuis le Master CV, ne garder que les éléments qui répondent aux mots-clés de l'offre.
                  \item \textbf{Reformuler :} Utiliser \textbf{le vocabulaire de l'offre} dans vos puces (ATS + Recruteur humain).
                  \item \textbf{Ordonner :} Mettre en \textbf{premier} (rubrique + premières puces) ce qui est \textbf{le plus pertinent} pour \emph{ce} poste.
                  \item \textbf{Quantifier :} Ajouter des chiffres (budget, effectif, \%, délai, volume, taux de réussite) là où c'est possible.
              \end{itemize}
        \item \textbf{Étape 4 : RELECTURE CROISÉE (Contrôle Qualité) :}
              \begin{itemize}
                  \item Orthographe / Grammaire (Antidote, LanguageTool, relecture à voix haute, tiers).
                  \item Cohérence dates (trous ? chevauchements ?).
                  \item Cohérence Titre / Accroche $\leftrightarrow$ Offre.
                  \item Lisibilité 30 sec (test : plier la feuille en 4, le haut de page suffit-il ?).
                  \item Format PDF, nom de fichier pro.
              \end{itemize}
    \end{enumerate}
\end{methode}

\begin{exoresolu}[Application ISA : Préparation pour une offre « Technicien SAV Itinérant »]
    \textbf{Énoncé :} Offre : \emph{« Société de distribution matériel médical recherche Technicien SAV Itinérant (H/F). Bac+2/3 Maintenance Biomédicale ou Électrotechnique. Missions : Dépannage sur site clients (hôpitaux/cliniques), maintenance préventive, relation client, reporting numérique. Permis B obligatoire. Déplacements régionaux 80\%. Anglais technique lu. »}
    Candidat : \textbf{Mme Mvondo Alice}, 25 ans, BTS Maintenance Biomédicale (2022), 1 an CDD Hôpital Central (maintenance interne), 6 mois Stage Usine Fabricant (réparation cartes électroniques), Habilitation B2V, Anglais B2, Permis B, Voiture perso. Compétences : Appareils anesthésie, respirateurs, scopes, sondes échographie, GMAO (Maintenance Assistant), Word/Excel.

    \tcblower
    \textbf{Solution détaillée :}

    \textbf{ÉTAPE 1 : INVENTAIRE (Extrait Master CV Alice) -- \textit{Non exhaustif, focus pertinent}}
    \begin{itemize}
        \item \textbf{Diplôme :} BTS Maintenance Biomédicale (2022) -- Projet : « Développement checklist maintenance préventive respirateurs ».
        \item \textbf{Expé 1 (2022-2023) :} Technicien Maintenance Biomédicale (CDD 1 an) -- Hôpital Central Yaoundé. Parc : 150 lits. Missions : MP/MC respirateurs, monitors, scopes, tables OP. GMAO : Maintenance Assistant. Formation utilisateurs (infirmiers). Astreintes WE.
        \item \textbf{Expé 2 (2021) :} Stagiaire SAV Usine -- MedTech Industries (Fabricant). Réparation niveau 3 cartes électroniques (soudure CMS, test banc). Lecture schémas composants. Anglais technique doc constructeur.
        \item \textbf{Compétences Tech :} Respirateurs (Hamilton, Dräger), Anesthésie (Datex, Fabius), Échographie (GE, Sonosite), Endoscopie (Olympus, Pentax), Électrochirurgie. Soudure CMS. Normes IEC 60601. GMAO Maintenance Assistant. Pack Office.
        \item \textbf{Certifs :} Habilitation B2V (2023), SST (2022), Formation constructeur Dräger (2023).
        \item \textbf{Langues :} Anglais B2 (TOEIC 780) -- Lu/Écrit/Parlé technique. Allemand A2.
        \item \textbf{Permis/Voiture :} Permis B (2019), Véhicule personnel.
        \item \textbf{Soft/Intérêts :} Trésorière Amicale Anciens BTS (gestion budget 500k FCFA), Course à pied (semi-marathon), Piano.
    \end{itemize}

    \textbf{ÉTAPE 2 : ANALYSE OFFRE -- Mots-clés cibles :}
    \cle{Itinérant}, \cle{Dépannage sur site}, \cle{Clients hôpitaux/cliniques}, \cle{Maintenance préventive}, \cle{Relation client}, \cle{Reporting numérique/GMAO}, \cle{Permis B}, \cle{Déplacements 80\%}, \cle{Anglais technique}, \cle{Biomédicale/Électrotechnique}.

    \textbf{ÉTAPE 3 : SÉLECTION \& ADAPTATION -- CV CIBLÉ « TECHNICIEN SAV ITINÉRANT »}

    \textbf{EN-TÊTE :} MVONDO Alice -- Yaoundé, Mvan -- +237 6XX XX XX XX -- alice.mvondo@email.com -- linkedin.com/in/alice-mvondo-biomed -- \textbf{Permis B + Véhicule} -- 25 ans -- Disponibilité immédiate -- \textit{Photo pro}

    \textbf{TITRE :} \textit{Technicienne SAV Itinérante Biomédicale -- BTS Maintenance Biomédicale -- Expérience Hôpital \& Usine -- Anglais Technique B2 -- Permis B + Véhicule -- Mobile 100\%}

    \textbf{EXPÉRIENCES PROFESSIONNELLES (Antéchronologique -- \textbf{Expé avant Formation} car 1 an expé significative) :}
    \begin{itemize}
        \item \textbf{10/2022 -- 10/2023 : Technicienne Maintenance Biomédicale (CDD)} \\
              \textit{Hôpital Central, Yaoundé (Parc 150 lits + Consultations externes)} \\
              \textbullet\ \textbf{Assurer} la maintenance préventive (planification GMAO \textbf{Maintenance Assistant}) et curative des équipements de réanimation/bloc : \textbf{Respirateurs (Hamilton, Dräger)}, Monitors, \textbf{Anesthésie (Datex, Fabius)}, Électrochirurgie \textbf{conformément aux} normes IEC 60601.
              \textbullet\ \textbf{Intervenir} en urgence (astreintes WE) \textbf{afin de} garantir la continuité des soins : diagnostic panne, remplacement modules, tests fonctionnels, remise en service \textbf{dans un délai < 2h} (objectif atteint 95\% cas).
              \textbullet\ \textbf{Gérer} la relation client interne (médecins, cadres infirmiers) : explication pannes, délais, formation utilisateurs aux checklists quotidiennes \textbf{afin de} réduire les fausses alertes (\textbf{-30\%} appels non fondés).
              \textbullet\ \textbf{Rédiger} les rapports d'intervention numériques (GMAO) et le reporting mensuel d'activité (KPI : taux dispo, MTTR, coûts pièces).
        \item \textbf{03/2021 -- 09/2021 : Stagiaire Technicienne SAV Niveau 3 (Usine Fabricant)} \\
              \textit{MedTech Industries, Douala (Fabricant matériel anesthésie/réanimation)} \\
              \textbullet\ \textbf{Réparer} en atelier les cartes électroniques défectueuses (niveau 3) : diagnostic composant, \textbf{soudure CMS}, test sur banc de test constructeur \textbf{grâce à} la lecture de schémas électroniques complexes \textbf{en anglais technique}.
              \textbullet\ \textbf{Participer} à la qualification de nouveaux sous-ensembles (protocoles test, documentation).
              \textbullet\ \textbf{Capitaliser} les modes opératoires de réparation courants \textbf{pour} enrichir la base de connaissances SAV.
    \end{itemize}

    \textbf{FORMATION :}
    \begin{itemize}
        \item \textbf{2022 -- BTS Maintenance Biomédicale} -- \textit{ISTA Yaoundé} -- Mention Bien. \textit{Mémoire : « Optimisation plan maintenance préventive parc respirateurs »}.
        \item \textbf{2019 -- Baccalauréat F3 (Électrotechnique)} -- \textit{Lycée Technique Mballa II} -- Mention AB.
    \end{itemize}

    \textbf{COMPÉTENCES TECHNIQUES (Mots-clés offre en \textbf{gras}) :}
    \begin{itemize}
        \item \textbf{Équipements critiques :} \textbf{Respirateurs} (Hamilton G5/S1, Dräger Evita), \textbf{Anesthésie} (Dräger Fabius, GE Datex), \textbf{Échographie} (GE Voluson, Sonosite), \textbf{Endoscopie} (Olympus CV-190), Monitors multiparamétriques.
        \item \textbf{Intervention :} \textbf{Dépannage sur site} (niveau 1/2), \textbf{Maintenance préventive} (planification, checklists), \textbf{Maintenance curative}, \textbf{Étallonnage/Tests fonctionnels} (analyseurs gaz, débitmètres).
        \item \textbf{Électronique/Élec :} Lecture schémas composants, \textbf{Soudure CMS}, Habilitation \textbf{B2V}, Normes \textbf{IEC 60601}.
        \item \textbf{Outils numériques :} \textbf{GMAO Maintenance Assistant} (Expert), Pack Office (Excel TCD, Word reporting), TeamViewer (téléassistance).
        \item \textbf{Langues :} \textbf{Anglais Technique B2} (TOEIC 780 -- Lu doc constructeur, rapport technique, oral client), Allemand A2.
    \end{itemize}

    \textbf{CERTIFICATIONS / ATOUTS MOBILITÉ :}
    \begin{itemize}
        \item Habilitation Électrique \textbf{B2V} (Valide), SST (Valide), Formation Constructeur \textbf{Dräger} (2023).
        \item \textbf{Permis B (2019) + Véhicule personnel} -- \textbf{Disponible pour déplacements régionaux 100\%}.
    \end{itemize}

    \textbf{CENTRES D'INTÉRÊT (Soft skills ciblés : Autonomie, Gestion, Endurance) :}
    \begin{itemize}
        \item \textbf{Trésorière Association Anciens BTS (2022--2024)} : Gestion budget 500k FCFA, organisation événements, relation adhérents (Rigueur, Gestion, Relationnel).
        \item \textbf{Course à pied (Semi-marathon 2023)} : Endurance, discipline, gestion effort (Adapté aux déplacements/astreintes).
    \end{itemize}

    \textbf{Conclusion :} La méthode ISA permet de passer d'un CV « généraliste » à un CV « miroir de l'offre ». L'expérience hôpital (client final, itinérance interne, relationnel, GMAO, astreintes) est mise en avant avant l'expérience usine. Le titre, la mobilité (voiture), l'anglais technique et la GMAO sont en « une ».
\end{exoresolu}

\courssub{Activités d'intégration : rédiger un CV complet}

\begin{methode}[Rédiger le CV final ciblé : de la page blanche au PDF prêt à envoyer]
    \textbf{Objectif :} Synthétiser toutes les étapes (Réception, Structure, Langue, Préparation) pour produire le livrable final évalué au Bac.
    \begin{enumerate}
        \item \textbf{Choisir le bon modèle (Template) :} Sobre, une colonne (préféré ATS) ou deux colonnes (barre latérale compétences/coordonnées) si design maîtrisé. \textbf{Éviter :} Canva trop graphique (colonnes complexes, graphiques en camembert pour compétences, photos de fond) $\rightarrow$ illisible par les ATS.
        \item \textbf{Rédiger l'Accroche (Titre + 1 phrase) :} \textit{[Métier visé] -- [Diplôme/Niveau] -- [Spécialité/Expertise clé] -- [Atout différenciant : Permis, Langue, Habilitation, Mobilité]}. Ex: \textit{Technicien Maintenance Industrielle -- BT MSMA -- Automatismes Siemens/Autocad -- Habilitation B2V -- Permis B -- Mobile National}.
        \item \textbf{Remplir les rubriques dans l'ordre stratégique (voir Méthode 3) :} Formation ou Expériences en premier. Pour chaque entrée : \textbf{Date -- Intitulé Gras -- Structure Gras -- Ville}. Puces : \textbf{Verbe infinitif} + \textbf{Action précise} + \textbf{Contexte/Outils} + \textbf{Résultat/Chiffre} + \textbf{Connecteur logique}.
        \item \textbf{Soigner la rubrique Compétences :} Grouper par familles. Indiquer le niveau (\textit{Expert / Autonome / Notions}). Mettre en \textbf{gras} les mots-clés de l'offre.
        \item \textbf{Relire « comme un recruteur » (Check-list finale 5 min) :}
              \begin{enumerate}
                  \item \textbf{Orthographe/Zéro faute :} Noms propres (entreprises, logiciels, diplômes), accords (participes passés, relatifs \textit{dont/que}).
                  \item \textbf{Cohérence chronologique :} Pas de trous inexpliqués > 6 mois (si trou : « Recherche active », « Formation », « Projet personnel », « Service civique »).
                  \item \textbf{Pertinence :} Chaque ligne sert-elle le poste visé ? Supprimer le superflu (colonie de vacances 2010, note de BEPC).
                  \item \textbf{Lisibilité :} Police unique, tailles cohérentes, alignements parfaits (tabulations, pas espaces), puces alignées, marges régulières.
                  \item \textbf{Identité numérique :} Email pro, LinkedIn à jour (cohérence CV/LinkedIn), pas de photo Facebook en lien.
              \end{enumerate}
        \item \textbf{Exporter :} \textbf{PDF} (Imprimer $\rightarrow$ Enregistrer sous PDF / Exporter). Nom : \texttt{Nom\_Prenom\_IntitulePoste\_Date.pdf}. Tester l'ouverture sur téléphone.
    \end{enumerate}
\end{methode}

\begin{exoresolu}[Production intégrée : CV complet pour « Technicien de Maintenance H/F » -- Sujet type Bac]
    \textbf{Énoncé :} « Vous êtes titulaire d'un Baccalauréat Technique (Série F4 Électrotechnique) obtenu en 2023 au Lycée Technique de Douala. Vous avez effectué un stage de 3 mois (Janv--Mars 2023) chez ENEO (Distribution) : maintenance préventive postes HTA/BT, relevé de charges, thermographie. Un job d'été Juil--Août 2022 chez CAMWATER (Plomberie/Élec) : réparation fuites, câblage pompes. Vous possédez l'Habilitation B2V, le Permis B, l'Anglais niveau B1. Vous maîtrisez AutoCAD Élec, Excel, Word. Vous êtes capitaine de l'équipe de basket du lycée. Vous postulez à l'offre : \emph{« ENEO recrute Techniciens Maintenance Réseau HTA/BT (Bac+2/3). Missions : Maintenance préventive/curative postes transformation, dépannage réseau, reporting GMAO. Permis B. Habilitation B2V. Déplacements. »} Rédigez votre CV complet sur une page. »

    \tcblower
    \textbf{Solution détaillée : CV FINAL PRÊT À ENVOYER (Représentation textuelle structurée)}

    \begin{center}
    \begin{tabularx}{\linewidth}{|X|}\hline
    \rowcolor{popblueL} \multicolumn{1}{|c|}{\textbf{\LARGE NJIKI Arnaud}} \\ \hline
    \rowcolor{popblue!10} \multicolumn{1}{|c|}{\textit{Technicien Maintenance Réseau HTA/BT -- BT Électrotechnique -- Habilitation B2V -- Permis B -- Mobile}} \\ \hline
    \textbf{Contact :} Douala, Bonapriso -- +237 6XX XX XX XX -- arnaud.njiki@email.com -- linkedin.com/in/arnaud-njiki -- 22 ans -- \textit{Photo pro} \\ \hline
    \rowcolor{popblueL} \multicolumn{1}{|c|}{\textbf{EXPÉRIENCES PROFESSIONNELLES}} \\ \hline
    \textbf{01/2023 -- 03/2023 (3 mois) : Stagiaire Technicien Maintenance Réseau HTA/BT} \\
    \textit{ENEO -- Direction Distribution Douala} \\
    \textbullet\ \textbf{Effectuer} la maintenance préventive planifiée (GMAO) sur postes de transformation HTA/BT (63/20 kV \& 20/0,4 kV) : contrôle serrage, nettoyage isolateurs, vérification relais protection, \textbf{thermographie} infrarouge connexions. \\
    \textbullet\ \textbf{Réaliser} les relevés de charges (ampèremétrique) et analyses d'huile (prélèvement, étiquetage, envoi labo) \textbf{afin de} détecter les surcharges et défauts naissants. \\
    \textbullet\ \textbf{Participer} aux interventions curatives sur réseau BT (recherche défauts, raccordement, mise en service) \textbf{en respectant} les procédures de consignation \textbf{Habilitation B2V}. \\
    \textbullet\ \textbf{Renseigner} les bons d'intervention et rapports quotidiens sur tablette (GMAO \textbf{Maximo}) \textbf{pour} assurer la traçabilité et le reporting hiérarchique. \\ \hline
    \textbf{07/2022 -- 08/2022 (2 mois) : Agent Technique Électricité/Plomberie (Job d'été)} \\
    \textit{CAMWATER -- Unité de Production Eau, Douala} \\
    \textbullet\ \textbf{Assister} l'équipe maintenance sur les pompages : réparation fuites réseaux (soudure PEHD, raccords mécanique), \textbf{câblage et raccordement moteurs pompes} (étoile/triangle, variateurs). \\
    \textbullet\ \textbf{Appliquer} les consignes sécurité (EPI, balisage, autorisation travail) \textbf{sur sites industriels} sensibles (traitement eau). \\ \hline
    \rowcolor{popblueL} \multicolumn{1}{|c|}{\textbf{FORMATION}} \\ \hline
    \textbf{2023 -- Baccalauréat Technique Série F4 (Électrotechnique)} -- \textit{Lycée Technique de Douala} -- \textbf{Mention Assez Bien}. \\
    \textit{Projet fin d'études : « Étude et dimensionnement poste transformation 20/0,4 kV pour zone industrielle » (AutoCAD Élec, Caneco BT).} \\
    \textbf{2020 -- Probatoire Série F4} -- \textit{Lycée Technique de Douala} -- Mention AB. \\
    \textbf{2018 -- BEPC} -- \textit{Collège Polyvalent Bonabéri}. \\ \hline
    \rowcolor{popblueL} \multicolumn{1}{|c|}{\textbf{COMPÉTENCES TECHNIQUES}} \\ \hline
    \textbf{Réseau HTA/BT :} Maintenance préventive/curative postes transformation, \textbf{Thermographie}, Relais protection (notions), Relevés charges, Consignation B2V. \\
    \textbf{Électricité Industrielle :} Câblage armoires, Moteurs (Démarrage, Variateurs), Lecture schémas unifilaires/multifilaires, Normes NF C 13-100/15-100. \\
    \textbf{Logiciels / CAO / GMAO :} \textbf{AutoCAD Electrical} (Avancé), \textbf{Caneco BT} (Dimensionnement), \textbf{Maximo/GMAO} (Utilisateur), Excel (TCD, RechercheV), Word. \\
    \textbf{Habilitations / Sécurité :} \textbf{Habilitation Électrique B2V / BR / BC} (Valide), SST (Sauveteur Secouriste Travail), Port EPI, Consignation. \\
    \textbf{Langues :} Français (Maternel), \textbf{Anglais B1} (Lu technique, Écrit courant, Parlé opérationnel -- TOEIC 650). \\ \hline
    \rowcolor{popblueL} \multicolumn{1}{|c|}{\textbf{CENTRES D'INTÉRÊT}} \\ \hline
    \textbf{Basket-ball (Capitaine Équipe Lycée Technique Douala -- Championnat Régional)} : Leadership, cohésion d'équipe, gestion stress, discipline, endurance. \\
    \textbf{Bricolage / Domotique (Projet personnel : Automatisation portail + éclairage extérieur)} : Veille technologique, mise en pratique compétences élec/programmation. \\ \hline
    \end{tabularx}
    \end{center}

    \textbf{Analyse de conformité (Grille d'évaluation Bac) :}
    \begin{itemize}
        \item \textbf{Structure externe :} Une page, police unique (simulée), titres colorés, puces alignées, gras sur mots-clés, espacement aéré. \checkmark
        \item \textbf{Structure interne :} Titre accrocheur (métier + atouts), Expériences avant Formation (stage long pertinent), Rubriques standardisées, Compétences groupées par familles. \checkmark
        \item \textbf{Langue :} Verbes infinitifs (\textit{Effectuer, Réaliser, Participer, Renseigner, Assister, Appliquer}), Connecteurs (\textit{afin de, en respectant, pour, sur}), Relatifs implicites (fluidité). \checkmark
        \item \textbf{Adaptation offre :} Mots-clés offre repris (\textit{HTA/BT, Thermographie, Relevés charges, GMAO/Maximo, B2V, Permis B, Déplacements}). Stage ENEO mis en \textbf{premier} et détaillé. \checkmark
        \item \textbf{Préparation :} Inventaire fait (job d'été gardé car électricité/pompes pertinent), sélection opérée (pas de collège détaillé), quantification (63/20kV, 20/0,4kV, 3 mois, 2 mois). \checkmark
    \end{itemize}

    \textbf{Conclusion :} Ce CV est \textbf{« ATS-friendly »} (mots-clés, structure simple), \textbf{lisible en 30 sec} (titre, stage ENEO, B2V, Permis B sautent aux yeux) et \textbf{preuve de professionnalisme} (vocabulaire technique précis, résultats chiffrés, soft skills prouvés).
\end{exoresolu}

\begin{attention}[Les 5 erreurs qui coûtent des points au Bac (CV)]
    \begin{enumerate}
        \item \textbf{Erreur de structure : Inversion Formation/Expérience non justifiée.}
              \textit{Piège :} Mettre la Formation en premier alors qu'on a 1 an d'expérience pertinente (ou inversement pour un jeune diplômé sans stage significatif). \textbf{Règle :} Le plus pertinent pour \emph{ce} poste en premier.
        \item \textbf{Erreur de langue : Style « télégraphique » ou « Je » / Phrases complètes sans connecteurs.}
              \textit{Piège :} « J'ai fait la maintenance. J'ai utilisé la valise. C'était bien. » ou « Maintenance. Valise. Bien. » \textbf{Exigence Bac :} Puces \textbf{Verbe infinitif} + \textbf{Connecteurs logiques} (\textit{afin de, grâce à, par conséquent, dont}) + \textbf{Précision technique}.
        \item \textbf{Erreur de contenu : CV « Passe-partout » non adapté à l'offre.}
              \textit{Piège :} Envoyer le même CV pour « Maintenance Industrielle » et « Bureau d'Études Élec » sans changer l'ordre des compétences, le titre, les mots-clés. \textbf{Sanction :} Hors sujet / Inadéquation profil/poste (0 à 2/5 en adéquation).
        \item \textbf{Erreur de forme : Mise en page illisible / Faute d'orthographe sur mot-clé.}
              \textit{Pièges :} Police fantaisie, photo de mauvaise qualité, marges irrégulières, puces décalées, \textbf{faute sur le nom de l'entreprise}, du logiciel (ex: « AutoCad », « Seeelectrical »), du diplôme (ex: « Baccalauréat Technicien »). \textbf{Zéro tolérance} sur les noms propres techniques.
        \item \textbf{Erreur de stratégie : Mensonge / Surévaluation / Oubli mobilité.}
              \textit{Pièges :} « Anglais Courant » (niveau A2 réel) $\rightarrow$ Échec test oral. « Maîtrise API Siemens » (vu 1h en cours) $\rightarrow$ Échec technique. Oublier \textbf{Permis B + Voiture} pour poste itinérant $\rightarrow$ Élimination administrative. \textbf{Honnêteté + Précision du niveau} (Expert/Autonome/Notions).
    \end{enumerate}
\end{attention}

\newpage

\coursec{Exercices}

\coursec{Produire des écrits utilitaires : le CV}

\courssub{Découverte et réception du CV}

\begin{definition}[Curriculum Vitae (CV)]
Document synthétique, normé et actualisé qui présente le parcours professionnel, la formation, les compétences et le profil d'un candidat dans le but d'obtenir un entretien d'embauche. C'est un \textbf{texte fonctionnel} à visée persuasive.
\end{definition}

\begin{aretenir}[Les trois grands types de CV]
\begin{itemize}
\item \textbf{Chronologique inverse (anti-chronologique) :} Expériences de la plus récente à la plus ancienne. \textit{Standard le plus utilisé, valorise la progression.}
\item \textbf{Fonctionnel (par compétences) :} Regroupe les expériences par compétences (gestion, technique, relationnel). \textit{Utile pour reconversion, trous dans le parcours, jeunes diplômés.}
\item \textbf{Mixte (ou combiné) :} Profil / Compétences clés en haut, puis parcours chronologique. \textit{Souvent le plus efficace pour profils expérimentés.}
\end{itemize}
\end{aretenir}

\begin{exemplebox}[Exemple d'en-tête standard camerounais]
\begin{center}
\begin{tabularx}{\linewidth}{|l|X|}\hline
\textbf{Nom Prénom} & \textbf{Intitulé du poste visé} \\
\hline
Adresse (Ville, Quartier) & Tél : +237 6XX XX XX XX \\
Email : prenom.nom@domaine.pro & LinkedIn : linkedin.com/in/prenomnom \\
\hline
\end{tabularx}
\end{center}
\textit{Photo d'identité professionnelle (3,5 x 4,5 cm, fond neutre) placée en haut à droite ou gauche.}
\end{exemplebox}

\courssub{Structure externe du CV}

\begin{methode}[Règles de mise en page (Standard 1 page A4)]
\begin{enumerate}
\item \textbf{Police :} Sobre, lisible (Calibri, Arial, Helvetica, Roboto) — Taille 10 à 11 pt (interlignage 1,15).
\item \textbf{Marges :} 1,5 à 2 cm de chaque côté.
\item \textbf{Mise en forme :} Titres de rubriques en \textbf{gras} (taille 12-14), éventuellement \textsc{petites capitales} ou couleur unique (bleu marine, anthracite). Puces \textbullet\ ou tirets pour les items.
\item \textbf{Ordre des rubriques :} Adapté au profil (Jeune diplômé : Formation avant Expérience ; Expérimenté : Expérience avant Formation).
\item \textbf{Fichier :} Export en \textbf{PDF} (nom : \texttt{Nom\_Prenom\_PosteVisé.pdf}).
\end{enumerate}
\end{methode}

\begin{attention}[Pièges à éviter]
\begin{itemize}
\item CV sur plusieurs pages (sauf cadre dirigeant >15 ans exp.).
\item Polices fantaisistes, couleurs multiples, tableaux complexes (illisibles par les ATS\footnote{Applicant Tracking Systems : logiciels de tri automatique des CV.}).
\item Informations personnelles sensibles : CNI, groupe sanguin, situation matrimoniale détaillée, enfants, religion, photo non professionnelle.
\item Fautes d'orthographe, de grammaire, de syntaxe (éliminatoires).
\end{itemize}
\end{attention}

\courssub{Structure interne du CV : les rubriques}

\begin{definition}[Rubriques obligatoires]
\begin{enumerate}
\item \textbf{En-tête :} Identité, contacts, lien LinkedIn, photo.
\item \textbf{Objectif professionnel (ou Profil) :} 2-3 lignes : Poste visé + Valeur ajoutée principale + Disponibilité.
\item \textbf{Formation :} Diplômes (du plus récent au plus ancien), établissement, année, mention, option, mémoire/projet pertinent.
\item \textbf{Expériences professionnelles / Stages :} Dates (Mois/Année), Intitulé exact, Structure, Localisation, \textbf{3 à 5 réalisations chiffrées} (verbes d'action).
\item \textbf{Compétences techniques (Hard skills) :} Logiciels, langues (niveau CECRL + score), permis, certifications, techniques métier.
\item \textbf{Compétences transverses (Soft skills) :} 3 à 5 qualités prouvées par un contexte (ex: « Gestion du stress : coordination chantier 15 pers. »).
\end{enumerate}
\end{definition}

\begin{definition}[Rubriques facultatives (selon pertinence)]
Centres d'intérêt (associatif, sport collectif, bénévolat), Références (Nom, Fonction, Tél, Email — avec accord), Projets personnels, Publications, Veille technologique.
\end{definition}

\begin{methode}[Rédiger une expérience : la méthode STAR adaptée]
\begin{enumerate}
\item \textbf{Situation / Contexte :} Entreprise, service, effectif, budget.
\item \textbf{Tâche / Mission :} Verbe d'action à l'infinitif ou participe passé (ex: \textit{Superviser, Concevoir, Optimiser, Négocier}).
\item \textbf{Action / Moyens :} Outils, méthodes, équipe, fréquence.
\item \textbf{Résultat / Indicateur :} \textbf{Chiffres clés} (%, FCFA, délais, volume, taux de satisfaction).
\end{enumerate}
\textit{Exemple :} « \textit{Piloté la maintenance préventive/curative (Action) d'un parc de 50 machines (Contexte) via GMAO (Moyen) $\rightarrow$ Disponibilité 95\% / Coûts -15\% (Résultat).} »
\end{methode}

\courssub{Langue : les outils de liaison dans la phrase}

\begin{definition}[Outils de liaison (Connecteurs logiques)]
Mots invariables ou pronoms qui assurent la \textbf{cohérence textuelle} en marquant la relation logique (cause, conséquence, opposition, concession, addition, illustration, temps) entre deux énoncés (phrases ou propositions).
\end{definition}

\begin{propriete}[Classification des principaux outils de liaison]
\begin{center}
\begin{tabularx}{\linewidth}{|l|X|X|}\hline
\rowcolor{popblueL} \textbf{Relation logique} & \textbf{Adverbes / Locutions adverbiales} & \textbf{Conjonctions / Pronoms relatifs} \\ \hline
\textbf{Addition / Renforcement} & de plus, en outre, de surcroît, par ailleurs, également & et, non seulement... mais aussi \\
\textbf{Conséquence / Conclusion} & par conséquent, dès lors, c'est pourquoi, en définitive, donc & car, parce que, puisque, si bien que \\
\textbf{Opposition / Concession} & toutefois, cependant, néanmoins, pourtant, en revanche & mais, or, alors que, bien que, quoique \\
\textbf{Illustration / Précision} & par exemple, notamment, en particulier, c'est-à-dire & qui, que, dont, où, lequel (relatifs) \\
\textbf{Chronologie / Temps} & ensuite, puis, ultérieurement, auparavant, simultanément & quand, lorsque, avant que, après que \\ \hline
\end{tabularx}
\end{center}
\end{propriete}

\begin{aretenir}[Usage dans le CV / Lettre de motivation]
\begin{itemize}
\item \textbf{Phrases nominales (CV) :} Privilégier la concision. Les connecteurs introduisent des compléments ou relient des items (ex: \textit{« Gestion budgétaire ; \textbf{par ailleurs}, reporting mensuel »}).
\item \textbf{Phrases verbales (Lettre/Objectif) :} Utiliser les pronoms relatifs (\textit{qui, que, dont, où}) pour éviter les répétitions et fluidifier : \textit{« Stage \textbf{où} j'ai acquis... » ; « Projet \textbf{dont} je suis fier. »}
\item \textbf{Variété :} Ne pas abuser de « et », « mais », « aussi ». Varier : \textit{de surcroît, néanmoins, en effet, par conséquent}.
\end{itemize}
\end{aretenir}

\courssub{Rédaction du CV : travail préparatoire}

\begin{methode}[Analyser une offre d'emploi : Tableau de correspondance]
\begin{enumerate}
\item \textbf{Surligner} dans l'offre : Diplômes, Expérience, Compétences techniques (mots-clés), Soft skills, Contraintes (permis, déplacements, langues).
\item \textbf{Construire le tableau} (5 à 7 lignes) : Colonne 1 = Exigence de l'offre ; Colonne 2 = Ma preuve / Mon atout (chiffré si possible) ; Colonne 3 = Rubrique CV concernée.
\item \textbf{Ordonner les rubriques} du CV selon les points forts identifiés.
\item \textbf{Rédiger l'accroche (Objectif)} en reprenant 2-3 mots-clés majeurs de l'offre.
\end{enumerate}
\end{methode}

\begin{exemplebox}[Exemple de ligne de tableau de correspondance]
\begin{center}
\begin{tabularx}{\linewidth}{|l|X|l|}\hline
\rowcolor{popblueL} \textbf{Exigence offre} & \textbf{Mon atout / Preuve} & \textbf{Rubrique CV} \\ \hline
Anglais technique (notices, formation client) & TOEIC 820 — Formation utilisateurs en anglais chez MedEquip (18 mois) & Compétences / Expériences \\ \hline
\end{tabularx}
\end{center}
\end{exemplebox}

\courssub{Activités d'intégration, compte rendu et remédiation}

\courssub{Je vérifie mes connaissances}

\begin{exercice}[QCM : Structure et vocabulaire du CV]
Pour chaque question, une seule réponse est correcte. Recopiez le numéro et la lettre choisie.
\begin{enumerate}
\item Dans un CV de type \textbf{chronologique inverse}, les expériences professionnelles sont présentées :
\begin{itemize}
\item[A.] De la plus ancienne à la plus récente.
\item[B.] De la plus récente à la plus ancienne.
\item[C.] Par ordre alphabétique des entreprises.
\item[D.] En fonction de leur pertinence uniquement.
\end{itemize}

\item La rubrique \textbf{« Objectif professionnel »} (ou « Profil ») doit :
\begin{itemize}
\item[A.] Résumer toute votre vie scolaire depuis le primaire.
\item[B.] Lister vos loisirs et centres d'intérêt.
\item[C.] Cibler le poste visé et mettre en avant votre valeur ajoutée en 2-3 lignes.
\item[D.] Détailler vos prétentions salariales.
\end{itemize}

\item Parmi les éléments suivants, lequel est \textbf{inutile} voire \textbf{déconseillé} dans un CV moderne au Cameroun (sauf cas particuliers) ?
\begin{itemize}
\item[A.] L'adresse e-mail professionnelle.
\item[B.] La photo d'identité récente.
\item[C.] Le numéro de carte nationale d'identité (CNI).
\item[D.] Le lien vers un profil LinkedIn à jour.
\end{itemize}

\item Les \textbf{adverbes de liaison} (par conséquent, toutefois, de plus, en effet) servent à :
\begin{itemize}
\item[A.] Remplacer les verbes conjugués.
\item[B.] Assurer la cohérence textuelle en marquant la relation logique entre deux phrases ou énoncés.
\item[C.] Introduire uniquement des phrases interrogatives.
\item[D.] Accorder le participe passé.
\end{itemize}
\end{enumerate}
\end{exercice}

\begin{exercice}[Vrai / Faux justifié]
Indiquez si les affirmations suivantes sont vraies ou fausses. Justifiez brièvement chaque réponse par une phrase.
\begin{enumerate}
\item Un CV doit impérativement tenir sur une seule page, quel que soit le nombre d'années d'expérience du candidat.
\item La rubrique « Centres d'intérêt » est facultative mais permet de révéler des qualités transverses (esprit d'équipe, engagement associatif).
\item L'utilisation de pronoms relatifs (qui, que, dont, où) alourdit le style d'un CV ; il faut les éviter au profit de phrases nominales.
\item Dans la rubrique « Formation », on mentionne uniquement les diplômes obtenus avec mention ; les échecs sont à masquer.
\end{enumerate}
\end{exercice}

\begin{exercice}[Définitions et terminologie technique]
Donnez une définition précise (une à deux phrases) pour chacun des termes suivants dans le contexte de la rédaction d'un CV :
\begin{enumerate}
\item \textbf{Mots-clés (keywords) :}
\item \textbf{Compétences transférables (soft skills / hard skills) :}
\item \textbf{Locution adverbiale :}
\item \textbf{CV fonctionnel (ou par compétences) :}
\end{enumerate}
\end{exercice}

\begin{exercice}[Outils de liaison : identification et classe grammaticale]
Dans les phrases suivantes, extraites de lettres de motivation ou de CV, soulignez l'outil de liaison (adverbe de liaison, locution adverbiale ou pronom relatif) et précisez sa nature grammaticale.
\begin{enumerate}
\item « J'ai dirigé une équipe de cinq personnes ; \textbf{par conséquent}, j'ai développé mon sens du leadership. »
\item « Le stage \textbf{que} j'ai effectué à la SONARA m'a permis de maîtriser les normes HSE. »
\item « Titulaire d'un BTS en Gestion ; \textbf{de surcroît}, je parle couramment l'anglais et l'allemand. »
\item « L'entreprise \textbf{dont} vous vantez les valeurs sociales correspond à mes aspirations. »
\end{enumerate}
\end{exercice}

\courssub{Je m'entraîne}

\begin{exercice}[Classement des informations dans les rubriques du CV]
Voici une liste d'informations concernant \textbf{M. Émile K., 24 ans, titulaire d'un DUT en Génie Civil (ENSTP Yaoundé, 2023)}, à la recherche de son premier emploi comme Conducteur de travaux junior.
Classez chaque item dans la rubrique appropriée du CV (En-tête, Objectif, Formation, Expériences/Stages, Compétences techniques, Compétences transverses/Langues, Centres d'intérêt, Références). Proposez ensuite l'ordre logique des rubriques pour ce profil « jeune diplômé ».
\begin{enumerate}
\item Permis B (voiture) + Permis C (poids lourds).
\item Stage de fin d'études (3 mois) : Suivi de chantier route Nationale 10 (Entreprise Razel-Bec).
\item Maîtrise d'AutoCAD, Revit, MS Project, Suite Office.
\item Né le 12/05/2000 à Douala — Célibataire — BP 1234 Yaoundé — 6XX XX XX XX — emile.k@email.com.
\item Baccalauréat F4 (Génie Civil) — Mention Assez Bien — Lycée Technique de Douala (2020).
\item Anglais : Niveau B2 (TOEIC 785) — Allemand : Notions (A1).
\item Membre actif de l'association « Bâtisseurs Solidaires » (rénovation écoles rurales) — Footballeur en club (milieu de terrain).
\item « Jeune ingénieur dynamique et rigoureux, je souhaite mettre ma formation en Génie Civil et mon expérience de terrain au service de vos chantiers en tant que Conducteur de travaux junior. »
\item M. Jean B., Chef de projet chez Razel-Bec (Tél : 6XX XX XX XX) — Mme Ngo M., Enseignante ENSTP (Tél : 6XX XX XX XX).
\item DUT Génie Civil — Option Bâtiment/Travaux Publics — ENSTP Yaoundé (2021-2023).
\item Stage ouvrier (1 mois) : Maçonnerie traditionnelle (Entreprise locale, Bafoussam).
\item Capacité à lire des plans d'exécution, gestion des approvisionnements, encadrement d'une petite équipe.
\end{enumerate}
\end{exercice}

\begin{exercice}[Rédaction de l'accroche « Objectif professionnel »]
Rédigez l'accroche (rubrique \textbf{Objectif professionnel / Profil}) pour les trois profils ci-dessous.
\textbf{Contraintes :} 3 phrases maximum par accroche. Utilisez au moins \textbf{deux pronoms relatifs} (qui, que, dont, où, lequel, etc.) au total dans chaque proposition. Mettez en gras les pronoms relatifs utilisés.
\begin{enumerate}
\item \textbf{Profil A :} Technicien de maintenance industrielle (BTS), 3 ans d'expérience en usine agroalimentaire (Chococam), cherche poste de Superviseur maintenance. Points forts : diagnostic pannes, GMAO, management d'équipe de 3 techniciens.
\item \textbf{Profil B :} Assistante de direction (DUT Secrétariat Bureautique), 5 ans d'expérience en cabinet d'avocats, bilingue français/anglais, cherche poste de Office Manager. Points forts : organisation agenda, gestion dossiers confidentiels, protocole.
\item \textbf{Profil C :} Développeur Web Full Stack (Licence Pro), freelance 2 ans (sites e-commerce, applis mobiles), cherche CDI en agence digitale. Points forts : React, Node.js, PHP, base de données, méthode Agile/Scrum.
\end{enumerate}
\end{exercice}

\begin{exercice}[Correction et amélioration : Les 12 erreurs]
Le CV ci-dessous (présenté sous forme linéaire) contient \textbf{12 erreurs} (structure, ordre, informations inutiles/redondantes, fautes de grammaire/syntaxe, manque de mots-clés). Identifiez-les toutes en les numérotant, puis proposez la version corrigée et mise en page (structure externe) sur une feuille annexe.

\textbf{CV de M. Paul D.}
\begin{quote}
\textsc{curriculum vitae}. Nom : PAUL D. — Né le 10/01/1995 à Bafia. Situation : Marié, 3 enfants. CNI n° 102030405. Groupe sanguin O+. Taille 1m75. Adresse : Quartier Mvog-Ada, Yaoundé. Tél : 690 00 00 00. Email : pauld1995@yahoo.fr.
\textbf{OBJECTIF} : Je veux un travail pour gagner ma vie et aider ma famille. Je suis courageux et honnête.
\textbf{EXPERIENCES} :
2018-2020 : Vendeur dans la boutique de mon oncle à Mfoundi. (Vente habits, caisse).
2021-2023 : Chauffeur livreur pour Société X. Conduite camion, livraison colis, chargement/déchargement.
2015-2017 : Étudiant à l'Université de Yaoundé I (Droit) — \textit{Échec en L2}.
\textbf{FORMATION} :
2023 : Permis poids lourds (C) obtenu.
2014 : Baccalauréat A4 (Lettres) — Mention Passable.
\textbf{COMPETENCES} : Conduite. Force physique. Je sais parler aux clients. Word un peu.
\textbf{LANGUES} : Français (langue maternelle). Ewondo (courant). Anglais (nul).
\textbf{LOISIRS} : Boire entre amis, regarder la télé, dormir, politique.
\textbf{REFERENCES} : Mon oncle (patron boutique). Mon curé.
\textit{Fait à Yaoundé, le 15/10/2023. Signature.}
\end{quote}
\end{exercice}

\begin{exercice}[Relier des énoncés par des outils de liaison]
Reliez les deux phrases de chaque item en une seule phrase complexe ou en deux phrases cohérentes, en utilisant l'outil de liaison imposé entre parenthèses. Faites les accords nécessaires.
\begin{enumerate}
\item Le candidat maîtrise parfaitement l'anglais. Il a travaillé deux ans à Londres. (\textbf{dont})
\item Elle a obtenu son BTS avec mention Bien. Elle a été major de sa promotion. (\textbf{par conséquent})
\item Ce poste exige une grande rigueur. Le profil de M. N. correspond parfaitement. (\textbf{en effet})
\item L'entreprise recherche un commercial bilingue. Mme K. parle couramment français et espagnol. (\textbf{or} / \textbf{cependant} — choisissez le sens logique)
\item Il a géré un budget de 50 millions FCFA. Il n'a jamais dépassé les prévisions. (\textbf{de surcroît})
\item Le stage s'est déroulé à Douala. Il a découvert la logistique portuaire. (\textbf{où})
\item Les délais étaient courts. L'équipe a livré le projet à temps. (\textbf{néanmoins})
\item Elle est titulaire d'un Master en Finance. Elle postule comme Trésorière. (\textbf{par ailleurs})
\item Le logiciel est complexe. Il a suivi une formation certifiante. (\textbf{car} / \textbf{puisque})
\item Nous attendons votre réponse. Nous restons à votre disposition pour un entretien. (\textbf{en attendant})
\end{enumerate}
\end{exercice}

\courssub{Je me prépare au Bac}

\begin{exercice}[Type Baccalauréat Technique — Sujet complet : Rédaction de CV à partir d'une offre]
\textbf{Durée conseillée : 1h30 — Coefficient 3}

\textbf{Document 1 : Offre d'emploi (parue dans \textit{Cameroon Tribune} du 12 mars 2024)}

\begin{center}
\begin{tabularx}{\linewidth}{|X|}\hline
\rowcolor{popblueL} \textbf{AVIS DE RECRUTEMENT — SOCIÉTÉ CAMEROUNAISE D'ÉQUIPEMENTS MÉDICAUX (SOCAMED)} \\
\hline
\textbf{Poste :} \textbf{Technico-Commercial Itinérant (H/F)} — Basé à Douala, déplacements nationaux (3 semaines / mois). \\
\textbf{Missions :}
\begin{itemize}
\item Prospecter et fidéliser une clientèle de cliniques, hôpitaux et pharmacies.
\item Présenter et démontrer les équipements (imagerie, laboratoire, bloc opératoire).
\item Négocier les contrats, assurer le suivi technique de premier niveau et le reporting d'activité hebdomadaire.
\end{itemize}
\textbf{Profil recherché :}
\begin{itemize}
\item Diplôme : BTS / DUT / Licence Pro en \textbf{Maintenance Biomédicale}, \textbf{Génie Biomédical} ou \textbf{Technico-Commercial} (spécialité médicale).
\item Expérience : 2 ans minimum sur un poste similaire (stage long validé accepté pour jeune diplômé brillant).
\item Compétences techniques : Connaissance matériel médical, normes IEC/ISO, anglais technique (notice, formation client).
\item Compétences transverses : Permis B obligatoire, autonomie, résistance au stress, sens du service, éthique médicale.
\item Maîtrise pack Office + CRM (Salesforce ou équivalent) un atout.
\end{itemize}
\textbf{Dossier :} CV + Lettre de motivation + Copies diplômes + Attestations de travail → \textbf{DRH SOCAMED, BP 456 Douala} ou \textbf{rh@socamed.cm} — \textit{Date limite : 30 mars 2024}. \\
\hline
\end{tabularx}
\end{center}

\textbf{Travail à faire :}
Vous êtes \textbf{Armand T., 26 ans}, titulaire d'une \textbf{Licence Professionnelle en Maintenance Biomédicale} (IUT Douala, 2022, Mention Bien). Vous avez effectué un \textbf{stage de fin d'études (6 mois)} au service maintenance du \textbf{CHU de Douala} (gestion parc matériel, interventions curatives/préventives, relation fournisseurs). Vous avez ensuite été embauché 18 mois comme \textbf{Technicien de maintenance itinérant} chez \textbf{MedEquip Services} (Douala) : interventions chez clients privés/cliniques, formation utilisateurs, rédaction rapports, utilisation CRM Salesforce. Vous êtes \textbf{bilingue français/anglais (TOEIC 820)}, titulaire du \textbf{Permis B}, véhicule personnel. Références : Chef de service Maintenance CHU Douala + Directeur Technique MedEquip.

\begin{enumerate}
\item \textbf{Analyse de l'offre (3 pts) :} Relevez 5 mots-clés (compétences, qualités, diplômes) sur lesquels vous devez faire apparaître votre adéquation dans votre CV.
\item \textbf{Travail préparatoire (4 pts) :} Complétez le tableau de correspondance « Exigences de l'offre $\leftrightarrow$ Mes preuves/Atouts » (5 lignes minimum).
\item \textbf{Rédaction du CV (13 pts) :} Rédigez votre \textbf{CV complet, mis en page} (structure externe : en-tête, rubriques ordonnées, mise en forme lisible ; structure interne : phrases nominales, verbes d'action, mots-clés, outils de liaison). Le CV doit tenir sur \textbf{une page A4}. Adaptez l'ordre des rubriques à votre profil (expérience significative). Soignez la rubrique « Objectif » (3 lignes max, 2 pronoms relatifs minimum).
\end{enumerate}
\end{exercice}

\begin{exercice}[Type Probatoire — Étude comparative de deux CV]
\textbf{Durée conseillée : 1h — Coefficient 2}

Deux candidats, \textbf{M. Aliou S.} et \textbf{Mme Chantal N.}, postulent au poste de \textbf{Comptable Unique} dans une PME agroalimentaire à Bafoussam (Offre : Diplôme BTS Comptabilité/Gestion, 3 ans exp., Maîtrise OHADA, Sage Saari, Excel avancé, Rigueur, Discrétion).

\textbf{CV 1 — Aliou S. (29 ans)}
\begin{itemize}
\item \textbf{Objectif :} « Comptable rigoureux, je veux intégrer votre structure pour évoluer. »
\item \textbf{Expériences :}
  \begin{itemize}
  \item 2021-2023 : Comptable Unique — SARL « AgroBois » (Bois). Tenue compta générale, déclarations fiscales (IS, TVA), paie 15 pers., clôture mensuelle. Logiciel : Sage Saari v10. Excel (TCD, RechercheV).
  \item 2019-2021 : Assistant Comptable — Cabinet « ExpertCom » (Yaoundé). Saisie, lettrage, révision dossiers clients divers. OHADA révisé.
  \end{itemize}
\item \textbf{Formation :} BTS Comptabilité-Gestion (ISTA Douala, 2018) — Mention AB. Bac G2 (2015).
\item \textbf{Compétences :} OHADA (SYSCOHADA), Sage Saari, Excel Avancé, Fiscalité camerounaise, Paie. Anglais technique (B1).
\item \textbf{Loisirs :} Lecture (économie), Marathon (discipline).
\item \textbf{Références :} Directeur Financier AgroBois + Associé ExpertCom.
\end{itemize}

\textbf{CV 2 — Chantal N. (32 ans)}
\begin{itemize}
\item \textbf{Objectif :} « Dynamique et sérieuse, je cherche un emploi stable près de chez moi. »
\item \textbf{Expériences :}
  \begin{itemize}
  \item 2018-2023 : Gérante de sa propre boutique « Chez Chantal » (Vêtements, Bafoussam). Gestion caisse, stocks, fournisseurs, 2 employés. \textit{(Pas de comptabilité formelle OHADA)}.
  \item 2015-2017 : Secrétaire Comptable — ONG « Santé Plus ». Saisie dépenses, suivi budget projet, reporting bailleur. Logiciel : Excel simple + QuickBooks (non OHADA).
  \end{itemize}
\item \textbf{Formation :} BTS Secrétariat Bureautique (Option Compta) — Lycée Technique Bafoussam (2014). Bac A4 (2011).
\item \textbf{Compétences :} Excel (moyen), QuickBooks, Gestion administrative, Accueil, Français courant. Anglais (notions).
\item \textbf{Loisirs :} Chant chorale, Cuisine, Association femmes entrepreneures.
\item \textbf{Références :} Pasteur de l'église + Président Association femmes.
\end{itemize}

\textbf{Travail demandé :}
\begin{enumerate}
\item \textbf{Grille d'analyse comparative (6 pts) :} Construisez un tableau à 3 colonnes (Critères de l'offre / Aliou S. / Chantal N.) évaluant l'adéquation sur 6 critères clés (Diplôme, Expérience pertinente, Maîtrise OHADA/Sage, Excel, Qualités transverses/Références, Présentation CV/Objectif). Cotez : \textbf{+++} (Parfait), \textbf{++} (Bon), \textbf{+} (Faible), \textbf{0} (Absent/Inadapté).
\item \textbf{Texte argumenté (14 pts) :} En vous appuyant \textbf{uniquement} sur la grille et les CV, rédigez un texte argumenté d'environ \textbf{15 lignes} (180-200 mots) pour recommander l'un des deux candidats. Structurez votre réponse : \textbf{Thèse claire} (choix), \textbf{Arguments chiffrés/factuels} (3 minimum), \textbf{Nuance/Point de vigilance} sur le retenu, \textbf{Conclusion professionnelle}. Utilisez au moins \textbf{3 connecteurs logiques variés} (adverbes de liaison, locutions, conjonctions).
\end{enumerate}
\end{exercice}

\begin{exercice}[Type Bac Technique — Rédaction intégrée et correction linguistique]
\textbf{Durée conseillée : 1h}

\textbf{Partie A : Rédaction de la rubrique « Compétences » enrichie (6 pts)}
À partir des notes brutes ci-dessous, rédigez la rubrique \textbf{« Compétences Techniques \& Fonctionnelles »} d'un CV pour un poste de \textbf{Chef de Projet Junior BTP}. Organisez par sous-catégories (Gestion de projet, Technique/Chantier, Outils numériques, Réglementation/Qualité). Utilisez des \textbf{phrases nominales percutantes}, intégrez des \textbf{mots-clés} du secteur et reliez les items par une ponctuation adaptée (points-virgules, deux-points) ou des outils de liaison (au moins 2 locutions adverbiales du type \textit{notamment, en particulier, par ailleurs}).

\textit{Notes brutes :} Planning (MS Project, Primavera P6) — Budget (suivi coûts, appels d'offres, DPGF) — Management (équipe 10-15 pers, sous-traitants, réunions hebdo) — Technique (Béton armé, VRD, lecture plans EXE/DAO AutoCAD/Revit) — Normes (NF DTU, Eurocodes, HQE, Sécurité HSE) — Reporting (tableaux de bord, reporting client mensuel, anglais technique B2) — Contractuel (Marchés publics/privés, CCAG, avenants).

\textbf{Partie B : Correction linguistique et stylistique (4 pts)}
Le paragraphe suivant est extrait de la rubrique « Expériences » d'un CV. Il contient 6 erreurs (syntaxe, accord, vocabulaire, liaison, registre). Recopiez le paragraphe \textbf{corrigé} en surlignant (gras) vos corrections.

\begin{quote}
« \textit{De 2020 à 2022, j'ai occupé le poste de Superviseur Maintenance à l'usine SOSUCAM. J'étais chargé de planifier les arrêts de production, gérer une équipe de 8 techniciens et de veiller au respect des normes sécurité. Grace à mon organisation, le taux de disponibilité des machines est passé de 85\% à 95\%. De plus, j'ai mis en place une GMAO qui a permis de réduire les coûts de maintenance de 15\%. Cependant, j'ai du partir pour raison familiale.} »
\end{quote}

\textbf{Partie C : Lien syntaxique (5 pts)}
Reliez les trois propositions suivantes en \textbf{une seule phrase complexe correcte et élégante} (style CV / rapport professionnel), en utilisant \textbf{au moins un pronom relatif} et \textbf{un adverbe de liaison}.
\begin{itemize}
\item Proposition 1 : Le candidat a piloté la construction d'un hôpital de 150 lits.
\item Proposition 2 : Ce chantier a été livré avec 2 semaines d'avance.
\item Proposition 3 : Le budget initial a été respecté à 99,5\%.
\end{itemize}
\end{exercice}

\newpage

\coursec{Corrigés des exercices}

\coursec{Corrigés des exercices — Le Curriculum Vitae (CV)}

\begin{corrige}[Exercice 1 — QCM : Structure et vocabulaire du CV]
\begin{enumerate}
\item \textbf{B.} Le CV chronologique inverse (anti-chronologique) présente les expériences de la plus récente à la plus ancienne : c'est le standard le plus utilisé car il valorise la progression et met en avant l'expérience actuelle.
\item \textbf{C.} L'objectif professionnel est une accroche de 2-3 lignes qui cible le poste visé et met en avant la valeur ajoutée du candidat. Il ne s'agit ni d'une autobiographie (A), ni d'une liste de loisirs (B), ni de prétentions salariales (D), qui ne figurent jamais dans un CV.
\item \textbf{C.} Le numéro de CNI est une information personnelle sensible, inutile et déconseillée dans un CV moderne (risque d'usurpation, non pertinent pour le recrutement). L'e-mail professionnel (A), la photo professionnelle (B) et le lien LinkedIn (D) sont au contraire recommandés.
\item \textbf{B.} Les adverbes de liaison assurent la \cle{cohérence textuelle} en marquant la relation logique (cause, conséquence, opposition, addition) entre deux énoncés.
\end{enumerate}
\end{corrige}

\begin{corrige}[Exercice 2 — Vrai / Faux justifié]
\begin{enumerate}
\item \textbf{FAUX.} La règle d'une page s'applique aux profils standards ; un cadre dirigeant avec plus de 15 ans d'expérience peut légitimement tenir sur deux pages, à condition que chaque information soit pertinente.
\item \textbf{VRAI.} Cette rubrique facultative permet de révéler des qualités transverses : un sport collectif prouve l'esprit d'équipe, l'engagement associatif montre le sens des responsabilités.
\item \textbf{FAUX.} Les pronoms relatifs (\emph{qui, que, dont, où}) fluidifient le style et évitent les répétitions ; ils sont utiles dans l'accroche et la lettre de motivation, même si le CV privilégie les phrases nominales.
\item \textbf{FAUX.} On mentionne tous les diplômes obtenus, pas seulement ceux avec mention ; en revanche, les années d'échec n'apparaissent pas car seuls les diplômes \emph{obtenus} figurent dans la rubrique Formation (on ne « masque » rien, on n'écrit que ce qui est validé).
\end{enumerate}
\end{corrige}

\begin{corrige}[Exercice 3 — Définitions et terminologie technique]
\begin{enumerate}
\item \textbf{Mots-clés (keywords) :} termes précis repris de l'offre d'emploi (diplômes, compétences, logiciels, qualités) que le candidat intègre dans son CV pour démontrer son adéquation au poste et être repéré par les logiciels de tri automatique (ATS).
\item \textbf{Compétences transférables :} compétences utilisables dans plusieurs contextes professionnels. On distingue les \emph{hard skills} (savoir-faire techniques mesurables : logiciels, langues, certifications) et les \emph{soft skills} (qualités comportementales : rigueur, esprit d'équipe, gestion du stress), qui doivent être prouvées par un contexte.
\item \textbf{Locution adverbiale :} groupe de mots invariable jouant le rôle d'un adverbe (ex. \emph{par conséquent, de surcroît, en revanche}) et exprimant une relation logique ou une circonstance.
\item \textbf{CV fonctionnel (ou par compétences) :} type de CV qui regroupe les expériences par grandes compétences (gestion, technique, relationnel) plutôt que par ordre chronologique ; il est adapté aux reconversions, aux parcours interrompus ou aux jeunes diplômés.
\end{enumerate}
\end{corrige}

\begin{corrige}[Exercice 4 — Outils de liaison : identification et classe grammaticale]
\begin{enumerate}
\item \textbf{par conséquent} : \emph{locution adverbiale} de conséquence (elle relie logiquement la cause — diriger une équipe — à la conséquence — développer le leadership).
\item \textbf{que} : \emph{pronom relatif} complément d'objet direct (représente « le stage », COD de « j'ai effectué »).
\item \textbf{de surcroît} : \emph{locution adverbiale} d'addition / renforcement (elle ajoute un atout supplémentaire au diplôme).
\item \textbf{dont} : \emph{pronom relatif} complément introduit par la préposition « de » (il représente « l'entreprise », complément de « vous vantez les valeurs \emph{de} »).
\end{enumerate}
\end{corrige}

\begin{corrige}[Exercice 5 — Classement des informations dans les rubriques du CV]
\begin{center}
\begin{tabularx}{\linewidth}{|l|X|}\hline
\rowcolor{popblueL} \textbf{Rubrique} & \textbf{Items} \\ \hline
En-tête & 4 (identité et contacts — \emph{à épurer} : on retire « né le 12/05/2000 », « célibataire » ; on garde nom, ville, téléphone, e-mail) \\ \hline
Objectif & 8 \\ \hline
Formation & 10 (DUT 2021-2023) puis 5 (Bac F4, 2020) \\ \hline
Expériences / Stages & 2 (stage ENSTP, 3 mois, Razel-Bec) puis 11 (stage ouvrier, 1 mois, Bafoussam) \\ \hline
Compétences techniques & 3 (AutoCAD, Revit, MS Project, Office) ; 1 (Permis B + C) ; 12 (lecture de plans, approvisionnements, encadrement) \\ \hline
Langues & 6 (Anglais B2 — TOEIC 785 ; Allemand A1) \\ \hline
Centres d'intérêt & 7 (association « Bâtisseurs Solidaires », football en club) \\ \hline
Références & 9 (M. Jean B., Mme Ngo M.) \\ \hline
\end{tabularx}
\end{center}

\textbf{Ordre logique pour un jeune diplômé :} En-tête $\rightarrow$ Objectif $\rightarrow$ \textbf{Formation} (placée avant l'expérience car c'est son point fort) $\rightarrow$ Expériences/Stages $\rightarrow$ Compétences techniques $\rightarrow$ Langues $\rightarrow$ Centres d'intérêt $\rightarrow$ Références.

\textbf{Point de vigilance :} l'item 4 contient des informations à supprimer (date de naissance, situation matrimoniale) : un CV moderne ne mentionne ni l'état civil détaillé ni l'âge.
\end{corrige}

\begin{corrige}[Exercice 6 — Rédaction de l'accroche « Objectif professionnel »]
Propositions corrigées (les pronoms relatifs sont en gras ; toute formulation équivalente respectant les contraintes est acceptée) :

\textbf{Profil A :} « Technicien de maintenance industrielle fort de trois ans d'expérience en milieu agroalimentaire, je maîtrise le diagnostic de pannes et la GMAO, outils \textbf{grâce auxquels} j'ai optimisé la disponibilité des lignes de production. Le poste de Superviseur maintenance \textbf{que} vous proposez correspond au management d'équipe \textbf{dont} j'ai fait la preuve auprès de trois techniciens. »

\textbf{Profil B :} « Assistante de direction bilingue français-anglais, je capitalise cinq années d'expérience en cabinet d'avocats, environnement exigeant \textbf{dans lequel} j'ai géré des agendas complexes et des dossiers confidentiels. Le poste d'Office Manager \textbf{que} vous offrez mobilise le sens de l'organisation et du protocole \textbf{dont} je suis reconnue. »

\textbf{Profil C :} « Développeur Web Full Stack freelance depuis deux ans, j'ai conçu des sites e-commerce et des applications mobiles \textbf{qui} ont fidélisé une clientèle variée. L'agence digitale \textbf{à laquelle} je postule trouvera en moi un profil maîtrisant React, Node.js et la méthode Agile, compétences \textbf{dont} mes réalisations témoignent. »

\textbf{Points de vigilance :} 3 phrases maximum ; au moins 2 pronoms relatifs par accroche ; mention du poste visé + valeur ajoutée ; registre professionnel (pas de « je veux », « je cherche un travail »).
\end{corrige}

\begin{corrige}[Exercice 7 — Correction et amélioration : Les 12 erreurs]
\textbf{Les 12 erreurs identifiées :}
\begin{enumerate}
\item \textbf{Informations sensibles inutiles :} CNI n° 102030405 et groupe sanguin O+ (à supprimer).
\item \textbf{Situation matrimoniale détaillée :} « Marié, 3 enfants » (à supprimer).
\item \textbf{Taille (1m75) :} information corporelle hors de propos (à supprimer).
\item \textbf{E-mail non professionnel :} « pauld1995@yahoo.fr » révèle l'année de naissance ; préférer \texttt{paul.d@domaine.pro}.
\item \textbf{Objectif mal rédigé :} « Je veux un travail pour gagner ma vie » est un registre familier, centré sur le candidat et non sur le poste ; il faut une accroche ciblée (poste + valeur ajoutée).
\item \textbf{Ordre chronologique fautif :} les expériences vont de 2018 à 2023 puis reviennent à 2015 ; il faut l'ordre \emph{anti-chronologique} (2021-2023, puis 2018-2020).
\item \textbf{Échec mentionné :} « Échec en L2 » ne figure jamais dans un CV ; seuls les diplômes obtenus apparaissent (la formation non validée peut être formulée « Études de Droit, niveau L2 » si utile, sinon omise).
\item \textbf{Compétences vagues et non professionnelles :} « Force physique », « Je sais parler aux clients », « Word un peu » (registre familier, non chiffré) ; reformuler : « Relation client », « Pack Office : Word (notions) ».
\item \textbf{Niveau de langue fautif :} « Anglais (nul) » est à supprimer ou à formuler « notions » ; une compétence nulle ne se mentionne pas.
\item \textbf{Loisirs inadaptés :} « Boire entre amis, dormir, politique » discréditent le candidat (la politique est un sujet sensible à éviter).
\item \textbf{Références non professionnelles :} « Mon oncle » (conflit d'intérêt familial) et « Mon curé » (religion, à proscrire) ; il faut des références professionnelles identifiables (fonction, téléphone).
\item \textbf{Mentions superflues :} « Fait à Yaoundé, le... Signature » est inutile dans un CV moderne (contrairement à la lettre manuscrite d'autrefois) ; le titre « curriculum vitae » en minuscules est aussi négligé.
\end{enumerate}

\textbf{Version corrigée (structure externe) :}
\begin{center}
\begin{tabularx}{\linewidth}{|X|}\hline
\rowcolor{popblueL} \textbf{PAUL D. — Chauffeur-Livreur / Vendeur polyvalent} \\
Quartier Mvog-Ada, Yaoundé — Tél : 690 00 00 00 — \texttt{paul.d@email.pro} \\ \hline
\textbf{OBJECTIF :} Chauffeur-livreur titulaire du permis C, je mets mon expérience de la livraison et de la vente au service de votre logistique. Rigoureux et ponctuel, je maîtrise la relation client. \\ \hline
\textbf{EXPÉRIENCES :} \\
2021-2023 : Chauffeur-livreur, Société X, Yaoundé — Livraison de colis, conduite de camion, chargement/déchargement. \\
2018-2020 : Vendeur, boutique Mfoundi, Yaoundé — Vente de prêt-à-porter, tenue de caisse, accueil client. \\ \hline
\textbf{FORMATION :} 2023 : Permis C (poids lourds). 2014 : Baccalauréat A4, mention Passable. \\ \hline
\textbf{COMPÉTENCES :} Conduite poids lourds ; Relation client ; Tenue de caisse ; Pack Office (Word : notions). \\ \hline
\textbf{LANGUES :} Français (courant) ; Ewondo (courant). \\ \hline
\textbf{CENTRES D'INTÉRÊT :} Sport, lecture. \\ \hline
\end{tabularx}
\end{center}
\end{corrige}

\begin{corrige}[Exercice 8 — Relier des énoncés par des outils de liaison]
\begin{enumerate}
\item Le candidat maîtrise parfaitement l'anglais, langue \textbf{dont} il a consolidé la pratique en travaillant deux ans à Londres. (\emph{dont} = complément du nom « la pratique ».)
\item Elle a obtenu son BTS avec mention Bien ; \textbf{par conséquent}, elle a été major de sa promotion. (Relation de conséquence.)
\item Ce poste exige une grande rigueur ; \textbf{en effet}, le profil de M. N. y correspond parfaitement. (\emph{en effet} introduit la confirmation/explication.)
\item L'entreprise recherche un commercial bilingue ; \textbf{or}, Mme K. parle couramment français et espagnol. (\emph{or} marque ici l'adéquation/l'argument décisif ; \emph{cependant} conviendrait si Mme K. ne parlait qu'une langue.)
\item Il a géré un budget de 50 millions FCFA ; \textbf{de surcroît}, il n'a jamais dépassé les prévisions. (Addition/renforcement.)
\item Le stage s'est déroulé à Douala, ville \textbf{où} il a découvert la logistique portuaire. (\emph{où} = pronom relatif de lieu.)
\item Les délais étaient courts ; l'équipe a \textbf{néanmoins} livré le projet à temps. (Concession/opposition.)
\item Elle est titulaire d'un Master en Finance ; \textbf{par ailleurs}, elle postule comme Trésorière. (Addition d'information.)
\item Le logiciel est complexe ; il a \textbf{donc} suivi une formation certifiante, \textbf{car} il voulait le maîtriser. — Ou plus simplement : Il a suivi une formation certifiante \textbf{car} le logiciel est complexe. (\emph{car} introduit la cause.)
\item \textbf{En attendant} votre réponse, nous restons à votre disposition pour un entretien. (Formule consacrée de correspondance professionnelle ; \emph{en attendant} est invariable.)
\end{enumerate}
\textbf{Point de vigilance :} le point-virgule précède les adverbes de liaison ; \emph{en attendant} s'écrit en deux mots et reste invariable.
\end{corrige}

\begin{corrige}[Exercice 9 — Type Baccalauréat Technique : Rédaction de CV à partir d'une offre]
\textbf{Barème indicatif :} Analyse de l'offre 3 pts ; tableau de correspondance 4 pts ; CV rédigé 13 pts (en-tête 1,5 ; objectif 2 ; expériences 3 ; formation 2 ; compétences/langues 2 ; mise en page et ordre 1,5 ; langue et outils de liaison 1).

\textbf{1. Analyse de l'offre (3 pts) — 5 mots-clés attendus (au choix) :} Maintenance biomédicale ; Technico-commercial / prospection et fidélisation ; Anglais technique ; Permis B + mobilité nationale ; CRM (Salesforce) ; normes IEC/ISO ; autonomie et résistance au stress.

\textbf{2. Tableau de correspondance (4 pts) :}
\begin{center}
\begin{tabularx}{\linewidth}{|l|X|l|}\hline
\rowcolor{popblueL} \textbf{Exigence de l'offre} & \textbf{Mon atout / Preuve} & \textbf{Rubrique CV} \\ \hline
Diplôme Maintenance biomédicale & Licence Pro Maintenance Biomédicale, IUT Douala, 2022, mention Bien & Formation \\ \hline
2 ans d'expérience (stage long accepté) & 18 mois Technicien itinérant MedEquip + stage 6 mois CHU Douala & Expériences \\ \hline
Anglais technique & TOEIC 820 ; formation d'utilisateurs en anglais chez MedEquip & Compétences / Langues \\ \hline
Permis B + déplacements & Permis B, véhicule personnel, disponible pour missions nationales & Compétences \\ \hline
CRM / Pack Office & Salesforce utilisé 18 mois (rapports hebdomadaires) ; Pack Office & Compétences \\ \hline
Relation client / négociation & Interventions chez cliniques privées, formation utilisateurs, relation fournisseurs (CHU) & Expériences \\ \hline
\end{tabularx}
\end{center}

\textbf{3. CV rédigé (13 pts) — corrigé type :}
\begin{center}
\begin{tabularx}{\linewidth}{|X|}\hline
\rowcolor{popblueL} \textbf{Armand T. — Technico-Commercial Itinérant (Maintenance Biomédicale)} \\
Douala — Tél : 6XX XX XX XX — \texttt{armand.t@email.pro} — Permis B, véhicule personnel \\ \hline
\textbf{OBJECTIF :} Technicien biomédical itinérant fort de deux ans d'expérience terrain, je maîtrise la maintenance d'équipements médicaux et la relation client, compétences \textbf{dont} j'ai fait la preuve au CHU de Douala et chez MedEquip. Le poste \textbf{que} vous proposez mobilise l'autonomie et l'anglais technique \textbf{qui} caractérisent mon parcours. Disponible immédiatement pour des missions nationales. \\ \hline
\textbf{EXPÉRIENCES PROFESSIONNELLES :} \\
\textbf{2022-2024 (18 mois) — Technicien de maintenance itinérant, MedEquip Services, Douala} \\
Interventions curatives et préventives chez cliniques privées ; formation des utilisateurs (français/anglais) ; rédaction de rapports hebdomadaires sous Salesforce ; taux de résolution en première visite : 90\%. \\
\textbf{2022 (6 mois) — Stage de fin d'études, Service maintenance, CHU de Douala} \\
Gestion du parc matériel biomédical ; maintenance préventive/curative ; relation fournisseurs. \\ \hline
\textbf{FORMATION :} \\
2022 : Licence Professionnelle Maintenance Biomédicale — IUT de Douala, mention Bien. \\ \hline
\textbf{COMPÉTENCES :} Maintenance biomédicale (normes IEC/ISO) ; CRM Salesforce ; Pack Office ; conduite longue distance. \\ \hline
\textbf{LANGUES :} Français (courant) ; Anglais technique (TOEIC 820). \\ \hline
\textbf{QUALITÉS :} Autonomie (tournées solo 18 mois) ; résistance au stress (urgences bloc opératoire) ; sens du service. \\ \hline
\textbf{RÉFÉRENCES :} Chef de service Maintenance, CHU Douala ; Directeur Technique, MedEquip Services. \\ \hline
\end{tabularx}
\end{center}

\textbf{Points de vigilance :} ordre « Expériences avant Formation » justifié (expérience significative de 18 mois) ; verbes d'action et résultats chiffrés ; 2 pronoms relatifs minimum dans l'objectif ; une seule page ; aucune information sensible (CNI, état civil) ; mots-clés de l'offre repris (biomédical, Salesforce, anglais technique, permis B).
\end{corrige}

\begin{corrige}[Exercice 10 — Type Probatoire : Étude comparative de deux CV]
\textbf{Barème indicatif :} grille 6 pts (1 pt par critère correctement coté) ; texte argumenté 14 pts (thèse 2 ; 3 arguments factuels 6 ; nuance 2 ; connecteurs et langue 2 ; conclusion 2).

\textbf{1. Grille d'analyse comparative :}
\begin{center}
\begin{tabularx}{\linewidth}{|l|X|X|}\hline
\rowcolor{popblueL} \textbf{Critère de l'offre} & \textbf{Aliou S.} & \textbf{Chantal N.} \\ \hline
Diplôme (BTS Compta/Gestion) & +++ (BTS Comptabilité-Gestion) & + (BTS Secrétariat, option compta) \\ \hline
Expérience pertinente (3 ans) & +++ (2 ans Comptable unique + 2 ans assistant) & + (pas de comptabilité OHADA formelle) \\ \hline
Maîtrise OHADA / Sage Saari & +++ (SYSCOHADA, Sage Saari v10) & 0 (QuickBooks, non OHADA) \\ \hline
Excel avancé & +++ (TCD, RechercheV) & + (Excel « moyen ») \\ \hline
Qualités transverses / Références & ++ (rigueur, références professionnelles directes) & + (références non professionnelles : pasteur) \\ \hline
Présentation CV / Objectif & ++ (objectif ciblé, à préciser davantage) & + (objectif égocentré : « près de chez moi ») \\ \hline
\end{tabularx}
\end{center}

\textbf{2. Texte argumenté — corrigé type (environ 190 mots) :}

« À l'issue de l'analyse comparative des deux dossiers, je recommande sans hésitation la candidature de M. Aliou S. pour le poste de Comptable Unique.

En effet, son profil satisfait pleinement aux exigences de l'offre. D'abord, il est titulaire du diplôme exact requis, un BTS Comptabilité-Gestion, \textbf{tandis que} Mme N. ne possède qu'un BTS de Secrétariat à option comptable. \textbf{Ensuite}, il justifie de quatre années d'expérience comptable, \textbf{notamment} deux ans comme Comptable unique, poste identique à celui proposé : il maîtrise donc la tenue de la comptabilité générale, les déclarations fiscales et la paie. \textbf{De surcroît}, il pratique le référentiel OHADA et le logiciel Sage Saari, outils explicitement demandés, \textbf{alors que} Mme N. ne connaît que QuickBooks, inadapté au référentiel en vigueur. Son niveau d'Excel avancé (tableaux croisés dynamiques) constitue un atout supplémentaire.

\textbf{Certes}, Mme N. témoigne d'une réelle polyvalence acquise dans la gestion de sa propre boutique ; \textbf{cependant}, l'absence de comptabilité formelle OHADA et ses références non professionnelles fragilisent sa candidature. Un point de vigilance demeure pour M. S. : son objectif, trop bref, gagnerait à être davantage personnalisé.

\textbf{Par conséquent}, M. Aliou S. est le candidat le plus opérationnel : son recrutement garantit une prise de fonction immédiate, sans formation préalable coûteuse. »

\textbf{Points de vigilance :} s'appuyer \emph{uniquement} sur les documents ; citer des faits (années, logiciels, diplômes) ; au moins 3 connecteurs variés ; une nuance sur le candidat retenu est exigée (ici : objectif perfectible).
\end{corrige}

\begin{corrige}[Exercice 11 — Type Bac Technique : Rédaction intégrée et correction linguistique]
\textbf{Barème indicatif :} Partie A : 6 pts (organisation en sous-catégories 2 ; phrases nominales et mots-clés 2 ; 2 locutions adverbiales et ponctuation 2). Partie B : 4 pts (6 erreurs corrigées, 0,5 pt chacune + 1 pt de présentation). Partie C : 5 pts (pronom relatif 1,5 ; adverbe de liaison 1,5 ; correction et élégance 2).

\textbf{Partie A — Rubrique « Compétences » corrigée :}

\textbf{COMPÉTENCES TECHNIQUES \& FONCTIONNELLES}
\begin{itemize}
\item \textbf{Gestion de projet :} planification (MS Project, Primavera P6) ; suivi des coûts, appels d'offres et DPGF ; management d'équipes de 10 à 15 personnes et coordination de sous-traitants ; \textbf{notamment} animation de réunions hebdomadaires de chantier.
\item \textbf{Technique / Chantier :} béton armé et VRD ; lecture de plans d'exécution (EXE) et DAO (AutoCAD, Revit).
\item \textbf{Outils numériques :} tableaux de bord de suivi ; reporting client mensuel ; \textbf{en particulier} reporting rédigé en anglais technique (niveau B2).
\item \textbf{Réglementation / Qualité :} normes NF DTU et Eurocodes ; démarche HQE et sécurité HSE ; \textbf{par ailleurs} maîtrise du contractuel : marchés publics et privés, CCAG, avenants.
\end{itemize}

\textbf{Partie B — Paragraphe corrigé (corrections en gras) :}

« De 2020 à 2022, j'ai occupé le poste de Superviseur Maintenance \textbf{à} l'usine SOSUCAM. J'étais chargé de planifier les arrêts de production, \textbf{de} gérer une équipe de 8 techniciens et de veiller au respect des normes \textbf{de} sécurité. \textbf{Grâce} à mon organisation, le taux de disponibilité des machines est passé de 85\% à 95\%. De plus, j'ai mis en place une GMAO qui a permis de réduire les coûts de maintenance de 15\%. \textbf{J'ai dû quitter ce poste pour des raisons familiales.} »

\textbf{Les 6 erreurs :}
\begin{enumerate}
\item « Grace » $\rightarrow$ \textbf{Grâce} (accent circonflexe obligatoire).
\item « normes sécurité » $\rightarrow$ normes \textbf{de} sécurité (complément du nom).
\item « gérer » non coordonné $\rightarrow$ « \textbf{de} gérer » (parallélisme des infinitifs après « chargé de »).
\item « j'ai du partir » $\rightarrow$ « j'ai \textbf{dû} » (accent circonflexe au participe passé du verbe devoir).
\item « Cependant, j'ai dû partir » $\rightarrow$ connecteur fautif : il n'y a pas d'opposition mais une simple information ; on supprime « cependant » (erreur de liaison logique).
\item « pour raison familiale » $\rightarrow$ « pour \textbf{des} raisons familiales » (article manquant) ; registre : on préfère « quitter ce poste » à « partir ».
\end{enumerate}

\textbf{Partie C — Phrase complexe corrigée :}

« Le candidat a piloté la construction d'un hôpital de 150 lits, chantier \textbf{qui} a été livré avec deux semaines d'avance ; \textbf{de surcroît}, le budget initial a été respecté à 99,5\%. »

Autre formulation acceptée : « Ayant piloté la construction d'un hôpital de 150 lits \textbf{dont} le budget initial a été respecté à 99,5\%, le candidat l'a \textbf{en outre} livré avec deux semaines d'avance. »

\textbf{Points de vigilance :} vérifier l'accord du participe passé (« livré » s'accorde avec « chantier ») ; le point-virgule précède l'adverbe de liaison ; une seule phrase complexe est exigée (pas trois phrases juxtaposées).
\end{corrige}

\newpage

\coursec{Fiche bilan}

\begin{popfigure}
\begin{tikzpicture}[mindmap, grow cyclic, every node/.style={concept}, root concept/.append style={font=\bfseries, fill=popblue, text=white, minimum size=3.2cm}, level 1 concept/.append style={level distance=3.6cm, sibling angle=60, font=\small}, concept color=popblue]
\node[root concept]{Le CV}
  child[concept color=poporange] { node {Définition et fonction} }
  child[concept color=popgreen] { node {Structure externe : mise en page, titre, photo} }
  child[concept color=poppurple] { node {Rubriques : état civil, formation, expérience} }
  child[concept color=poppink] { node {Outils de liaison : adverbes, relatifs} }
  child[concept color=popgold] { node {Rédaction : travail préparatoire} }
  child[concept color=popblue!60] { node {Adapter à l'offre d'emploi} };
\end{tikzpicture}
\legende{Carte mentale de la leçon : les six points essentiels du Curriculum Vitae.}
\end{popfigure}

\begin{bilan}
\textbf{Définitions clés}
\begin{itemize}
\item Le \cle{Curriculum Vitae} (CV) : texte \cle{fonctionnel} qui présente de manière claire et synthétique le parcours, les compétences et les qualités d'un candidat à un emploi, à un stage ou à une formation.
\item \cle{Structure externe} : la présentation matérielle (papier A4, une page, titre, mise en page aérée, rubriques visibles).
\item \cle{Structure interne} : le contenu organisé en \cle{rubriques}.
\item Outils de \cle{liaison} : mots et groupes de mots qui relient les idées dans la phrase et entre les phrases.
\end{itemize}

\textbf{Les rubriques du CV (ordre conseillé)}
\begin{center}
\begin{tabularx}{\linewidth}{|l|X|}\hline
\rowcolor{popblueL} \textbf{Rubrique} & \textbf{Contenu} \\ \hline
État civil & Nom, prénom, date et lieu de naissance, nationalité, situation de famille \\ \hline
Coordonnées & Adresse, téléphone, adresse électronique \\ \hline
Formation / diplômes & Du plus récent au plus ancien (ordre antichronologique) \\ \hline
Expérience professionnelle & Stages, emplois, avec dates et fonctions \\ \hline
Compétences & Langues, informatique, aptitudes techniques \\ \hline
Centres d'intérêt & Sports, loisirs, activités associatives \\ \hline
Références & Personnes à contacter (facultatif) \\ \hline
\end{tabularx}
\end{center}

\textbf{Outils de liaison à connaître par cœur}
\begin{center}
\begin{tabularx}{\linewidth}{|l|X|}\hline
\rowcolor{popblueL} \textbf{Outil} & \textbf{Exemples} \\ \hline
Adverbes de liaison & d'abord, ensuite, enfin, donc, cependant, ainsi \\ \hline
Locutions adverbiales & en effet, c'est-à-dire, par conséquent, en outre \\ \hline
Pronoms relatifs & qui, que, où, dont, lequel, ce qui, ce que \\ \hline
\end{tabularx}
\end{center}

\textbf{Méthode express de rédaction}
\begin{enumerate}
\item Lire attentivement l'offre d'emploi et relever les qualités demandées.
\item Faire l'inventaire de ses informations (travail préparatoire) : diplômes, stages, compétences.
\item Classer les informations par rubriques, du plus récent au plus ancien.
\item Rédiger en phrases courtes ou en rubriques nominales, avec des verbes d'action.
\item Relier les idées avec les outils de liaison ; relire et corriger l'orthographe.
\end{enumerate}
\end{bilan}

\begin{aretenir}[Checklist avant le Bac]
\begin{itemize}
\item[$\square$] Je sais définir le CV et dire à quoi il sert.
\item[$\square$] Je connais les qualités de la structure externe : une page, clarté, aération, titre.
\item[$\square$] Je cite sans hésiter les rubriques principales du CV.
\item[$\square$] Je présente la formation et l'expérience du plus récent au plus ancien.
\item[$\square$] Je distingue adverbes de liaison, locutions adverbiales et pronoms relatifs.
\item[$\square$] J'emploie correctement \emph{qui}, \emph{que}, \emph{où}, \emph{dont} dans mes phrases.
\item[$\square$] Je sais adapter mon CV à une offre d'emploi précise.
\item[$\square$] Je rédige des formulations brèves, sans faute d'orthographe ni abréviations.
\item[$\square$] Je vérifie que mes coordonnées sont complètes et exactes.
\item[$\square$] Je relis mon CV avant de le rendre et je soigne la présentation.
\end{itemize}
\end{aretenir}

\findecours

\end{document}
